\documentclass[11pt,reqno]{amsart}

\usepackage[margin=1in]{geometry}
\usepackage{amsmath,amssymb,amsthm}
\usepackage{array}
\setlength{\emergencystretch}{1em}

\newcommand{\R}{\mathbb{R}}
\newcommand{\C}{\mathbb{C}}
\newcommand{\Z}{\mathbb{Z}}
\newcommand{\N}{\mathbb{N}}
\newcommand{\E}{\mathbb{E}}
\newcommand{\Sph}{\mathbb{S}^2}
\newcommand{\TV}{d_{\mathrm{TV}}}
\newcommand{\Var}{\operatorname{Var}}
\newcommand{\gap}{\operatorname{gap}}
\newcommand{\Ran}{\operatorname{Ran}}
\newcommand{\dd}{\mathrm{d}}
\newcommand{\cs}{c_{\mathrm{s}}}
\newcommand{\ct}{c_{\mathrm{t}}}
\newcommand{\Pext}{\mathcal{P}_{\mathrm{ext}}}
\newcommand{\Dcons}{\mathcal{D}_{\mathrm{cons}}}
\newcommand{\certpending}[1]{\textbf{[certificate check pending]}}

\theoremstyle{plain}
\newtheorem{theorem}{Theorem}[section]
\newtheorem{proposition}[theorem]{Proposition}
\newtheorem{lemma}[theorem]{Lemma}
\newtheorem{corollary}[theorem]{Corollary}
\newtheorem{hypothesis}[theorem]{Hypothesis}
\theoremstyle{definition}
\newtheorem{definition}[theorem]{Definition}
\theoremstyle{remark}
\newtheorem{remark}[theorem]{Remark}

\numberwithin{equation}{section}

\begin{document}

\title{The spectral gap of the spin-2 AKLT model on the square lattice}

\begin{abstract}
We prove that the spin-2 AKLT Hamiltonian on the square-lattice torus
has a unique ground state for every pair of integers $L,N\ge3$ and a
spectral gap bounded below by one positive constant for all integers
$L,N\ge L_0$, for some fixed $L_0$, with either parity allowed for
each period. The proof is computer-assisted. Its computational inputs
are local operator inequalities, an open-trail growth bound, and a
type-dependent positive-vector bound for a response kernel. These
inputs are checked using exact matrix positivity certificates, a
finite-memory test, and stationary Bellman inequalities with finite
prefix sums. The accompanying formalization uses the Lean kernel
to verify the finite certificate checks; pending checks are marked
explicitly.
Three finite coherent-state frames compare the quantum gap with the
gap of a finite-state heat-bath dynamics. A sign change on each lifted
square permits arbitrary torus parities. A direct periodic block
argument reduces the heat-bath gap to window estimates, which follow
from a separating-cut factorization, a walk expansion, and a
contracting response recursion. Uniqueness follows independently from
polynomial divisibility.
\end{abstract}

\maketitle
\setcounter{tocdepth}{1}
\tableofcontents

\section{The result and conventions}\label{s:result}

We study the spin-\(2\) square-lattice AKLT model of Affleck, Kennedy,
Lieb, and Tasaki~\cite{AKLT88,KLT88}. We prove uniqueness of its ground
state on every torus with integer periods \(L,N\ge3\) and a uniform
positive spectral gap for all integers \(L,N\ge L_0\), with no parity
restriction, for one fixed threshold \(L_0\). The computational inputs
are local operator inequalities, an open-trail growth bound, and a
type-dependent positive-vector bound obtained from stationary Bellman
inequalities and finite prefix sums. The accompanying formalization
uses the Lean kernel to verify the finite certificate checks; pending
checks are marked explicitly. The polynomial realization of the spin space,
including its Hilbert space normalization, is fixed in
Definition~\ref{d:poly}.

\begin{definition}\label{d:model}
For integers \(L,N\ge3\), let \(\Lambda=\Lambda_{L,N}=\mathbb Z_L\times\mathbb Z_N\)
with its simple, unoriented nearest-neighbour edge set \(E_{L,N}\).
The Hilbert space is \(\mathcal H_\Lambda=\bigotimes_{x\in\Lambda}V_4\).
Each site carries the spin-\(2\) space \(V_4\), and \(P_{\{x,y\}}\)
is the orthogonal projection onto total spin \(4\) in \(V_4\otimes V_4\).
Fix one orientation of each edge and set
\[
 q_{xy}=z_{x,1}z_{y,2}-z_{x,2}z_{y,1},\qquad
 B_{L,N}=\prod_{e\in E_{L,N}}q_e,\qquad
 H_{L,N}=\sum_{e\in E_{L,N}}P_e.
\]
Every edge is counted once. The spectral gap of \(H_{L,N}\) means its
smallest positive eigenvalue.
\end{definition}

Here an edge joins two distinct sites whose difference is one of
\((1,0),(-1,0),(0,1),(0,-1)\), with coordinates reduced modulo the
periods. The edge set is a set of unordered pairs, so that an edge
is included only once even if it has more than one description.
The notation \(q_e\) uses the chosen orientation; \(P_e\) does not.
The polynomial \(B_{L,N}\) is the valence-bond state.

\begin{theorem}\label{t:main}
For every \(L,N\ge3\), \(B_{L,N}\ne0\) and
\(\ker H_{L,N}=\mathbb C B_{L,N}\).
There exist \(c>0\) and an integer \(L_0\ge3\) such that, for all
integers \(L,N\ge L_0\), every nonzero eigenvalue of \(H_{L,N}\)
is at least \(c\).
\end{theorem}

\begin{proof}
Proposition~\ref{p:toruskernel} proves the nonvanishing and kernel
assertions for every \(L,N\ge3\). For the uniform gap,
Theorem~\ref{t:inputs} supplies (H1), (H2) with
\[
 C_V=10\ge1,\qquad
 0<q=\frac{1119948287353}{1120437614225}<1,
\]
and (H3\(_h\)) for the prescribed type-dependent kernel and fixed
orderings. The common thresholds are \(c_{\rm s}=961/1000\) and
\(c_{\rm t}=187/200\).
These are exactly the hypotheses of Theorem~\ref{t:rec}, which
gives finite \(A\ge2\) and \(0<\theta<1\) such that
\(\mathsf P(A,\theta)\) holds.
Theorem~\ref{t:cutMOS} then gives
\(\mathsf W(64,p_0,M_\star)\) for the plane frame family, with
one \(p_0\in(0,1]\) and one integer \(M_\star\ge1\) valid for
every proper plane three-colouring and every translate.
This is the hypothesis of the periodic assembly theorem,
Theorem~\ref{t:mainW}, proved in Section~4, for the plane frame
family of Definition~\ref{d:framemodel}. Its direct quantum
conclusion supplies \(c>0\) and an integer \(L_0\ge3\) with the
asserted spectral bound for all integers \(L,N\ge L_0\).
\end{proof}

\subsection{Structure of the proof}

The uniqueness assertion has a separate algebraic proof. The kernel
of each edge projection consists of multiples of its edge determinant.
Unique factorization combines these divisibilities, and degree four
at every torus site leaves a scalar quotient. This argument applies
also when either period is three.

For the gap, finite fourth moments in Section~2 give tight frames of
coherent states. In Section~3 we assign three explicit rotated
icosahedral frames according to a proper torus colouring. Their
separation makes the squared modulus of the frame image of
\(B_{L,N}\) strictly positive. The resulting antiferromagnetic
probability law on the torus has a single-site heat bath whose gap
bounds the quantum gap from below. A conditional physical square
lifts to the plane, where a coordinatewise sign change converts its
law to the positive-interaction plane specification. This sign
change is defined on the lift.

Section~4 uses a uniform estimate for window marginals to choose one
square scale. A stopped coupling bounds the Hamming cost of changing
one exterior letter. Counting all periodic translates of that square
gives a gap for block updates, and comparison of conditional block
and site Dirichlet forms gives the torus site gap.
Theorem~\ref{t:mainW} combines this with the quantum comparison of
Theorem~\ref{t:QD} to obtain the uniform quantum gap.
Section~5 proves the required window estimate from the pair-response
bound \(\mathsf P(A,\theta)\), uniformly over plane colourings and
placements. Sections~6--8 prove that response bound using a walk
expansion, a contracting positive-vector recursion, and the exact
computational inputs. Thus the gap implication is
\[
 \begin{gathered}
 \text{exact computational inputs}
 \ \Longrightarrow\ \mathsf P
 \ \Longrightarrow\ \mathsf W\\
 \ \Longrightarrow\ \text{periodic block gap}
 \ \Longrightarrow\ \text{site gap}
 \ \Longrightarrow\ \text{quantum gap}.
 \end{gathered}
\]
The finite probability and comparison tools are proved in Appendix~A;
the planar topology and the computational certificate arguments are
proved in Appendices~B--D. All scale choices precede the choice of the
torus periods, colouring, block placement and exterior configuration.

\subsection{Conventions}

All Hilbert spaces are complex, with inner products conjugate linear
in the first argument. An inequality between self-adjoint operators
means the corresponding inequality of quadratic forms.
We abbreviate \(B=B_{L,N}\) and \(H=H_{L,N}\) when the periods are
fixed. Products over empty sets are one and sums over empty sets
are zero.

On the plane,
\[
 d(x,y)=|x_1-y_1|+|x_2-y_2|,\qquad
 |x-y|_\infty=\max\{|x_1-y_1|,|x_2-y_2|\}.
\]
We write \(x\sim y\) for nearest neighbours and
\(d(x,S)=\min_{y\in S}d(x,y)\), with the value \(+\infty\) if
\(S=\varnothing\). Distances in plane estimates are these ambient
distances, even if sites have been removed. A torus distance, when
needed, is stated explicitly.
Lattice translations lie in \(\mathbb Z^2\), and a radius used as
an endpoint of a lattice interval is a nonnegative integer.

For probabilities \(\nu,\nu'\) on a finite set,
\[
 d_{\rm TV}(\nu,\nu')
 =\frac12\sum_s|\nu(s)-\nu'(s)|
 =\max_A|\nu(A)-\nu'(A)|.
\]
For two configurations on the same product of site alphabets, their
Hamming distance is the number of sites at which their letters
differ. Conditional expectations and probabilities are finite sums
on positive conditioning events. Variances use squared complex
moduli. A claimed Poincare gap \(g>0\) means
\(\mathcal D(f)\ge g\operatorname{Var}(f)\) for every \(f\).
This convention includes a one-point configuration space, for which
both sides vanish; the infimum over positive-variance functions is
then \(+\infty\). Every heat-bath update rate is one per physical
site unless a different set of update labels is explicitly specified.

\section{Finite moment algebra and the response input}\label{s:moments}

We use finite product sums with explicitly specified moments through degree
four. Pins are arbitrary unit vectors on the original legs; deleting a
site creates no new pin. All degree bounds below count the inserted
observable together with the interaction factors.

\begin{definition}\label{d:moment}
Write \(\mathbb S^2=\{a\in\mathbb R^3:|a|=1\}\), put
\(\phi=(1+\sqrt5)/2\), \(r=\sqrt{\phi+2}\), and set
\[
 D=\frac1r\{(0,\pm1,\pm\phi),(\pm1,\pm\phi,0),(\pm\phi,0,\pm1)\},
 \qquad
 \mathcal M(f)=\frac1{12}\sum_{a\in D}f(a).
\]
All indicated signs are independent. Define finite probability laws
\(\lambda_{\mathrm U}=\frac1{12}\sum_{a\in D}\delta_a\) and
\(\lambda_{\mathrm D}=\frac35\lambda_{\mathrm U}+\frac25\delta_0\).
Products of these functionals always mean finite sums.
For plane sites, \(d\) denotes ambient \(\ell^1\) distance and
\(\operatorname{dist}_\infty\) denotes ambient \(\ell^\infty\) distance.
Distances to an empty set are \(+\infty\); a product over an empty site
set has its one empty configuration. For finite probabilities, use
\[
 d_{\rm TV}(p,p')=\tfrac12\sum_z|p(z)-p'(z)|,\qquad
 \operatorname{Var}_\mu f=\mu(|f-\mu f|^2),\qquad
 d_H(n,n')=\sum_x\mathbf1_{\{n_x\ne n'_x\}} .
\]
\end{definition}

For a vector-valued polynomial, the functional is applied to each
coordinate. In particular the expressions for means below are finite
vector sums.

\begin{lemma}\label{l:moments}
The set \(D\) has twelve distinct unit vectors and is antipodal.
For \(i,j,k,l\in\{1,2,3\}\),
\[
 \begin{gathered}
 \mathcal M(1)=1,\qquad \mathcal M(n_i)=\mathcal M(n_in_jn_k)=0,\qquad
 \mathcal M(n_in_j)=\delta_{ij}/3,\\
 \mathcal M(n_in_jn_kn_l)=
 (\delta_{ij}\delta_{kl}+\delta_{ik}\delta_{jl}+\delta_{il}\delta_{jk})/15.
 \end{gathered}
\]
For every polynomial \(p\) of total degree at most four and every
\(R\in\mathrm O(3)\), \(\mathcal M(p\circ R)=\mathcal M(p)\).
The same covariance holds for \(\lambda_{\mathrm D}\); its constant
moment is one and its moments of positive degree at most four are
\(3/5\) times those displayed above.
These assertions hold site by site for polynomials of separate degree
at most four in a finite product of spins.
\end{lemma}

\begin{proof}
The identities \(\phi^2=\phi+1\), \(r^2=\phi+2=1+\phi^2\) show that
every displayed vector has norm one. Each vector has precisely one zero
coordinate; its location specifies one of the three displayed groups,
and its two nonzero signs specify one of four distinct vectors within
that group. Thus there are twelve distinct vectors. Reversing both signs
gives the antipodal vector.

The set is preserved by changing any individual coordinate sign and by
cyclically permuting the coordinates. A monomial with any odd coordinate
exponent consequently has zero average. For each coordinate \(i\), and
distinct coordinates \(i,j\), direct summation over the three groups gives
\[
 \mathcal M(n_i^2)=\frac{1+\phi^2}{3r^2}=\frac13,\qquad
 \mathcal M(n_i^4)=\frac{1+\phi^4}{3r^4}=\frac15,\qquad
 \mathcal M(n_i^2n_j^2)=\frac{\phi^2}{3r^4}=\frac1{15}.
\]
Here \(r^4=5\phi^2\) and \(1+\phi^4=3\phi^2\).
These cases, the constant, and the vanishing odd-exponent cases exhaust
the monomials through degree four. For four indices, the all-equal case
has three nonzero Kronecker products, the two-equal-pairs case has one,
and every other case has zero; this gives the stated tensor.

For \(R\in\mathrm O(3)\), substituting
\((Rn)_i=\sum_aR_{ia}n_a\) into the second tensor gives
\(\sum_aR_{ia}R_{ja}/3=\delta_{ij}/3\).
Substitution into the fourth tensor gives the three products of the
same contractions, divided by \(15\). The first and third tensors
remain zero, and the constant remains one. Linearity proves covariance
for every polynomial in the stated degree range.

For \(\lambda_{\mathrm D}\), a positive-degree monomial vanishes at
zero, so its moment is \(3/5\) times the corresponding moment of
\(\mathcal M\); the constant moment is \(3/5+2/5=1\).
Zero is fixed by \(R\), so covariance follows as well.
Finally a polynomial of separate degree at most four is a finite sum
of products of single-site monomials of degrees at most four.
Its product functional factors on every such product. Applying the
single-site conclusions to each factor proves the last assertion,
including different orthogonal substitutions at different sites.
\end{proof}

\begin{definition}\label{d:pinmodel}
For finite \(G\subset\mathbb Z^2\), let
\[
 \mathcal L(G)=\{(x,z):x\in G,\ z\notin G,\ d(x,z)=1\},\qquad
 S(G)=\{x\in G:d(x,\mathbb Z^2\setminus G)\le2\}.
\]
A pin assignment is \(\tau\in(\mathbb S^2)^{\mathcal L(G)}\).
For \(Y\subset G\), write \(D_Y=G\setminus Y\) and define
\[
 Z^\tau_{D_Y}[f]=
 \sum_{n\in D^{D_Y}}12^{-|D_Y|}f(n)
 \prod_{\substack{\{x,y\}\subset D_Y\\d(x,y)=1}}(1+n_x\cdot n_y)
 \prod_{\substack{(x,z)\in\mathcal L(G)\\x\in D_Y}}
       (1+\tau_{xz}\cdot n_x).
\]
The free sum \(Z^0_{D_Y}[f]\) omits all leg factors.
Write \(Z^\bullet_{D_Y}=Z^\bullet_{D_Y}[1]\),
\(\mathbb E^\bullet_{D_Y}f=Z^\bullet_{D_Y}[f]/Z^\bullet_{D_Y}\),
and \(\Psi_{D_Y}(\tau)=Z^\tau_{D_Y}/Z^0_{D_Y}\).
Here and below \(\bullet\) is either \(\tau\) or \(0\).
The empty partition sum is one. There are no new factors on legs into \(Y\).
Different original legs may carry unrelated unit pins, even if their
outside endpoints coincide.
\end{definition}

The argument \(f\) is a function of the surviving spins. In the
partition sum every surviving site has at most four factors: each
factor uses one of its four original lattice directions.

\begin{lemma}\label{l:finitefacts}
Consider a finite network with site laws \(\lambda_{\mathrm U}\) or
\(\lambda_{\mathrm D}\), retained distinct-endpoint bonds
\(1+n_x\cdot n_y\), and retained original unit-pin factors
\(1+\tau_\ell\cdot n_x\), with at most four factors incident to each site.
Its pinned and free partition sums are strictly positive, including
the empty network. Non-neighbours are conditionally independent given
the other spins on every conditioning event of positive probability.
If \(f\) is a product of scalar products of spins and the product of
\(f\) with the bond factors has separate degree at most four at every
site, then \(Z^0[f]\ge0\), hence \(\mathbb E^0 f\ge0\).
All normalized spins have norm at most one.
\end{lemma}

\begin{proof}
Every site law has finite support in the closed unit ball, and every
bond or pin factor lies in \([0,2]\). Thus every sum is finite and
nonnegative. To prove strict positivity, order the nonremoved sites
as \(x_1,\ldots,x_m\), and choose at each site a letter in the
twelve-point set \(D\). Such a letter has positive mass under either
law, namely \(1/12\) or \(1/20\).

At stage \(i\), each pin at \(x_i\) excludes at most its antipodal
letter, and each bond to an already chosen site excludes at most the
antipode of that site's letter. There are at most four incident
factors in total, so at most four of the twelve letters are excluded.
Choose any remaining letter. A bond to a later site will be made
strictly positive when that site's letter is chosen. At the end every
retained factor is strictly positive, and the product site mass is
positive. This single summand proves strict positivity. The same
argument with pins omitted proves positivity of the free sum; the
empty network has sum one.

For conditional independence, let \(A\) be any set of sites with no
retained bond between its distinct members, and fix all spins off
\(A\) at a configuration of positive marginal probability. The
unnormalized conditional weight has the form
\[
 C\prod_{x\in A}
 \left(\lambda_{\mathrm{st}(x)}(n_x)
       \prod_{\substack{\text{retained factors}\\\text{incident to }x}}
                     (1+n_x\cdot b)\right).
\]
Here each \(b\) is a fixed neighbouring spin or pin, and factors
wholly off \(A\) are included in \(C\). Positive marginal probability
implies \(C>0\) and that every local sum in this product is positive.
The conditional normalizer is \(C\) times the product of those sums.
Dividing therefore gives a product of the single-site conditional
laws. Taking \(A\) to consist of two non-neighbours gives the stated
assertion.

For free positivity expand each retained factor \(1+n_x\cdot n_y\)
into its two terms, and expand every scalar product, including those
in \(f\), as \(\sum_{i=1}^3n_{x,i}n_{y,i}\). Every coefficient in
these finite expansions is nonnegative. For any resulting choice
of coordinate indices, the sum over spins is a product of
single-site moment tensors. The assumed separate degree bound
allows only degrees \(0,1,2,3,4\).
By Lemma~\ref{l:moments}, a degree-one or degree-three tensor is zero,
and the tensors of degrees zero, two and four are sums of products
of Kronecker deltas with nonnegative coefficients.
The same is true for \(\lambda_{\mathrm D}\), with the extra
positive factor \(3/5\) at each site of positive even degree.
Consequently every term after these substitutions is nonnegative.
This also covers scalar products with coincident endpoints, counting
both appearances in the degree. Thus \(Z^0[f]\ge0\); division by the
strictly positive \(Z^0\) proves the normalized assertion.
The supports give \(|n_x|\le1\) pointwise, and hence the same bound
for any conditional or unconditional mean.
\end{proof}

\begin{lemma}\label{l:35}
For \(0\le k\le3\) and \(|a_j|\le1\), the probability law proportional to
\(\prod_{j=1}^k(1+n\cdot a_j)\lambda_{\mathrm U}(\mathrm dn)\)
has mean of norm at most \(3/5\). The same bound holds with
\(\lambda_{\mathrm D}\) in place of \(\lambda_{\mathrm U}\).
Moreover,
\[
 \left\{
 \frac{\mathcal M\!\left(n\prod_{j=1}^3(1+n\cdot a_j)\right)}
      {\mathcal M\!\left(\prod_{j=1}^3(1+n\cdot a_j)\right)}
       :a_1,a_2,a_3\in\mathbb S^2
 \right\}=\{b\in\mathbb R^3:|b|\le3/5\}.
\]
Consequently an unpinned site having at most three present factors
has conditional mean of norm at most \(3/5\), on every positive
conditioning fibre.
\end{lemma}

\begin{proof}
First take three unit vectors \(a,b,c\), and put
\[
 p=a\cdot b+a\cdot c+b\cdot c,\qquad
 A=\mathcal M\bigl((1+n\cdot a)(1+n\cdot b)(1+n\cdot c)\bigr).
\]
The second moment gives \(A=1+p/3\). Since
\(0\le|a+b+c|^2=3+2p\), we have \(A\ge1/2\).
The second and fourth moments give
\[
 15\,\mathcal M\bigl(n(1+n\cdot a)(1+n\cdot b)(1+n\cdot c)\bigr)
 =V,\qquad
 V=5(a+b+c)+a(b\cdot c)+b(a\cdot c)+c(a\cdot b).
\]
Thus the mean is \(V/(5(3+p))\). Set
\(X=1-b\cdot c\), \(Y=1-a\cdot c\), \(Z=1-a\cdot b\);
each lies in \([0,2]\). Expanding the scalar products in \(|V|^2\)
and using \(|a|^2=|b|^2=|c|^2=1\) gives the polynomial identity
\[
 (9+3p)^2-|V|^2
 =4\bigl((X-Y)^2+(X-Z)^2+(Y-Z)^2\bigr)+6XYZ.
\]
The right-hand side is nonnegative, and \(9+3p\ge9/2>0\).
Taking square roots gives \(|V|\le3(3+p)\), which proves the
\(3/5\) bound for three unit fields.

We spell out the reduction to at most three fields in the closed
unit ball. For any such field \(a\), choose a unit vector \(e\)
parallel to \(a\) if \(a\ne0\), and any unit vector if \(a=0\).
With \(p_a=(1+|a|)/2\),
\[
 1+n\cdot a=p_a(1+n\cdot e)+(1-p_a)(1-n\cdot e).
\]
An absent factor can likewise be inserted as
\(1=\tfrac12(1+n\cdot e)+\tfrac12(1-n\cdot e)\).
Expand these identities until there are exactly three unit fields
in every term. The result is a finite convex combination
\(W=\sum_\eta t_\eta W_\eta\), with \(t_\eta\ge0\),
\(\sum_\eta t_\eta=1\), and
\(\mathcal M(W_\eta)=A_\eta\ge1/2\). The normalized mean for \(W\)
is the convex combination of the means for \(W_\eta\) with weights
\(t_\eta A_\eta/\sum_\xi t_\xi A_\xi\).
These weights are nonnegative and sum to one, proving the bound
and positivity of its normalizer in every case \(0\le k\le3\).

For the same product \(W\), let \(A=\mathcal M(W)>0\) and let
\(\mu=\mathcal M(nW)/A\). Since \(W(0)=1\), the mean under the
normalized \(\lambda_{\mathrm D}\) law is
\[
 \frac{(3/5)A}{(3/5)A+2/5}\,\mu.
\]
The prefactor belongs to \((0,1)\), so its norm is at most \(3/5\).

It remains to show that every vector in the ball occurs with three
unit fields under \(\mathcal M\). Take
\[
 a=e_1,\quad b_t=\cos t\,e_1+\sin t\,e_2,\quad
 c_t=\cos t\,e_1-\sin t\,e_2,\qquad 0\le t\le2\pi/3.
\]
The displayed rational formula for the mean has denominator
\(5(3+p)\ge15/2\), so it is continuous along this path.
At \(t=0\), the formula gives mean \(3e_1/5\).
At \(t=2\pi/3\), the three pairwise scalar products are \(-1/2\)
and \(a+b_t+c_t=0\), so \(V=0\).
The continuous norm therefore takes every value in \([0,3/5]\).
The moment formulas also show directly that simultaneously replacing
the three fields by \(Ra,Rb,Rc\) replaces the mean by \(R\mu\),
for every \(R\in\mathrm O(3)\).
For any nonzero desired vector, choose a point on the path with
its norm and an orthogonal matrix sending that mean to the desired
vector; the zero vector already occurs at the endpoint.
The preceding norm bound excludes vectors outside the ball.
This covariance argument concerns only the displayed polynomials,
of degrees three and four.

On a positive conditioning fibre at an unpinned site, its incident
factors give precisely one of the laws just considered, with
neighbouring spins of norm at most one. Positivity is as in
Lemma~\ref{l:finitefacts}. Applying the corresponding single-site
bound proves the final assertion.
\end{proof}

\begin{lemma}\label{l:exact}
At each site of a finite set choose any positive weighted unit-vector
law whose moments through degree four equal those of \(\mathcal M\).
A polynomial of separate degree at most four has the same product
functional under these laws as under \(\lambda_{\mathrm U}\) at all
sites. The analogous assertion holds after replacing any chosen site
laws by \(3/5\) times that law plus \(2/5\,\delta_0\).
In particular this applies to partition products of factors
\(1\pm n_x\cdot n_y\) and \(1\pm n_x\cdot a\), \(|a|=1\), whenever
each variable occurs at most four times, and to inserted observables
only when their degrees are included in this same bound.
\end{lemma}

\begin{proof}
Write a polynomial on sites \(x_1,\ldots,x_m\) as a finite linear
combination of monomials
\[
 \prod_{j=1}^m\prod_{i=1}^3 n_{x_j,i}^{\,\gamma_{j,i}},
 \qquad \sum_{i=1}^3\gamma_{j,i}\le4\quad(1\le j\le m).
\]
The product functional of such a monomial is the product of its
single-site moments. Each factor agrees with its
\(\lambda_{\mathrm U}\) value by hypothesis. This proves the first
assertion by linearity, including the empty product.
At a replaced site, the constant moment is unchanged and every
positive-degree moment is multiplied by \(3/5\), since that
monomial vanishes at zero. These are exactly the moments of
\(\lambda_{\mathrm D}\). The same monomial argument proves the
assertion for any specified subset of replaced sites.

A bond factor is affine in each distinct endpoint, and a pin factor
is affine in its spin. Expanding their product therefore has the
claimed separate degree bound when each spin has at most four
appearances, with multiplicity. If an observable is inserted, its
degree is added at each site before this test is applied.
\end{proof}

This replacement identity compares summed-out polynomial expressions.
In particular it supplies the partition-function comparisons for
different finite alphabets with these moments.

\begin{definition}[(P)]\label{h:P}
For \(A\ge2\) and \(0<\theta<1\), property \(\mathsf P(A,\theta)\)
means that, for every finite \(G\subset\mathbb Z^2\), every unit-pin
assignment on \(\mathcal L(G)\), every
\(Y\subset G\setminus S(G)\), and all \(u,v\in G\setminus Y\) with
\(d(u,v)=2\) and a common neighbour in \(Y\),
\[
 \left|\mathbb E^\tau_{G\setminus Y}(n_u\cdot n_v)
             -\mathbb E^0_{G\setminus Y}(n_u\cdot n_v)\right|
 \le A\theta^{\min\{d(u,S(G)),d(v,S(G))\}} .
\]
The distances are ambient plane distances to the original \(S(G)\);
removed sites are not added to the boundary.
\end{definition}



\section{The quantum model, uniqueness, and positive frames}
\label{s:quantum-frames}

\subsection{Polynomial normalization and the edge kernel}

We use the Pauli matrices
\[
 \sigma_1=\begin{pmatrix}0&1\\1&0\end{pmatrix},\qquad
 \sigma_2=\begin{pmatrix}0&-i\\i&0\end{pmatrix},\qquad
 \sigma_3=\begin{pmatrix}1&0\\0&-1\end{pmatrix},
 \qquad a\cdot\sigma=\sum_{j=1}^3a_j\sigma_j.
\]

\begin{definition}\label{d:poly}
Let \(V_d\) be the homogeneous binary polynomials of degree \(d\).
Its monomials \(z_1^{d-j}z_2^j\), \(0\le j\le d\), are orthogonal with
squared norms \(\binom dj^{-1}\). The action is
\((U_gp)(z)=p(g^{-1}z)\). Use \(V_4\) at each physical site.
For each of the finitely many frame letters \(a\) choose one unit
spinor \(\zeta(a)\in\mathbb C^2\) with
\(\zeta(a)\zeta(a)^*=(I+a\cdot\sigma)/2\).
No spinor choice on the whole sphere is part of the definitions.
\end{definition}

Here \(d\) is a nonnegative integer and \(g\in\mathrm{SU}(2)\).
We record explicitly why this is the spin-\(d/2\) Hilbert space.
In the \(d\)-fold tensor product of the dual of \(\mathbb C^2\),
average the \(\binom dj\) distinct elementary tensors having \(j\)
second-coordinate covectors. Evaluation of this average on
\(z^{\otimes d}\) is \(z_1^{d-j}z_2^j\), and its squared norm is
\(\binom dj\binom dj^{-2}=\binom dj^{-1}\).
These averages are orthogonal for different \(j\).
The dual action of a unitary matrix is unitary, as is its tensor
power; restriction to symmetric tensors therefore gives precisely
the invariant inner product of Definition~\ref{d:poly}.

For completeness, irreducibility follows directly from the monomials
\(m_j=z_1^{d-j}z_2^j\). The diagonal subgroup
\(g_t=\operatorname{diag}(e^{it/2},e^{-it/2})\) acts on \(m_j\)
by \(e^{it(j-d/2)}\). When \(d\ge1\), take
\(t=2\pi/(d+1)\); the \(d+1\) eigenvalues are distinct, so each
monomial projection is a polynomial in \(U_{g_t}\).
A nonzero invariant subspace consequently contains some \(m_j\).
For the skew-Hermitian trace-zero matrices
\[
 A=\begin{pmatrix}0&1\\-1&0\end{pmatrix},\qquad
 C=\begin{pmatrix}0&i\\i&0\end{pmatrix},
\]
differentiating the action along \(\exp(tA)\) and \(\exp(tC)\)
gives, respectively,
\(dU(A)=z_1\partial_{z_2}-z_2\partial_{z_1}\) and
\(dU(C)=-i(z_1\partial_{z_2}+z_2\partial_{z_1})\).
The subspace is stable under these operators and hence under their
complex linear combinations
\(z_1\partial_{z_2}=(dU(A)+i\,dU(C))/2\) and
\(z_2\partial_{z_1}=(-dU(A)+i\,dU(C))/2\).
Their actions are
\[
 z_1\partial_{z_2}m_j=j m_{j-1}\quad(j\ge1),\qquad
 z_2\partial_{z_1}m_j=(d-j)m_{j+1}\quad(j<d).
\]
The indicated coefficients are nonzero. Iterating reaches all
monomials, proving irreducibility. For \(d=0\) it is immediate.
The weights are \(-d/2,-d/2+1,\ldots,d/2\), which identifies the
spin. Differentiation here preserves the invariant subspace because
it is finite dimensional and hence closed.

If \(|a|=1\), the Hermitian matrix \((I+a\cdot\sigma)/2\) is a
rank-one orthogonal projection: its trace is one and
\((a\cdot\sigma)^2=I\). A unit vector in its range supplies the
spinor required in Definition~\ref{d:poly}. We make this choice only
for the finite sets under consideration. Changing such a choice by
a phase does not change any squared modulus below.
The tensor product \(\mathcal H_\Lambda\) is identified with
polynomials homogeneous of degree four in each pair
\(z_x=(z_{x,1},z_{x,2})\).

\begin{lemma}\label{l:edge}
For \(e=\{x,y\}\) and \(\varphi\in\mathcal H_\Lambda\),
\(P_e\varphi=0\) if and only if
\(\varphi=q_e\chi\), where \(\chi\) is homogeneous of degree three
in \(z_x,z_y\) and degree four in every other site variable.
\end{lemma}

\begin{proof}
First consider just the two endpoint sites. Define the diagonal
evaluation map
\[
 \Delta:V_4\otimes V_4\longrightarrow V_8,\qquad
 (\Delta p)(z)=p(z,z).
\]
It is equivariant for the diagonal \(\mathrm{SU}(2)\) action.
For \(0\le k\le8\), choose \(j=\max\{0,k-4\}\); then both
\(j\) and \(k-j\) lie in \(\{0,\ldots,4\}\), and
\[
 \Delta\bigl(z_{x,1}^{4-j}z_{x,2}^j
              z_{y,1}^{4-(k-j)}z_{y,2}^{k-j}\bigr)
       =z_1^{8-k}z_2^k.
\]
Thus \(\Delta\) is onto and
\(\dim\ker\Delta=25-9=16\).

Multiplication by \(q_{xy}\) injects \(V_3\otimes V_3\) into
\(V_4\otimes V_4\), since the polynomial ring is an integral
domain and \(q_{xy}\ne0\). Its image is contained in
\(\ker\Delta\), because \(q_{xy}(z,z)=0\). Its dimension is
\(4\cdot4=16\), so
\[
 \ker\Delta=q_{xy}(V_3\otimes V_3).
\]

We identify the orthogonal complement with the physical total-spin
space, using the inner products just fixed. Since \(\Delta\) is
onto, \(\Delta^*\) is injective. Unitarity and equivariance imply
that \(\Delta^*\) also intertwines the actions. Hence
\(\operatorname{Ran}\Delta^*=(\ker\Delta)^\perp\) is an
irreducible copy of \(V_8\), namely spin four.
In \(V_4\otimes V_4\), the weight of \(m_j\otimes m_k\) is
\(j+k-4\). Its weight-four subspace is one dimensional, spanned
by \(m_4\otimes m_4\). Every spin-four copy contains a nonzero
weight-four vector. Two such copies therefore intersect nontrivially,
and irreducibility forces them to coincide. Consequently
\(\operatorname{Ran}\Delta^*\) is the entire total-spin-four
space onto which \(P_e\) projects. Thus \(\ker P_e=\ker\Delta\).
Reversing the edge orientation changes the determinant by \(-1\)
and leaves its image unchanged.

Finally, expand a general \(\varphi\) in a tensor monomial basis
of the other sites. The equation \(P_e\varphi=0\) is equivalent
to the same equation for each endpoint coefficient. Each coefficient
is \(q_e\) times a polynomial of bidegree \((3,3)\), by the
two-site argument. Summing those quotients against the other-site
monomials gives exactly the stated \(\chi\). The converse follows
coefficient by coefficient as well.
\end{proof}

\begin{lemma}\label{l:divisibility}
Let \(L,N\ge3\), \(R\subset\Lambda_{L,N}\), and
\(E_R=\{e\in E_{L,N}:e\cap R\ne\varnothing\}\).
If \(\varphi\in\mathcal H_\Lambda\) and \(P_e\varphi=0\)
for every \(e\in E_R\), then
\(\varphi=(\prod_{e\in E_R}q_e)\chi\), where \(\chi\) is independent
of every \(z_x\), \(x\in R\).
\end{lemma}

\begin{proof}
The four sites \(x+(1,0),x-(1,0),x+(0,1),x-(0,1)\) are
different from \(x\) and pairwise distinct. Indeed, within a
coordinate the equality of the two neighbours would make its period
divide two; between coordinates an equality would make a period
divide one. Neither is possible for periods at least three.
Thus there are exactly four distinct incident edges at every site.

Work in the unique factorization domain
\(\mathcal R=\mathbb C[z_{u,1},z_{u,2}:u\in\Lambda]\).
For two distinct sites \(x,y\), regard
\(q_{xy}=z_{x,1}z_{y,2}-z_{x,2}z_{y,1}\) as a polynomial in
\(z_{y,2}\) over the polynomial ring in all the other variables.
Its two coefficients are relatively prime: the variable
\(z_{x,1}\) does not divide \(z_{x,2}z_{y,1}\).
It is primitive and is linear and nonconstant over the fraction
field of that coefficient ring. Gauss's lemma therefore makes it
irreducible in \(\mathcal R\). This argument already includes all
other site variables. Each \(q_e\), with either orientation, is
accordingly an irreducible and hence a prime element.

Two different unoriented edges have different endpoint sets.
Their determinants depend nontrivially on different sets of site
variables and cannot differ by a nonzero scalar, the only kind of
unit in \(\mathcal R\). They are therefore nonassociate.
Lemma~\ref{l:edge} says that every \(q_e\), \(e\in E_R\),
divides \(\varphi\). Unique factorization for these distinct prime
factors implies that their product, denoted \(b_R\), divides
\(\varphi\).

If \(\varphi=0\), take \(\chi=0\). Otherwise write
\(\varphi=b_R\chi\) and decompose \(\chi\) into its finitely
many multihomogeneous components. Write \(d_R(u)\) for the number
of edges of \(E_R\) at \(u\), so that \(b_R\) has degree
\(d_R(u)\) in \(z_u\). Multiplication by the nonzero \(b_R\)
is injective, and distinct multidegrees remain distinct after adding
\((d_R(u))_u\). Since \(\varphi\) has multidegree \((4)_u\),
the only possible component of \(\chi\) has degree
\(4-d_R(u)\) at each \(u\). For \(u\in R\), all four incident
edges are in \(E_R\), so this degree is zero. A polynomial of
degree zero in both variables at each such site is independent of
those variables. This also covers \(R=\varnothing\), where
\(b_R=1\) and the independence assertion is vacuous.
\end{proof}

\begin{proposition}\label{p:toruskernel}
For every \(L,N\ge3\), the torus is simple and has four distinct
neighbours at each site, \(B_{L,N}\ne0\), and
\(\ker H_{L,N}=\mathbb C B_{L,N}\).
Changing edge orientations does not change this one-dimensional space.
\end{proposition}

\begin{proof}
The neighbour calculation in Lemma~\ref{l:divisibility} proves
the asserted degree and the absence of loops, including at period
three; the edge set in Definition~\ref{d:model} has no repeated
unordered pairs. Each determinant is nonzero, so their product is
nonzero in the polynomial integral domain. Exactly four factors
meet each site; thus \(B_{L,N}\) has degree four there and is a
vector of \(\mathcal H_\Lambda\).
It is divisible by \(q_e\) for every \(e\), so
Lemma~\ref{l:edge} gives \(P_eB_{L,N}=0\).

For any \(\varphi\), orthogonality of the projections gives
\[
 \langle\varphi,H_{L,N}\varphi\rangle
       =\sum_e\|P_e\varphi\|^2.
\]
Consequently \(\ker H_{L,N}=\bigcap_e\ker P_e\).
Apply Lemma~\ref{l:divisibility} to a vector in this intersection
with \(R=\Lambda\). The quotient by \(B_{L,N}\) is independent
of every site variable and is a scalar. This proves both inclusions
in the kernel equality. Reversing any chosen orientations multiplies
\(B_{L,N}\) by \((-1)^k\), where \(k\) is the number reversed,
and therefore preserves its span.
\end{proof}

\begin{lemma}\label{l:star}
Let \(h_\star=\sum_{j=1}^4P_{\{0,j\}}\) on \(V_4^{\otimes5}\).
Its smallest positive eigenvalue \(g_\star\) exists and is positive.
For \(L,N\ge3\), set \(H_x=\sum_{e\ni x}P_e\) and
\(\mathcal S_x=\ker H_x\). Then
\[
 H_x\ge g_\star(I-\Pi_{\mathcal S_x}),\qquad
 \sum_xH_x=2H_{L,N}.
\]
\end{lemma}

\begin{proof}
The operator \(h_\star\) is positive semidefinite on the
\(5^5\)-dimensional Hilbert space. It is nonzero: the tensor
\((z_2^4)^{\otimes5}\) has each endpoint pair in the
one-dimensional weight-four subspace used in Lemma~\ref{l:edge},
so every one of the four projections fixes it. Its eigenvalue for
\(h_\star\) is four. The finite set of positive eigenvalues is
therefore nonempty and has a strictly positive minimum \(g_\star\).
The spectral theorem gives
\(h_\star\ge g_\star(I-\Pi_{\ker h_\star})\).

The centre \(x\) and its four neighbours are distinct, so relabelling
those five tensor factors identifies \(H_x\) with
\(h_\star\) tensored with the identity on the other sites.
Edges between the neighbours do not occur in \(H_x\) and do not
affect this identification. Tensoring the preceding inequality
proves the first assertion for \(H_x\). Each unoriented edge has
two different endpoints and occurs once at each endpoint in
\(\sum_xH_x\), proving the second assertion.
\end{proof}

\subsection{Finite frames and their separation}

The symbols \(D\), \(\phi\), \(r\), and \(\mathcal M\) have the
meaning fixed in Definition~\ref{d:moment}. In particular,
\(\phi=(1+\sqrt5)/2\), \(r^2=\phi+2\), and
\(\mathcal M\) is the uniform average on the twelve displayed
icosahedral vertices. We use its degree-four moment identities and
polynomial covariance from Lemma~\ref{l:moments}.

\begin{definition}\label{d:frame}
A frame is a finite set \(A\subset\mathbb S^2\) with positive weights
\(w(a)\) summing to one such that
\(\sum_{a\in A}w(a)p(a)=\mathcal M(p)\)
for every real polynomial \(p\) of degree at most four.
\end{definition}

The defining identity also holds for complex polynomials by applying
it to their real and imaginary parts.

\begin{lemma}\label{l:tight}
For every frame \((A,w)\) and every \(p\in V_4\),
\[
 \sum_{a\in A}5w(a)|p(\zeta(a))|^2=\|p\|_{V_4}^2.
\]
For site frames \((D_x,w_x)\), set \(\Omega=\prod_xD_x\),
\(\pi(s)=\prod_x5w_x(s_x)\), and
\(\iota\varphi(s)=\varphi((\zeta(s_x))_x)\).
Then \(\iota:\mathcal H_\Lambda\to\ell^2(\Omega,\pi)\) is an isometry.
\end{lemma}

\begin{proof}
Let \(m_j(z)=z_1^{4-j}z_2^j\), \(0\le j\le4\).
For \(j\le k\), the spinor projection identity gives
\[
 \overline{m_j(\zeta(a))}\,m_k(\zeta(a))
 =P_{jk}(a):=
 \frac{(1+a_3)^{4-k}(1-a_3)^j(a_1+i a_2)^{k-j}}{16}.
\]
Indeed, the paired factors are
\(|\zeta_1|^2=(1+a_3)/2\),
\(|\zeta_2|^2=(1-a_3)/2\), and
\(\overline{\zeta_1}\zeta_2=(a_1+i a_2)/2\).
For \(j>k\), use \(P_{jk}=\overline{P_{kj}}\).
Every \(P_{jk}\) has degree at most four. If \(j<k\), rotation
through \(\pi/(k-j)\) about the third coordinate multiplies
\(P_{jk}\) by \(-1\). Polynomial covariance in
Lemma~\ref{l:moments} implies \(\mathcal M(P_{jk})=0\).
The same follows by conjugation for \(j>k\).

For diagonal entries only the moments of \(t=a_3\) are needed:
\(\mathcal M(1)=1\), \(\mathcal M(t)=\mathcal M(t^3)=0\),
\(\mathcal M(t^2)=1/3\), and \(\mathcal M(t^4)=1/5\).
Expanding the five polynomials yields
\[
 \begin{array}{c|ccccc}
 j&0&1&2&3&4\\ \hline
 \mathcal M(P_{jj})&1/5&1/20&1/30&1/20&1/5
 \end{array}
 \qquad
 5\mathcal M(P_{jj})=\binom4j^{-1}.
\]
For example, the three different even parts of their numerators
are \(1+6t^2+t^4\), \(1-t^4\), and
\(1-2t^2+t^4\), respectively. The frame identity for these
degree-four polynomials now gives
\[
 \sum_{a\in A}5w(a)
       \overline{m_j(\zeta(a))}m_k(\zeta(a))
       =\begin{cases}\binom4j^{-1},&j=k,\\0,&j\ne k.\end{cases}
\]
Expanding \(p\) in the monomials proves its norm identity with
exactly the normalization of Definition~\ref{d:poly}.

For a finite collection of sites, the evaluation map is the tensor
product of the single-site evaluation maps. The product monomials
form an orthogonal basis, and the finite sums of their pairwise
inner products factor into the single-site sums just computed.
Their tensor inner products are therefore preserved, which proves
the isometry on all of \(\mathcal H_\Lambda\).
\end{proof}

The weight \(\pi\) in this lemma is a positive reference measure
of total mass \(5^{|\Lambda|}\), rather than a probability.
The use of rotations in the proof was solely on polynomials of
degree at most four.

\begin{lemma}\label{l:threeframes}
Let
\[
 T=\begin{pmatrix}99/101&-20/101&0\\20/101&99/101&0\\0&0&1\end{pmatrix},
 \quad D^{(j)}=T^jD\ (j=0,1,2),\quad
 b=\frac{5-\sqrt5}{505},\quad r_b=\frac b{2-b}.
\]
The matrix \(T\) is in \(\mathrm{SO}(3)\).
Each \(D^{(j)}\), with uniform weights \(1/12\), is an antipodal frame.
For \(i\ne j\), \(a\in D^{(i)}\), \(a'\in D^{(j)}\),
\[
 0<b<1,\qquad b\le1\pm a\cdot a'\le2-b,\qquad0<r_b<1.
\]
The common lower bound \(b\) is attained, and distinct frames are
disjoint even up to antipodal signs.
\end{lemma}

\begin{proof}
The equality \(99^2+20^2=101^2\) proves \(T^TT=I\) and
\(\det T=1\). The set \(D\) consists of twelve distinct unit
vectors and is antipodal by Lemma~\ref{l:moments}.
For a polynomial \(p\) of degree at most four,
\[
 \frac1{12}\sum_{a\in D^{(j)}}p(a)
       =\mathcal M(p\circ T^j)=\mathcal M(p).
\]
Here \(T^j\) is bijective and covariance applies on precisely
the stated degree range. Thus each rotated set is an antipodal
frame of twelve letters.

Write \(\vartheta\) for the rotation angle of \(T\), so that
\(\cos\vartheta=99/101\) and \(\sin\vartheta=20/101\).
We have
\[
 1-\cos\vartheta=\frac2{101},\qquad
 1-\cos(2\vartheta)=\frac{800}{10201}\ge\frac2{101},
 \qquad
 b=\frac{2}{101(\phi+2)}.
\]
The last identity follows by rationalizing
\(\phi+2=(5+\sqrt5)/2\). In particular \(0<b<1\).
By applying a common inverse rotation and, if needed, swapping
the two vectors, it suffices to bound
\(1-u\cdot T^k v\) for \(u,v\in D\) and \(k\in\{1,2\}\).

First let \(u=v\). Direct multiplication of the planar rotation
matrix gives
\[
 1-u\cdot T^ku
       =(1-\cos(k\vartheta))(u_1^2+u_2^2).
\]
The three possible values of \(u_1^2+u_2^2\), from the three
coordinate families of \(D\), are
\(1/r^2,1,\phi^2/r^2\). Since \(\phi>1\), their minimum is
\(1/r^2\). It follows that \(1-u\cdot T^ku\ge b\).
Equality holds for \(k=1\) and \(u=(0,1,\phi)/r\).

Next suppose \(u\ne v\). For vertices in the same coordinate
family, their scalar product has the form
\((\epsilon\epsilon'+\phi^2\delta\delta')/r^2\) with
\(\epsilon,\epsilon',\delta,\delta'\in\{\pm1\}\).
Apart from the equal-vertex case, its values are
\(-1\) and \(\pm\phi/r^2\). For vertices in different
families, the only shared nonzero coordinate has product
\(\pm\phi\), so the scalar product is \(\pm\phi/r^2\).
The identity \(\phi/r^2=1/\sqrt5\) thus proves
\[
 u\cdot v\in\{-1,1/\sqrt5,-1/\sqrt5\},\qquad
 |u-v|^2\ge2-\frac2{\sqrt5}>1.
\]
On the first two coordinates the operator norm of a rotation minus
the identity is \(\sqrt{2(1-\cos(k\vartheta))}\), while it is
zero on the third coordinate. Hence
\[
 \|T-I\|=\frac2{\sqrt{101}}\le\frac{40}{101},\qquad
 \|T^2-I\|=\frac{40}{101}<\frac25.
\]
The triangle inequality gives \(|u-T^kv|>3/5\), and therefore
\[
 1-u\cdot T^kv=\tfrac12|u-T^kv|^2>\frac9{50}>b.
\]
For the last inequality, \(\phi+2>2\) implies
\(b<1/101<9/50\). This proves \(1-a\cdot a'\ge b\)
in every case. Replace \(a'\) by \(-a'\), which belongs to the
same rotated frame, to obtain \(1+a\cdot a'\ge b\) as well.
These two lower bounds give both upper bounds \(2-b\).
The equality case already displayed attains \(b\); changing the
sign of its second vector also attains the bound with the other
choice of sign. An equality \(a=\pm a'\) would make one of
\(1\pm a\cdot a'\) zero, contradicting \(b>0\).
Finally, \(0<b<1\) implies \(0<b/(2-b)<1\).
\end{proof}

\begin{lemma}\label{l:colour}
For \(m\ge3\) and representatives \(0\le i<m\), put
\[
 a_m(i)=\begin{cases}2,&m\text{ odd and }i=m-1,\\i\bmod2,&\text{otherwise}.
 \end{cases}
\]
Then \(\kappa(i,j)=a_L(i)+a_N(j)\pmod3\) is a proper
three-colouring of \(\Lambda_{L,N}\) for every \(L,N\ge3\).
\end{lemma}

\begin{proof}
If \(m\) is even, the two colours zero and one alternate along
the cycle, and the wraparound edge joins the odd index \(m-1\)
to index zero. If \(m\) is odd, they alternate on indices
\(0,\ldots,m-2\); the last of these indices is odd and has
colour one. The remaining vertex \(m-1\) has colour two and is
adjacent to that vertex and to index zero, whose colour is zero.
This covers \(m=3\) as well. Thus every cycle edge changes
\(a_m\) by a nonzero element of \(\mathbb Z_3\).
A torus edge changes exactly one summand in the formula for
\(\kappa\), leaving the other fixed, so its endpoint colours
are different modulo three.
\end{proof}

\begin{lemma}\label{l:qmod}
For unit spinors \(\zeta,\zeta'\) with
\(\zeta\zeta^*=(I+a\cdot\sigma)/2\) and
\(\zeta'\zeta'^*=(I+a'\cdot\sigma)/2\),
\[
 |\zeta_1\zeta'_2-\zeta_2\zeta'_1|^2=(1-a\cdot a')/2 .
\]
\end{lemma}

\begin{proof}
Expansion of the two-by-two determinant gives the Gram identity
\[
 |\zeta_1\zeta'_2-\zeta_2\zeta'_1|^2
       =\|\zeta\|^2\|\zeta'\|^2-|\zeta^*\zeta'|^2
       =1-|\zeta^*\zeta'|^2.
\]
The displayed Pauli matrices satisfy
\(\operatorname{tr}\sigma_j=0\) and
\(\operatorname{tr}(\sigma_j\sigma_k)=2\delta_{jk}\).
Consequently
\[
 |\zeta^*\zeta'|^2
 =\operatorname{tr}(\zeta\zeta^*\zeta'\zeta'^*)
 =\frac14\operatorname{tr}\bigl((I+a\cdot\sigma)
                               (I+a'\cdot\sigma)\bigr)
 =\frac{1+a\cdot a'}2.
\]
Substitution proves the assertion.
\end{proof}

\subsection{Frame specifications and the quantum comparison}

\begin{definition}\label{d:framemodel}
Given any proper three-colouring \(\kappa\) of a plane or torus,
put \(D_x=D^{(\kappa(x))}\), with weights \(1/12\).
On the torus define
\[
 \mu^{\mathrm{af}}_{\kappa,L,N}(s)=
 \frac1{Z^{\mathrm{af}}}\prod_x\frac1{12}
       \prod_{\{x,y\}\in E_{L,N}}(1-s_x\cdot s_y).
\]
On the plane, for finite \(W\subset\mathbb Z^2\) and
\(\eta_z\in D_z\) at exterior neighbour sites, define
\[
 \mu^{+,\eta}_{\kappa,W}(n)=\frac1{Z^{+,\eta}_{\kappa,W}}
 \prod_{x\in W}\frac1{12}
 \prod_{\substack{\{x,y\}\subset W\\d(x,y)=1}}(1+n_x\cdot n_y)
 \prod_{\substack{x\in W,z\notin W\\d(x,z)=1}}(1+n_x\cdot\eta_z).
\]
Only the finitely many needed exterior letters are specified.
These laws have full support. Write \(E_U\) for conditioning off \(U\),
\(\operatorname{Var}_U f=E_U|f-E_Uf|^2\), and
\[
 \mathcal D_\mu(f)=\sum_x\mu(\operatorname{Var}_x f),\qquad
 \lambda(\mu)=\inf_{\operatorname{Var}_\mu f>0}
             \frac{\mathcal D_\mu(f)}{\operatorname{Var}_\mu f}.
\]
All generators have one update per physical site. A gap lower bound is
always equivalently a Poincare inequality, including one-point spaces.
\end{definition}

Indeed, every interaction factor in this definition lies in
\([b,2-b]\), by proper colouring and Lemma~\ref{l:threeframes}.
There are finitely many configurations, each of strictly positive
weight, so every normalizer is finite and positive. For
\(W=\varnothing\) the normalizer is one and the configuration
space has one point.
On a positive finite law the conditional expectation \(E_U\) is
the orthogonal projection onto the functions of the letters off
\(U\): this follows by summing on each fixed exterior fibre.
In particular,
\[
 \mu(\operatorname{Var}_U f)=\|(I-E_U)f\|_{L^2(\mu)}^2,
 \qquad \mathcal L_\mu=\sum_x(E_x-I).
\]
If \(\mathcal D_\mu(f)=0\), full support makes \(f\) constant
on every single-site fibre. Any two product configurations can be
joined by changing their coordinates one at a time, so \(f\) is
constant. The finite-dimensional spectral theorem consequently
gives \(\lambda(\mu)>0\) whenever the configuration space has
at least two points. On a one-point space every Poincare inequality
holds by the convention of Section~1.

We abbreviate the torus law by \(\mu^{\rm af}\) and the valence-bond
state by \(B\). Every evaluation of an incident determinant is
nonzero: its squared modulus is at least \(b/2\) by
Lemmas~\ref{l:threeframes} and~\ref{l:qmod}.

\begin{lemma}\label{l:quantumfibre}
Fix \(x\), write \(q_e=\epsilon_{e,x}q_{xy}\) for \(e=\{x,y\}\),
and set
\[
 c_x=\prod_{e\ni x}\epsilon_{e,x},\quad
 \beta_x(t;\sigma)=\prod_{y\sim x}q_{xy}(\zeta(t),\zeta(\sigma_y)),
 \quad \rho_x=c_x\prod_{e\not\ni x}q_e .
\]
All products are evaluated at the corresponding frame spinors.
Then \(\iota B=\beta_x\rho_x\), and if
\(\psi=(\prod_{e\ni x}q_e)\chi\) with \(\chi\) independent of \(z_x\),
then \(\iota\psi=c_x\beta_x\,\iota\chi\).
Let \(\mathcal P_x\) be the fibrewise orthogonal projection in
\(\ell^2(\Omega,\pi)\) onto the span of \(\beta_x\), and put
\(Uf=f\,\iota B/\|B\|\).
The map \(U:L^2(\mu^{\rm af})\to\ell^2(\Omega,\pi)\) is unitary,
\(UE_x=\mathcal P_xU\), and
\[
 \|(I-\mathcal P_x)\iota\varphi\|
            \le\|(I-\Pi_{\mathcal S_x})\varphi\|.
\]
\end{lemma}

\begin{proof}
Here \(\epsilon_{e,x}=1\) when the chosen orientation of \(e\)
starts at \(x\), and \(\epsilon_{e,x}=-1\) otherwise, so
\(c_x\in\{1,-1\}\). Separating the four incident edge factors
in the product for \(B\) gives \(\iota B=\beta_x\rho_x\).
Evaluation of a product is the product of its evaluations, which
also gives \(\iota\psi=c_x\beta_x\,\iota\chi\) in the stated
case. The notation \(\iota\chi\) in this equality means evaluation
of the residual polynomial; no assertion about an isometry for its
different degrees is used. Both \(\rho_x\) and this evaluation
are independent of the coordinate at \(x\).

We first verify the probability identity that is needed for
unitarity. Lemma~\ref{l:qmod} and the product reference weights give
\begin{equation}\label{e:frame-ground-weight}
 \pi(s)|\iota B(s)|^2
 =5^{|\Lambda|}2^{-|E_{L,N}|}
       \prod_x\frac1{12}
       \prod_{\{u,v\}\in E_{L,N}}(1-s_u\cdot s_v).
\end{equation}
By Lemma~\ref{l:tight}, summing the left side gives \(\|B\|^2\).
Thus
\(\|B\|^2=5^{|\Lambda|}2^{-|E_{L,N}|}Z^{\rm af}\), and
\begin{equation}\label{e:frame-ground-law}
 \mu^{\rm af}(s)=\frac{\pi(s)|\iota B(s)|^2}{\|B\|^2}.
\end{equation}
The product \(\iota B\) is nowhere zero. Equation
\eqref{e:frame-ground-law} proves that \(U\) preserves the norm,
and its inverse on the full finite function space is
\(h\mapsto\|B\|h/(\iota B)\). Therefore it is unitary.

Fix an exterior configuration \(s_{\Lambda\setminus\{x\}}\)
and let \(\sigma\) denote its neighbour letters. The conditional
probabilities, by either the factorization or
\eqref{e:frame-ground-law}, are
\[
 \mu^{\rm af}(t\mid s_{\Lambda\setminus\{x\}})
 =\frac{5w_x(t)|\beta_x(t;\sigma)|^2}
        {\sum_{u\in D_x}5w_x(u)|\beta_x(u;\sigma)|^2}.
\]
The denominator is strictly positive. On the same fibre the
orthogonal projection formula in \(\ell^2(D_x,5w_x)\) is
\[
 (\mathcal P_xh)(t,s_{\Lambda\setminus\{x\}})
 =\beta_x(t;\sigma)
  \frac{\sum_{u\in D_x}5w_x(u)\overline{\beta_x(u;\sigma)}
                         h(u,s_{\Lambda\setminus\{x\}})}
       {\sum_{u\in D_x}5w_x(u)|\beta_x(u;\sigma)|^2}.
\]
Since the outside factor of \(\pi\) is constant on that fibre,
these orthogonal fibre projections combine to an orthogonal
projection on \(\ell^2(\Omega,\pi)\).
Substitute \(h=Uf=\beta_x\rho_x f/\|B\|\) in the formula.
The result is \(\beta_x\rho_x(E_xf)/\|B\|=U(E_xf)\), proving
the intertwining.

If \(\psi\in\mathcal S_x\), positivity of its four projection
terms implies \(P_e\psi=0\) for all \(e\ni x\).
Lemma~\ref{l:divisibility}, with \(R=\{x\}\), gives the
factorization of \(\psi\) used in the statement with a quotient
independent of \(z_x\). Its evaluated fibre is consequently a
multiple of \(\beta_x\). This proves the inclusion
\(\iota\mathcal S_x\subset\operatorname{Ran}\mathcal P_x\).
It follows that
\[
 \begin{split}
 \|(I-\mathcal P_x)\iota\varphi\|
 &=\|(I-\mathcal P_x)\iota(I-\Pi_{\mathcal S_x})\varphi\|\\
 &\le\|\iota(I-\Pi_{\mathcal S_x})\varphi\|
  =\|(I-\Pi_{\mathcal S_x})\varphi\|.
 \end{split}
\]
The inequality is the contraction property of an orthogonal
projection, and the last equality is the physical-space isometry
of Lemma~\ref{l:tight}.
\end{proof}

\begin{theorem}\label{t:QD}
For every \(L,N\ge3\) and every proper torus three-colouring,
\[
 \mu^{\mathrm{af}}(s)=
 \frac{\pi(s)|\iota B_{L,N}(s)|^2}{\|B_{L,N}\|^2},\qquad
 H_{L,N}\big|_{B_{L,N}^{\perp}}\ge
       \frac{g_\star}{2}\lambda(\mu^{\mathrm{af}})I .
\]
In particular
\(\operatorname{gap}(H_{L,N})\ge(g_\star/2)\lambda(\mu^{\mathrm{af}})\).
\end{theorem}

\begin{proof}
The probability identity was proved in
\eqref{e:frame-ground-law}. For \(\varphi\perp B\), set
\(f=U^{-1}\iota\varphi\). Since
\(U1=\iota B/\|B\|\) and both maps preserve inner products,
\[
 \mu^{\rm af}(f)
 =\langle U1,\iota\varphi\rangle
 =\frac{\langle B,\varphi\rangle}{\|B\|}=0,\qquad
 \operatorname{Var}_{\mu^{\rm af}}f=\|\varphi\|^2.
\]
Using the Poincare definition of \(\lambda\), the intertwining
and norm inequality of Lemma~\ref{l:quantumfibre}, and
Lemma~\ref{l:star}, we obtain
\[
 \begin{split}
 \lambda(\mu^{\rm af})\|\varphi\|^2
 &\le\sum_x\|(I-E_x)f\|_{L^2(\mu^{\rm af})}^2\\
 &=\sum_x\|(I-\mathcal P_x)\iota\varphi\|_{\ell^2(\pi)}^2\\
 &\le\sum_x\langle\varphi,(I-\Pi_{\mathcal S_x})\varphi\rangle\\
 &\le\frac1{g_\star}\sum_x\langle\varphi,H_x\varphi\rangle
  =\frac2{g_\star}\langle\varphi,H_{L,N}\varphi\rangle.
 \end{split}
\]
For \(\varphi=0\) the same inequality is immediate.
The space \(B^\perp\) is invariant under \(H_{L,N}\), since
\(H_{L,N}\) is self-adjoint and kills \(B\).
Proposition~\ref{p:toruskernel} identifies its spectrum with the
positive spectrum of \(H_{L,N}\), proving the last assertion.
\end{proof}

\subsection{Finite energy and the gauge on a lifted square}

\begin{lemma}\label{l:planeenergy}
All the frame specifications above are consistent under conditioning
and have nearest-neighbour Markov dependence. Uniformly over proper
colourings, finite volumes, admissible frame-valued boundaries, and
fully specified one-site fibres, every letter has conditional
probability at least
\[
 \alpha=\frac1{12}r_b^4\in(0,1].
\]
The same lower bound holds after conditioning on fewer letters.
Any prescribed assignment on \(k\) sites, under arbitrary consistent
partial conditioning off those sites, has conditional probability
at least \(\alpha^k\).
For a finite plane block \(Q\), the single-site heat bath of every conditional block law in \(Q\) has Poincare gap at least
\(r_b^{|E(Q)|}\), where \(E(Q)\) is its internal edge set.
\end{lemma}

\begin{proof}
Let \(U\) be a subset of the unfixed sites in a finite specification
and fix the letters on its complement. Every site or edge factor
not meeting \(U\) is then constant and cancels between the
conditional weight and its normalizer. The remaining factors are
exactly the site factors on \(U\), the internal edge factors of
\(U\), and the crossing factors with the newly fixed neighbour
letters. They have the same sign of interaction as the initial
specification. This proves consistency, both for the torus law and
for the plane laws, and shows that only exterior nearest-neighbour
letters affect the conditional law. The assertion also holds for
\(U=\varnothing\), whose conditional space has one point.

For a fully specified fibre at a site \(x\), the unnormalized
weight of a proposed letter \(t\in D_x\) is \(1/12\) times
at most four factors, each in \([b,2-b]\). If their number is
\(d_x\le4\), the numerator is at least \(b^{d_x}/12\),
and the sum of all twelve weights is at most \((2-b)^{d_x}\).
Thus
\[
 \Pr(s_x=t\mid\hbox{all other unfixed letters})
 \ge\frac1{12}r_b^{d_x}\ge\frac1{12}r_b^4=\alpha.
\]
Lemma~\ref{l:threeframes} gives \(0<r_b<1\), hence
\(0<\alpha\le1\). Conditioning on fewer letters takes a finite
weighted average of the fully specified conditional probabilities;
every summand has this same lower bound.

More explicitly, let \(A=\{x_1,\ldots,x_k\}\) consist of
distinct sites and prescribe \(t_i\in D_{x_i}\). Condition on
any positive event specifying letters off \(A\). At stage \(i\)
of the chain rule, that event and the previous \(i-1\) prescribed
letters condition only off \(x_i\), so the conditional probability
of \(t_i\) is at least \(\alpha\). The preceding lower bounds
ensure that every intermediate event has positive probability.
Multiplying gives \(\alpha^k\). For \(k=0\), the probability
and the asserted bound are both one. The argument also applies to
a positive event determined by fewer outside letters, by averaging.

For the block-gap assertion, fix a finite plane block \(Q\) and
all its exterior neighbour letters \(\eta\). Absorb every fixed
boundary factor into its endpoint site weight by setting
\[
 h_x(t)=\frac1{12}
       \prod_{\substack{z\notin Q\\d(x,z)=1}}
                    (1+t\cdot\eta_z),\qquad
 \nu_x(t)=\frac{h_x(t)}{\sum_{u\in D_x}h_x(u)},\qquad
 \nu_Q=\bigotimes_{x\in Q}\nu_x.
\]
Every \(h_x(t)\) is positive, since all its boundary factors are
in \([b,2-b]\). Thus \(\nu_Q\) is a strictly positive product
probability on the finite product \(\prod_{x\in Q}D_x\). Let
\[
 R(n)=\prod_{\{x,y\}\in E(Q)}(1+n_x\cdot n_y),\qquad
 e=|E(Q)|.
\]
The conditional block law is \(d\mu=(R/\nu_Q(R))\,d\nu_Q\),
and \(b^e\le R\le(2-b)^e\). In particular,
\(0<\nu_Q(R)<\infty\), and its density has the bounds
\[
 0<m_Q:=\frac{b^e}{\nu_Q(R)}
 \ \le\ \frac{R(n)}{\nu_Q(R)}
 \ \le\ M_Q:=\frac{(2-b)^e}{\nu_Q(R)}<\infty.
\]
Apply the density comparison of Lemma~\ref{l:density}, proved in
Appendix~A, with its product reference law \(\pi=\nu_Q\),
density \(\varrho=R/\nu_Q(R)\), and constants \(m=m_Q\),
\(M=M_Q\). All its hypotheses have just been verified. Its
conditional-variance comparison and product Poincare inequality give,
for every complex \(f\),
\[
 \mathcal D_\mu(f)
 \ge m_Q\mathcal D_{\nu_Q}(f)
 \ge m_Q\operatorname{Var}_{\nu_Q}f
 \ge\frac{m_Q}{M_Q}\operatorname{Var}_\mu f.
\]
Consequently the Poincare gap is at least
\(m_Q/M_Q=r_b^e\). Only internal edges occur in this ratio;
all fixed boundary factors have been retained in \(\nu_Q\).
If some sites of \(Q\) are already fixed, apply this construction
to the free subset \(U\subset Q\). Since
\(E(U)\subset E(Q)\) and \(0<r_b<1\), the resulting bound
\(r_b^{|E(U)|}\) is at least the displayed bound for \(Q\).
If the free subset is empty, the Poincare inequality holds because
its variance is zero. In particular the argument includes empty
blocks and blocks with no internal edges.
\end{proof}

The finite-energy bound in this lemma concerns admissible frame
letters on neighbour sites. Arbitrary unit-vector pins, which are
used in the separate response estimates, are not included in its
quantifiers.

\begin{lemma}\label{l:boxlift}
Let \(Q'=a+\{-M,\ldots,M\}^2\subset\mathbb Z^2\), \(M\ge1\),
and \(L,N\ge2M+3\). Reduction modulo the periods is injective
on the \(\ell^\infty\) one-neighbourhood of \(Q'\), preserves exactly
the internal nearest-neighbour edges of \(Q'\), and identifies all its
crossing bonds with the torus crossing bonds of its image \(Q\).
For every torus outside configuration \(\sigma\), pull back the
colouring and set
\[
 n_z=(-1)^{z_1+z_2}s_{p(z)}\quad(z\in Q'),\qquad
 \eta_z=(-1)^{z_1+z_2}\sigma_{p(z)}
       \quad(z\notin Q',\ d(z,Q')=1).
\]
This coordinatewise bijection carries the conditional
antiferromagnetic law in \(Q\) to
\(\mu^{+,\eta}_{\kappa\circ p,Q'}\), intertwines every site update,
and preserves Hamming distances between configurations in each fibre.
\end{lemma}

\begin{proof}
Let \(p:\mathbb Z^2\to\mathbb Z_L\times\mathbb Z_N\) be
coordinatewise reduction. The indicated one-neighbourhood is
\[
 A'=a+\{-M-1,\ldots,M+1\}^2.
\]
The difference between two of its first coordinates has absolute
value at most \(2M+2<L\), and the analogous bound is
\(2M+2<N\) for second coordinates. If their reductions agree,
both differences are multiples of their periods, so both are zero.
This proves injectivity on \(A'\).

Plane nearest neighbours reduce to distinct torus nearest neighbours.
Conversely, if \(z,w\in Q'\) have nearest-neighbour reductions,
the congruent coordinate differences are \(\pm1\) and zero,
in one order. A difference in \([-2M,2M]\) congruent to zero
modulo its period must be zero. A difference in that interval
congruent to \(1\) or \(-1\) must be that very integer:
every other representative has absolute value at least
\(\min\{L,N\}-1\ge2M+2\). Thus \(d(z,w)=1\), which
proves the exact correspondence of internal edges.

For a torus edge leaving \(Q=p(Q')\), lift its endpoint in \(Q\)
uniquely to \(z\in Q'\). Its other endpoint has the form
\(p(z+v)\), with a unique
\(v\in\{(1,0),(-1,0),(0,1),(0,-1)\}\); the four torus
neighbours are distinct. Then \(z+v\in A'\), and it is outside
\(Q'\) since the torus endpoint is outside \(Q\).
Conversely, a plane crossing bond from \(z\in Q'\) to
\(w\notin Q'\) has \(w\in A'\); injectivity on \(A'\)
ensures \(p(w)\notin Q\). These operations are inverse on
crossing bonds. They also show that all required exterior sites
are represented without identifications.

The function \(\kappa\circ p\) is a proper plane colouring:
each plane nearest-neighbour edge maps to a torus edge whose
colours differ. Put \(\varepsilon_z=(-1)^{z_1+z_2}\) for the
plane sites involved in the lift. Each alphabet is antipodal with
uniform weights, so multiplication by \(\varepsilon_z\) is a
weight-preserving bijection on its alphabet. The displayed \(\eta\)
are consequently admissible exterior frame letters.
For a plane internal or crossing bond \(\{z,w\}\),
\(\varepsilon_z\varepsilon_w=-1\). Substituting
\(s_{p(z)}=\varepsilon_z n_z\) at a free site and
\(\sigma_{p(w)}=\varepsilon_w\eta_w\) at a fixed site changes
its torus factor as follows:
\[
 1-s_{p(z)}\cdot s_{p(w)}=1+n_z\cdot n_w,
 \qquad
 1-s_{p(z)}\cdot\sigma_{p(w)}=1+n_z\cdot\eta_w.
\]
Factors wholly outside \(Q\) cancel in the conditional torus
law. The exact internal and crossing-bond correspondences show
that all remaining factors are precisely those in
\(\mu^{+,\eta}_{\kappa\circ p,Q'}\). Site weights are unchanged,
and the configuration bijection makes their normalizing sums
equal. This proves the asserted transport of laws.

For \(x=p(z)\), let \(Jf(s)=f((\varepsilon_w s_{p(w)})_{w\in Q'})\)
be pullback on the corresponding conditional function spaces.
It is unitary by the transported law. In a single-site fibre all
other coordinates are fixed, and summing over the letter at \(x\)
is exactly summing over its image letter at \(z\) with the
transported conditional probabilities. Therefore
\(E_x^{\rm af}J=J E_z^+\). This proves the update intertwining.
Finally \(s_{p(z)}=t_{p(z)}\) if and only if
\(\varepsilon_zs_{p(z)}=\varepsilon_zt_{p(z)}\), so the
coordinate bijection preserves the number of unequal coordinates.
All signs used here are defined on plane lifts; no sign function
on the torus is needed.
\end{proof}

\section{Physical windows and the periodic gap}
\label{s:periodic-gap}

We pass from a uniform estimate on plane window marginals to a gap on
the physical torus.  All shell sets below consist of physical sites.
The finite coupling and Poincare tools used in this section are proved
in Appendix~\ref{a:finite}.  In particular, all heat-bath sums count
each indicated update at rate one.

\begin{definition}\label{d:shells}
For \(a\in\mathbb Z^2\), integer \(u\ge1\), put
\(C_u(a)=a+\{-u,\ldots,u\}^2\).
For \(\operatorname{dist}_\infty(y,C_u(a))=1\), let
\(s=u^{1/4}\), \(k_i=\lfloor is\rfloor\), \(i_1=\lfloor s\rfloor+1\).
For an integer \(l_0\ge1\) define
\[
 \begin{split}
 F_i&=\{x\in C_u(a):|x-y|_\infty>(i-1)s\},\\
 D_i&=\{x\in C_u(a):|x-y|_\infty=k_i\},\\
 A_i&=\{x\in D_i:\operatorname{dist}_\infty
                    (x,\mathbb Z^2\setminus C_u(a))\le l_0\},\qquad
 B_i=D_i\setminus A_i .
 \end{split}
\]
Let \(I_0\) consist of the integers \(2\le i\le i_1\) for which
some point of \(D_i\) is at distance strictly less than \(s/3\)
from a corner site of \(C_u(a)\).
All these are sets of physical sites.
\end{definition}

\begin{definition}[(W)]\label{d:W}
The family of plane frame specifications has \(\mathsf W(l_0,p_0,M_\star)\)
if \(l_0,M_\star\) are positive integers, \(0<p_0\le1\), and,
for every proper plane three-colouring, every translate \(C_u(a)\)
with integer \(u\ge M_\star\), every \(y\) as in Definition~\ref{d:shells},
every integer \(2\le i\le u^{1/4}\) outside \(I_0\), and every two
admissible frame-valued boundary conditions \(\eta,\xi\) for \(F_i\),
\[
 d_{\mathrm{TV}}\!\left(
       \mu^{+,\eta}_{\kappa,F_i,B_i},
       \mu^{+,\xi}_{\kappa,F_i,B_i}\right)\le1-p_0 .
\]
The subscript \(B_i\) denotes a marginal. Boundaries are arbitrary
within their site alphabets; they need not agree anywhere.
\end{definition}

For a positive specification with other finite site alphabets, the
phrase ``window bound (W)'' means the displayed marginal inequality
with that specification's laws in place of the frame laws, with the
same quantifiers over translates, exterior sites and permitted indices.
Consistency means that conditioning a volume law on its letters off a
subset gives the specification law on that subset with the induced
boundary.  Range one means that the latter law depends only on exterior
letters within \(\ell^\infty\) distance one of the subset.

\begin{lemma}\label{l:stopped}
Suppose a consistent positive plane specification with site-dependent
finite alphabets has range one in \(\ell^\infty\), one-site probabilities
at least \(\alpha\in(0,1]\), and the window bound (W) at every translate.
Put \(\delta=p_0\alpha^{16l_0}\) and
\(M_{\mathrm g}(l_0)=\max\{16,6(l_0+1)\}^4\).
For \(u\ge\max\{M_\star,M_{\mathrm g}(l_0)\}\), boundaries of \(C_u(a)\)
agreeing off one exterior site \(y\), and any integer
\(4\le h\le\lfloor u^{1/4}\rfloor\), there is a coupling of the
full laws such that
\[
 \mathbb P(n_{Q_h}\ne n'_{Q_h})\le(1-\delta)^{h-3},
 \qquad Q_h=\{x\in C_u(a):|x-y|_\infty>hu^{1/4}\}.
\]
If \(y\) is farther than range one, the full laws are equal.
\end{lemma}

\begin{proof}
Write \(C=C_u(a)\), \(\mu=\mu_C^\eta\) and
\(\mu'=\mu_C^\xi\) for the two volume laws.  If the exterior site
\(y\) is farther than distance one from \(C\), the range-one property
gives \(\mu=\mu'\); sampling the same configuration in both
components proves the conclusion.  Hence assume
\(\operatorname{dist}_\infty(y,C)=1\), and use the shells of
Definition~\ref{d:shells}.  The size assumption gives
\(s=u^{1/4}\ge\max\{16,6(l_0+1)\}\).
Lemma~\ref{l:shellgeom} supplies
\[
 |I_0|\le1,\qquad D_i\subset F_i,
 \qquad |A_i|\le8l_0\quad(i\notin I_0),
 \qquad 2\le i\le i_1.
\]
It also says that every site of \(C\setminus F_i\) within distance
one of \(F_i\) belongs to \(D_{i-1}\), and that \(y\) is not
within distance one of \(F_i\), for \(i\ge2\).
Since \(h+1\le\lfloor s\rfloor+1=i_1\), these geometric statements
apply to every surgery index used below.

We first record the consequence of finite energy needed for the
surgeries.  In any one of the positive finite-volume laws under
consideration, a one-site conditional probability after specifying
only some other letters is an average of fully specified one-site
conditional probabilities.  Thus it is at least \(\alpha\).
Order any finite set \(A\) and apply the finite chain rule to a
prescribed assignment \(t_A\), conditional on any fixed letters off
\(A\).  Every factor is at least \(\alpha\), so that
\[
 \mathbb P(n_A=t_A\mid\text{those fixed letters})
                      \ge\alpha^{|A|}.
\]
Positivity guarantees every conditioning denominator used in this
argument is nonzero.  For \(A=\varnothing\) the assertion reads
\(1\ge1\).

Start from the product coupling \(\mu\otimes\mu'\).  For the fixed
indices \(i=2,\ldots,h+1\), in increasing order, use the following
operation.  If \(i\in I_0\), leave the pair unchanged.  Otherwise,
for the current pair \((\omega,\omega')\), let \(p_i,p'_i\) be its
two specification laws on \(F_i\), with the respective complements
and original boundaries fixed.  If
\(\omega_{D_{i-1}}=\omega'_{D_{i-1}}\), these laws coincide:
all their exterior neighbours inside \(C\) belong to the agreeing
separator \(D_{i-1}\), and their exterior neighbours outside \(C\)
have identical boundary letters because none is \(y\).  In this
case draw one sample from the common law and use it in both components
on \(F_i\).

If the two configurations on \(D_{i-1}\) differ, maximally couple
the \(B_i\) marginals of \(p_i,p'_i\), using
Lemma~\ref{l:finitecoupling}.  Given the two sampled marginals,
complete them independently to \(F_i\) using their respective
conditional laws on \(F_i\setminus B_i\).  Finally resample \(A_i\)
independently in the two components, each with its conditional law
given the rest of its own configuration.  The initial completion has
the correct two marginals \(p_i,p'_i\).  The last resampling preserves
each of these laws, by Lemma~\ref{l:CE} and consistency.  Thus in
either branch the output on \(F_i\), for every fixed input pair,
has marginals \(p_i,p'_i\).  Lemma~\ref{l:surgery} shows that every
step preserves the original full marginals \(\mu,\mu'\).

Suppose now \(i\le h\) and \(i\notin I_0\).
Then \(i\le h\le\lfloor s\rfloor\), so (W) implies agreement
on \(B_i\) with probability at least \(p_0\) in the maximal-coupling
branch.  Fix an arbitrary assignment \(t_{A_i}\) using the alphabet
at each of its sites.  Given the entire result before the final
independent resamplings, the probability that both choose this
assignment is at least \(\alpha^{2|A_i|}\).
Those resamplings leave \(B_i\) fixed.  Therefore, conditional on
any history preceding step \(i\),
\[
 \mathbb P(\omega_{D_i}=\omega'_{D_i}
                     \text{ after step }i\mid\text{history})
       \ge p_0\alpha^{2|A_i|}
       \ge p_0\alpha^{16l_0}=\delta.
\]
In the identical-resampling branch the probability is one because
\(D_i\subset F_i\), so this bound holds in that branch as well.
At step \(h+1\), if \(h+1>s\), the coupling construction still
exists but no success bound from (W) is used.

Let
\[
 J=\{i\in\{2,\ldots,h\}:i\notin I_0,
                                      \ i+1\notin I_0\}.
\]
There are \(h-1\) possible trial indices, and the at most one
exceptional index excludes at most two of them.  Hence
\(|J|\ge h-3\ge1\).
If a trial \(i\in J\) produces agreement on \(D_i\), the
unskipped next step \(i+1\) uses identical resampling, giving agreement
on \(F_{i+1}\).  This agreement persists.  Indeed, for a later index
\(j\ge i+2\), \(F_j\subset F_{i+1}\), and
\[
 \lfloor(j-1)s\rfloor\ge(j-1)s-1
                      \ge(i+1)s-1>is,
\]
so \(D_{j-1}\subset F_{i+1}\).
A skipped step changes nothing, and any unskipped step therefore
again uses identical resampling on \(F_j\), preserving agreement
everywhere in \(F_{i+1}\).  Since
\(Q_h=F_{h+1}\subset F_{i+1}\), one successful trial in \(J\)
ensures final agreement on \(Q_h\).

To bound the chance of no success, list \(J=\{j_1<\cdots<j_r\}\)
and let \(H_i\) denote the full history through step \(i\).
The event that the first \(k-1\) trials fail is measurable with
respect to \(H_{j_k-1}\).  The displayed conditional success bound
therefore yields, successively,
\[
 \mathbb P(\text{first }k\text{ trials fail})
       \le(1-\delta)
           \mathbb P(\text{first }k-1\text{ trials fail}).
\]
Induction gives probability at most \((1-\delta)^r\), which is
at most \((1-\delta)^{h-3}\).  This proves the required bound for
the final coupling of the full laws.  If \(\delta=1\), the first
trial fails with probability zero, so the same argument gives zero.
An empty \(B_i\) has its unique empty configuration and is coupled
with agreement probability one; an empty \(A_i\) needs no resampling.
Empty volumes occurring in a step and singleton alphabets are covered
by the same formulas.  No comparison of letters at distinct sites
has been made.
\end{proof}

\begin{corollary}\label{p:surgery}
Under the hypotheses and size threshold of Lemma~\ref{l:stopped},
the same two laws have marginal total variation at most
\((1-\delta)^{u^{1/4}-4}\) on
\(\{x\in C_u(a):|x-y|_\infty>u^{1/2}\}\).
When \(\delta=1\), this upper bound is zero.
\end{corollary}

\begin{proof}
If the changed exterior site is farther than range one, the full laws
are equal by Lemma~\ref{l:stopped}.  Otherwise set
\(s=u^{1/4}\) and \(h=\lfloor s\rfloor\).
The size threshold gives \(s\ge16\), so \(h\ge4\).
Since \(hs\le s^2=u^{1/2}\), the target set of this corollary
is contained in \(Q_h\).  Lemmas~\ref{l:stopped}
and~\ref{l:finitecoupling} thus bound its marginal total variation by
\((1-\delta)^{h-3}\).
For \(0<\delta<1\), the inequality
\(\lfloor s\rfloor-3\ge s-4\) implies
\((1-\delta)^{h-3}\le(1-\delta)^{s-4}\).
If \(\delta=1\), the coupling already gives agreement with
probability one and \(s-4>0\), so the asserted upper bound is zero.
\end{proof}

In the following lemma, ``as above'' retains all the specification
hypotheses of Lemma~\ref{l:stopped}: consistency, full support on
nonempty finite site alphabets, range one in \(\ell^\infty\), a
uniform lower bound \(\alpha\in(0,1]\) on fully conditional
one-site probabilities, and (W) at every translate.

\begin{lemma}\label{l:hamming}
Assume (W) and \(\alpha\in(0,1]\) as above, and put
\(d_0=\min\{1/2,p_0\alpha^{16l_0}\}\).
Choose an integer \(n\) satisfying
\[
 n\ge44,\quad n^8\ge M_\star,\quad
 n^2\ge\max\{50000,6(l_0+1)\},\quad
 n^{12}d_0^{20}\ge50000\,2^{20}\,20!.
\]
Such an integer exists. Set \(M=n^8\).
Every translated \(C_M(a)\), with any two admissible boundaries
agreeing off one exterior site, admits a full-law coupling with
expected Hamming distance at most \(M/1000\).
One sufficient bound is
\[
 25n^6+25n^{16}(1-d_0)^{n-3}\le M/1000.
\]
\end{lemma}

\begin{proof}
The number \(d_0\) is strictly positive and at most \(1/2\).
All four requirements on \(n\) are lower bounds by finite constants;
for example, an integer at least
\[
 \max\left\{44,\ M_\star^{1/8},\
       \sqrt{\max\{50000,6(l_0+1)\}},\
       \left(50000\,2^{20}\,20!\,d_0^{-20}\right)^{1/12}\right\}
\]
satisfies them.  This proves existence before any torus periods or
boundary conditions are chosen.

We first prove a finite tail inequality.  For \(0<d\le1/2\) and
integer \(n\ge44\), put \(r=n-3\).  Then
\((1-d)^{-1}\ge1+d\), and the binomial theorem gives
\[
 \begin{split}
 (1-d)^{-r}
 &\ge(1+d)^r\ge\binom r{20}d^{20}\\
 &\ge\frac{(r-19)^{20}}{20!}d^{20}
   =\frac{(n-22)^{20}}{20!}d^{20}
   \ge\frac{n^{20}d^{20}}{2^{20}20!}.
 \end{split}
\]
Here \(r\ge41\), and each of the twenty factors in
\(r(r-1)\cdots(r-19)\) is at least \(r-19>0\).
Taking reciprocals proves
\[
 (1-d)^{n-3}\le\frac{2^{20}20!}{d^{20}n^{20}}.
\]

Fix a translated square with radius \(M=n^8\).  Its shell parameter
is \(s=M^{1/4}=n^2\), which is at least both \(16\) and
\(6(l_0+1)\).  Thus \(M\ge M_{\mathrm g}(l_0)\), and the other
size assumption gives \(M\ge M_\star\).
Apply Lemma~\ref{l:stopped} with \(h=n\); its index restriction
holds because \(4\le n\le n^2=\lfloor s\rfloor\).
If the changed exterior site is farther than distance one, use the
identical coupling and there is nothing to prove.  Otherwise the
lemma gives a full-law coupling agreeing outside the near region
\(\{x:|x-y|_\infty\le n^3\}\), except on an event of probability
at most \((1-\delta)^{n-3}\le(1-d_0)^{n-3}\).
The near region contains at most
\((2n^3+1)^2\le25n^6\) sites, and the square has
\((2n^8+1)^2\le25n^{16}\) sites.  Consequently the expected
Hamming cost of this same coupling is at most
\[
 25n^6+25n^{16}(1-d_0)^{n-3}.
\]
The condition \(n^2\ge50000\) gives \(25n^6\le n^8/2000\).
Using the tail inequality and the last requirement on \(n\) gives
\[
 25n^{16}(1-d_0)^{n-3}
   \le\frac{25\,2^{20}20!}{d_0^{20}n^4}
   \le\frac{n^8}{2000}.
\]
Adding proves the asserted bound \(M/1000\), uniformly over every
translate and boundary pair.  The truncation in the definition of
\(d_0\) also covers \(\delta=1\) without division by \(1-\delta\).
\end{proof}

In the following criterion the square radius and both torus periods
are integers.  The torus \(\ell^\infty\) distance is the minimum
over coordinate lifts.

\begin{theorem}\label{t:periodic}
Let \(M\ge1\), \(L,N\ge2M+3\), and let \(\mu>0\) be a
probability on a finite product over \(\Lambda_{L,N}\) whose
conditional law on any set depends only on letters in its exterior
\(\ell^\infty\) one-neighbourhood.
For every translated square \(Q\) of side \(2M+1\), suppose that:
the single-site heat bath of every conditional law on \(Q\) has Poincare gap at least
\(g>0\); and any two conditional block laws whose outside
configurations agree off one site admit a coupling with expected
Hamming distance in \(Q\) at most \(K\ge0\), whenever that site is
outside and at \(\ell^\infty\) distance one from \(Q\).
Put \(T_M=(2M+1)^2\), \(S_M=8M+8\), and suppose \(S_MK<T_M\).
Then
\[
 \sum_{Q}(I-E_Q)\text{ has Poincare gap at least }T_M-S_MK,\qquad
 \lambda(\mu)\ge g\left(1-\frac{S_MK}{T_M}\right).
\]
The sum is over all torus translates, each at rate one.
In particular, \(K\le M/1000\) implies
\(\lambda(\mu)\ge(499/500)g\).
\end{theorem}

\begin{proof}
For a centre label \(a\in\Lambda_{L,N}\), let
\(Q_a=a+\{-M,\ldots,M\}^2\), with coordinates reduced modulo
the periods.  There are \(A=LN\) labels.  We count these labels
without making any assumption on symmetries of \(\mu\).

Fix a site \(x\).  The centre labels for which \(x\in Q_a\) are
exactly \(a=x-z\), where \(z\in\{-M,\ldots,M\}^2\).
Reduction is injective on this offset square, so their number is
\(T_M=(2M+1)^2\).  Likewise, the labels with
\(\operatorname{dist}_\infty(x,Q_a)\le1\), using cyclic coordinate
distances on the torus, are exactly \(a=x-z\) with
\(z\in\{-M-1,\ldots,M+1\}^2\).  This follows by adding an
offset in \(\{-1,0,1\}^2\) to an offset in the smaller square;
their sum is precisely the larger square.  Its reduction is still
injective: a coordinate difference between two of its points has
absolute value at most \(2M+2\), strictly less than either relevant
period.  Hence exactly
\[
 (2M+3)^2-(2M+1)^2=8M+8=S_M
\]
labels exclude \(x\) but have it at distance one.  This count
includes the four corner offsets, which may have no interaction with
\(Q_a\) for a nearest-neighbour specification.  In particular, the
maximal site multiplicity of the labelled block family is \(T_M\).

Set \(P=A^{-1}\sum_aE_{Q_a}\).  Consider two input configurations
differing only at site \(x\), and choose the same uniformly random
block label in both components.  If \(x\in Q_a\), their exterior
configurations agree, so their conditional block laws are equal.
Use an identical sample on \(Q_a\); the output pair then agrees
everywhere.  If \(x\notin Q_a\) but its distance from \(Q_a\)
is one, use the assumed coupling on \(Q_a\).  The original
disagreement remains at \(x\), and its expected number of additional
disagreements is at most \(K\); the total is at most \(1+K\).
For every remaining label, the range-one assumption makes the two
block laws equal.  Identical resampling then leaves just the one
disagreement at \(x\).  The resulting mixture is a coupling of
the two \(P\)-output laws with
\[
 \mathbb E d_H
    \le\frac{S_M(1+K)+(A-T_M-S_M)}A
    =1-\frac{T_M-S_MK}A.
\]
By assumption \(c=T_M-S_MK>0\), and
\(c\le T_M\le A\) since \(K\ge0\).
Lemma~\ref{l:contraction} proves a block Poincare gap at least \(c\).
Applying Lemma~\ref{l:comparison} with the conditional block gap
\(g\) and multiplicity \(t_{\max}=T_M\) yields
\[
 \lambda(\mu)\ge\frac{cg}{T_M}
             =g\left(1-\frac{S_MK}{T_M}\right).
\]
These are the two claimed bounds with unnormalized block rates.
They also hold on one-point product spaces by the Poincare convention.
Finally, if \(K\le M/1000\), then
\[
 \frac{S_MK}{T_M}
 \le\frac{8M(M+1)}{1000(2M+1)^2}
 =\frac1{500}\left(1-\frac1{(2M+1)^2}\right)
 <\frac1{500}.
\]
In particular the strict block-contraction assumption is satisfied,
and the final stated lower bound follows.
\end{proof}



\begin{theorem}\label{t:mainW}
If the plane frame family satisfies (W), there exist \(c>0\) and
an integer \(L_0\ge3\) such that
\(\ker H_{L,N}=\mathbb C B_{L,N}\) and every nonzero eigenvalue is
at least \(c\), for all integers \(L,N\ge L_0\).
\end{theorem}

\begin{proof}
Fix constants \((l_0,p_0,M_\star)\) witnessing (W).
Lemma~\ref{l:planeenergy} supplies a consistent positive,
nearest-neighbour plane family with a common one-site bound
\(\alpha=12^{-1}r_b^4\in(0,1]\).  Choose \(n\) by
Lemma~\ref{l:hamming}, once and for all, and put
\[
 M=n^8,\qquad L_0=2M+3,\qquad
 g_M=r_b^{\,2(2M+1)(2M)}>0.
\]
These choices precede the torus periods, colouring, block placement
and exterior fibre.

Let \(L,N\ge L_0\).  In particular \(L,N\ge3\), so
Lemma~\ref{l:colour} provides the proper torus three-colouring
\[
 \kappa(i,j)=a_L(i)+a_N(j)\pmod3,\qquad
 a_m(i)=\begin{cases}
 2,&m\text{ odd and }i=m-1,\\
 i\bmod2,&\text{otherwise},
 \end{cases}
 \quad 0\le i<m.
\]
Consider
\(\mu=\mu^{\mathrm{af}}_{\kappa,L,N}\) from
Definition~\ref{d:framemodel}.  This measure has full support and
nearest-neighbour conditional dependence by
Lemma~\ref{l:planeenergy}.  The periodic pullback
\(\bar\kappa(z)=\kappa(p(z))\), where
\(p:\mathbb Z^2\to\Lambda_{L,N}\) is reduction, is a proper plane
colouring.  Indeed each plane nearest-neighbour edge reduces to a
torus edge between distinct vertices since \(L,N\ge3\).

For every labelled square \(Q\), choose its lift
\(Q'=a+\{-M,\ldots,M\}^2\) to the plane, and put
\(Q'^+=a+\{-M-1,\ldots,M+1\}^2\).
Here \(M=n^8\ge1\) and \(L,N\ge2M+3\), precisely the size
hypotheses of Lemma~\ref{l:boxlift}.  It gives injectivity of
\(p\) on \(Q'^+\) and the exact internal and crossing-bond
correspondences.  For an exterior torus configuration \(\sigma\),
the local sign change in that lemma is
\[
 \varepsilon_z=(-1)^{z_1+z_2},\qquad
 n_z=\varepsilon_zs_{p(z)}\quad(z\in Q'),\qquad
 \eta_z=\varepsilon_z\sigma_{p(z)}
                                  \quad(z\in Q'^+\setminus Q').
\]
Injectivity ensures \(p(z)\notin Q\) for
\(z\in Q'^+\setminus Q'\), so the displayed boundary values
are defined.  Only nearest-neighbour exterior letters enter the law.
Lemma~\ref{l:boxlift} identifies the conditional torus law under
this map with \(\mu^{+,\eta}_{\bar\kappa,Q'}\).
Antipodality of each alphabet
ensures that the sign change is a bijection using the same site
alphabets.  This bijection intertwines single-site updates and
preserves Hamming distances.  It also preserves variances and
Dirichlet forms, by the change of variables in the finite sums.
The number of internal edges of \(Q'\) is
\(2(2M+1)(2M)\): there are \(2M\) bonds in each of its
\(2M+1\) horizontal rows and the same number vertically.
Lemma~\ref{l:planeenergy}, equivalently the density comparison in
Lemma~\ref{l:density}, therefore gives conditional single-site
Poincare gap at least \(g_M\), uniformly over all such fibres.

Now take two torus exterior configurations for \(Q\) differing at
just one site \(y\) at \(\ell^\infty\) distance one.
The image \(p(Q'^+)\) is the full torus one-neighbourhood of \(Q\):
every one-step coordinate displacement lifts to
\(\{-1,0,1\}^2\), and
\(Q'^+=Q'+\{-1,0,1\}^2\).
Thus the one-neighbourhood injection in Lemma~\ref{l:boxlift}
gives an existing, unique exterior lift of \(y\).
Under the same local
sign change the two plane boundary assignments differ at most at
that lift.  Only nearest-neighbour exterior letters are needed; if
necessary, extend the other exterior letters identically and
arbitrarily in their site alphabets.  For a corner site that has no
crossing nearest-neighbour bond the two conditional laws are already
equal.  In all other cases Lemma~\ref{l:hamming}, applied to this
plane colouring and placement, provides a coupling of expected
Hamming cost at most \(K=M/1000\).  The inverse coordinatewise
bijection carries it to a coupling of the two conditional torus laws
with exactly the same Hamming cost.

All hypotheses of Theorem~\ref{t:periodic} now hold, giving
\[
 \lambda(\mu^{\mathrm{af}}_{\kappa,L,N})
                        \ge\frac{499}{500}g_M.
\]
Use the constant defined and proved positive in Lemma~\ref{l:star}:
\(g_\star\) is the smallest positive eigenvalue of
\(h_\star=\sum_{j=1}^4P_{\{0,j\}}\) on \(V_4^{\otimes5}\).
It is independent of the torus periods.  Theorem~\ref{t:QD} applies
to every \(L,N\ge3\) and every proper torus three-colouring using
the frames of Definition~\ref{d:framemodel}; these hypotheses hold
for the periods and colouring just chosen.  Its precise operator
inequality is
\[
 H_{L,N}\big|_{B_{L,N}^{\perp}}
       \ge\frac{g_\star}{2}
               \lambda(\mu^{\mathrm{af}}_{\kappa,L,N})I.
\]
The factor \(1/2\) uses Lemma~\ref{l:star}:
\(\sum_xH_x=2H_{L,N}\), where \(H_x=\sum_{e\ni x}P_e\),
with one heat-bath update per site.  Consequently every eigenvalue on
\(B_{L,N}^{\perp}\) is at least
\[
 c=\frac{g_\star}{2}\,\frac{499}{500}g_M
                         =\frac{499}{1000}g_\star g_M>0.
\]
Proposition~\ref{p:toruskernel}, whose hypothesis is only
\(L,N\ge3\), supplies both \(B_{L,N}\ne0\) and
\(\ker H_{L,N}=\mathbb C B_{L,N}\).  Thus \(B_{L,N}^{\perp}\)
is the direct sum of the eigenspaces for nonzero eigenvalues.
The constants \(c,L_0\) are independent of \(L,N\).
The sign change was defined separately on each lifted square, so this
argument applies to both even and odd periods.
\end{proof}


We now prove the two conclusions of Theorem~\ref{t:main}, stated in
the result section: nonvanishing of \(B_{L,N}\) and the kernel
identity for every \(L,N\ge3\), and the existence of one
\(c>0\) and integer \(L_0\ge3\) bounding every nonzero eigenvalue
below by \(c\) for all integers \(L,N\ge L_0\).

\begin{proof}[Proof of Theorem~\ref{t:main}]
Proposition~\ref{p:toruskernel} proves \(B_{L,N}\ne0\) and
\(\ker H_{L,N}=\mathbb C B_{L,N}\) for every \(L,N\ge3\).
For the uniform positive gap, Theorem~\ref{t:inputs} gives (H1),
(H2) with \(C_V=10\) and
\(0<q=1119948287353/1120437614225<1\), and (H3\(_h\)), all at
the same pair \((c_{\rm s},c_{\rm t})=(961/1000,187/200)\).
Theorem~\ref{t:rec} thus gives \(A\ge2\), \(0<\theta<1\) and
\(\mathsf P(A,\theta)\).
Theorem~\ref{t:cutMOS} supplies
\(\mathsf W(64,p_0,M_\star)\) with common
\(p_0\in(0,1]\) and positive integer \(M_\star\), valid for
every proper plane colouring and every translate.
Theorem~\ref{t:mainW} then gives \(c>0\) and an integer
\(L_0\ge3\) such that all nonzero eigenvalues are at least \(c\)
whenever \(L,N\ge L_0\), as required.
\end{proof}

\section{Physical windows from a separating cut}\label{s:cut}

We use the shells of Definition~\ref{d:shells}, the plane
specifications of Definition~\ref{d:framemodel}, and the original-leg
partition sums of Definition~\ref{d:pinmodel}.  In a sum with subscript
\(G\setminus Y\), the original volume remains \(G\): the retained
legs come from \(\mathcal L(G)\), and the protected layer is always
\(S(G)\), with no new pins at removed sites.  The only response
assumption in this section is \(\mathsf P(A,\theta)\) of
Definition~\ref{h:P}.  Throughout, \(d\) is ambient \(\ell^1\)
distance.  Distances in the square geometry, including the distances
to corners in the shell definition, are ambient \(\ell^\infty\)
distances.  An exterior half-plane is understood through its lattice
sites; for example, the left exterior half-plane of \([-u,u]^2\) is
\(\{z\in\mathbb Z^2:z_1\le-u-1\}\).

\subsection{Shell geometry}

The parallel-side assertion below uses the same nonexceptional index
condition as the bound on \(A_i\).  This condition is necessary.  For
example, take \(l_0=64\), an integer \(s\ge390\), \(u=s^4\), \(a=0\),
\(y=(2s-u,-u-1)\), and \(i=2\).  The left side of the radius-\(k_2\)
square has first coordinate \(-u\), meets \(C_u(0)\), and has distance
one from the left exterior half-plane.  The corner \((-u,-u)\)
belongs to \(D_2\), so this index is exceptional.

\begin{lemma}\label{l:shellgeom}
For \(u\ge M_{\mathrm g}(l_0)=\max\{16,6(l_0+1)\}^4\),
\(|I_0|\le1\), and \(D_i\subset F_i\) for \(2\le i\le i_1\).
For \(2\le i\le i_1\) outside \(I_0\), \(|A_i|\le8l_0\).
For the same indices \(2\le i\le i_1\) outside \(I_0\), every side of
\(y+[-k_i,k_i]^2\) parallel to an edge of \(C_u(a)\)
and meeting \(C_u(a)\) is at distance at least \(s/3\) from the
exterior half-plane beyond that edge.
For \(i\ge2\), every point of \(C_u(a)\setminus F_i\)
within \(\ell^\infty\) distance one of \(F_i\) lies in \(D_{i-1}\);
\(y\) is not within that distance of \(F_i\).
\end{lemma}

\begin{proof}
Translate the coordinates so that \(a=0\).  We may interchange the
coordinates and reverse their signs for this geometric proof.  Since
\(\operatorname{dist}_\infty(y,C_u(0))=1\), we can consequently write
\[
 y=(\beta,-u-1),\qquad -u-1\le\beta\le u+1.
\]
These changes of coordinates concern sets and distances only.
They make no assertion about the frame specification.

We have \(u=s^4\), \(s\ge16\), and \(s\ge6(l_0+1)\).  In particular,
\[
 s^2+\tfrac43s<u,\qquad s^2+s<u/2,\qquad
 \lfloor s\rfloor>2s/3,\qquad s/3>l_0.
\]
For example, \(s^2+s\le(17/16)s^2<s^4/2\),
\(s^2+4s/3\le(13/12)s^2<s^4\), and
\(\lfloor s\rfloor\ge s-1>2s/3\).
For \(2\le i\le i_1\), we also have \(k_i\le s^2+s\).

Both upper corners of \(C_u(0)\) have distance \(2u+1\) from
\(y\).  At least one lower corner has distance at least \(u\)
from \(y\), because
\( |\beta+u|+|\beta-u|\ge2u\).
Thus at most one corner can have distance strictly less than \(u\)
from \(y\).  If a corner \(c\) has distance less than \(s/3\)
from a point \(z\in D_i\), then
\[
 |c-y|_\infty\le |c-z|_\infty+k_i
       <s^2+\tfrac43s<u.
\]
Only this one possible corner can therefore cause an exceptional
index.  For distinct indices the radii differ by at least
\(\lfloor s\rfloor>2s/3\).  Two points lying within \(s/3\)
of the same corner have distance strictly less than \(2s/3\),
whereas points of the two different spheres have distance at least
the difference of their radii.  This proves \(|I_0|\le1\).

Fix now \(2\le i\le i_1\), \(i\notin I_0\), and put \(k=k_i\).
The bottom horizontal side of the radius-\(k\) square misses
\(C_u(0)\).  Its top side has second coordinate \(-u-1+k\).
If this side meets \(C_u(0)\), its distances from the bottom and
top exterior half-planes are respectively \(k\) and \(2u+2-k\).
They are at least \(s\) and greater than \(u\), respectively,
since \(k\ge2s-1\ge s\) and \(k<u/2\).

Suppose the left vertical side meets \(C_u(0)\).  Its first
coordinate is \(\beta-k\in[-u,u]\), and its lowest point in
\(C_u(0)\) is \((\beta-k,-u)\in D_i\).  Its distance to the
lower left corner is \(\beta-k+u\).  Nonexceptionality gives
\(\beta-k+u\ge s/3\), so the whole side has distance
\(\beta-k+u+1\ge s/3+1\) from the left exterior half-plane.
Its distance from the right exterior half-plane is
\(u+1-\beta+k\ge k\ge s\).
For the right vertical side the lowest point is
\((\beta+k,-u)\); its distance to the lower right corner is
\(u-\beta-k\).  Thus its distance from the right exterior
half-plane is at least \(s/3+1\), and its distance from the
left exterior half-plane is \(\beta+k+u+1\ge k\ge s\).
This proves the parallel-side assertion, including when
\(\beta=\pm(u+1)\).

A point of \(A_i\) has distance at most \(l_0\) from at least
one of the four exterior half-planes.  A shell side parallel to
that half-plane cannot contain this point, by the clearance just
proved.  On a perpendicular side at most \(l_0\) lattice sites
have that distance at most \(l_0\).  There are four shell sides
and two perpendicular half-planes for each side.  Counting with
possible repetitions therefore gives \(|A_i|\le8l_0\).

For every \(2\le i\le i_1\),
\(k_i\ge is-1>(i-1)s\), proving \(D_i\subset F_i\).
Finally let \(i\ge2\), \(x\in C_u(0)\setminus F_i\), and
\(z\in F_i\) with \(|x-z|_\infty\le1\).  With \(r=(i-1)s\),
integrality and the Lipschitz property of \(|\,\cdot-y|_\infty\)
give
\[
 \lfloor r\rfloor\le |z-y|_\infty-1
       \le |x-y|_\infty\le\lfloor r\rfloor.
\]
Consequently \(x\in D_{i-1}\).  Each point of \(F_i\) has
distance greater than \((i-1)s\ge s>1\) from \(y\).
If \(F_i\) is empty both final assertions are vacuous, with
distance to it equal to \(+\infty\).
\end{proof}

\subsection{Factorization over a cut}

\begin{lemma}\label{l:cut}
Let \(F\subset\mathbb Z^2\) be finite, \(B\subset F\) have no
neighbour outside \(F\), and \(G=F\setminus B\).
Suppose \(X\subset G\) and \(G\setminus X=G_1\sqcup G_2\), with
no bond between the pieces, no site of \(G_1\) adjacent to \(F^c\),
and no site of \(G_2\) adjacent to \(B\).
For the original legs of \(G\), use pins \(\eta\) into \(F^c\)
and window letters \(s\) into \(B\). Put
\[
 e(\eta,s)=\log\Psi_{G\setminus X}(\eta s)-\log\Psi_G(\eta s).
\]
If \(|e(\eta,s)|\le E<\infty\) for all admissible \(\eta,s\),
then there is a strictly positive probability \(\nu\) on window
configurations, independent of \(\eta\), such that
\[
 e^{-2E}\nu(s)\le\mu^{+,\eta}_{\kappa,F}(n_B=s)
                         \le e^{2E}\nu(s).
\]
Consequently any two window laws have total variation at most
\(1-e^{-2E}\).
One may take \(\nu(s)\) proportional to \(W_B(s)Z_{G_1}^{s}\),
where \(W_B\) is the actual product of frame site weights and
internal window bond weights. Only original legs are retained in
the partition functions after removing \(X\).
\end{lemma}

\begin{proof}
Fix a proper colouring \(\kappa\).  The window configuration space
is the nonempty finite product \(\Omega_B=\prod_{x\in B}D_x\).
For \(s\in\Omega_B\) define
\[
 W_B(s)=\prod_{x\in B}\frac1{12}
       \prod_{\substack{\{x,z\}\subset B\\d(x,z)=1}}
                       (1+s_x\cdot s_z).
\]
All these weights are positive by Lemma~\ref{l:threeframes}.
The absence of bonds from \(B\) to \(F^c\) implies, on summing
the actual frame weights over \(G\), that
\begin{equation}\label{e:windowfactor}
 \mu^{+,\eta}_{\kappa,F}(n_B=s)
       =\frac{W_B(s)Z_G^{\eta s}}
                    {Z^{+,\eta}_{\kappa,F}}.
\end{equation}
Here the sum over \(G\) equals \(Z_G^{\eta s}\) by
Lemma~\ref{l:exact}: every summed spin has at most four incident
bond or original-leg factors.  This use of quadrature concerns
only a partition polynomial, with the letters \(s\) held fixed.
All partition sums in this proof are finite and strictly positive
by Lemma~\ref{l:finitefacts}; the actual frame normalizer is also
strictly positive by Definition~\ref{d:framemodel}.

For \(r=1,2\), define the partition sums on \(G_r\) by retaining
only the original legs of \(G\) whose tails lie in \(G_r\).
In particular, do not introduce factors on legs into \(X\).
The hypotheses on adjacency say that the retained pins on \(G_1\)
are all window letters and those on \(G_2\) are all exterior
letters.  The absence of bonds between the pieces and the product
site functional give
\begin{equation}\label{e:cutfactorization}
 Z_{G\setminus X}^{\eta s}=Z_{G_1}^{s}Z_{G_2}^{\eta},
 \qquad
 Z_{G\setminus X}^{0}=Z_{G_1}^{0}Z_{G_2}^{0}.
\end{equation}
The empty-piece convention makes these identities valid when
one or both pieces are empty.  Using the definition of \(e\),
we obtain
\[
 Z_G^{\eta s}
   =e^{-e(\eta,s)}
       \frac{Z_G^0 Z_{G_2}^{\eta}}{Z_{G_1}^0Z_{G_2}^0}
       Z_{G_1}^{s}.
\]
Set
\[
 Q=\sum_{s\in\Omega_B}W_B(s)Z_{G_1}^{s},\qquad
 \nu(s)=Q^{-1}W_B(s)Z_{G_1}^{s}.
\]
Every term in \(Q\) is positive and \(\Omega_B\) is nonempty,
so \(0<Q<\infty\), and \(\nu\) is strictly positive and does
not depend on \(\eta\).  Equations~\eqref{e:windowfactor}
and~\eqref{e:cutfactorization} imply
\[
 \mu^{+,\eta}_{\kappa,F}(n_B=s)
       =c(\eta)\nu(s)e^{-e(\eta,s)},\qquad
 c(\eta)^{-1}=\sum_{s\in\Omega_B}\nu(s)e^{-e(\eta,s)}.
\]
Admissible boundary and window assignments exist, so the assumed
absolute-value bound also implies \(E\ge0\).  The last sum lies
in \([e^{-E},e^E]\); its reciprocal lies in the same interval.
This proves the two pointwise comparisons.

For any two finite probability laws \(p,p'\),
\[
 d_{\mathrm{TV}}(p,p')
    =1-\sum_s\min\{p(s),p'(s)\}.
\]
Both window laws are at least \(e^{-2E}\nu\), so their overlap
sum is at least \(e^{-2E}\).  The asserted total variation bound
follows.  When \(B=\varnothing\), \(\Omega_B\) has its single
empty configuration, \(\nu\) assigns it mass one, and the same
proof gives total variation zero.
\end{proof}

\subsection{The cost of removing a connected set}

\begin{lemma}\label{l:removal}
Assume \(\mathsf P(A,\theta)\), and put
\(\mathsf F(d)=A\theta^{\max\{d,0\}}\).
Let \(x_1,\ldots,x_p\) be distinct points of \(G\), \(p\ge1\),
with \(b_G(x_j)=d(x_j,G^c)\ge3\), such that each \(x_j\), \(j\ge2\),
has a nearest neighbour among its predecessors.
For \(X=\{x_1,\ldots,x_p\}\) and every original unit-pin assignment,
\[
 |\log\Psi_{G\setminus X}(\tau)-\log\Psi_G(\tau)|
 \le1+4\mathsf F(b_G(x_1)-3)
                  +2\sum_{j=2}^p\mathsf F(b_G(x_j)-3).
\]
\end{lemma}

\begin{proof}
Here \(G\) is finite, as in Definition~\ref{d:pinmodel}.
Put \(Y_0=\varnothing\), \(Y_j=\{x_1,\ldots,x_j\}\), and
\(D_j=G\setminus Y_j\).  The depth assumption implies
\(Y_j\subset G\setminus S(G)\).  At the \(j\)-th removal let
\(u_1,\ldots,u_k\) be the surviving neighbours of \(x_j\).
The site \(x_j\) has no original leg, all four of its original
neighbours belong to \(G\), and \(k=4\) for \(j=1\).
For \(j\ge2\) one predecessor is a neighbour, so \(0\le k\le3\).

Integrate the spin at \(x_j\) in the finite functional.  By
Lemma~\ref{l:moments}, the resulting polynomial is
\begin{equation}\label{e:removalpolynomial}
 \begin{split}
 I_j&=\begin{cases}
 1,&k=0,1,\\
 1+\frac13\displaystyle\sum_{1\le r<t\le k}n_{u_r}\cdot n_{u_t},
       &k=2,3,\\
 1+\frac13\displaystyle\sum_{1\le r<t\le4}n_{u_r}\cdot n_{u_t}
                    +\frac1{15}Q_j,&k=4,
 \end{cases}\\
 Q_j&=(n_{u_1}\cdot n_{u_2})(n_{u_3}\cdot n_{u_4})
      +(n_{u_1}\cdot n_{u_3})(n_{u_2}\cdot n_{u_4})
      +(n_{u_1}\cdot n_{u_4})(n_{u_2}\cdot n_{u_3}).
 \end{split}
\end{equation}
The displayed definition of \(Q_j\) is made only in the case
\(k=4\); no \(u_4\) or quartic polynomial is formed in the other
cases.  For either
\(\bullet=\tau\) or \(\bullet=0\), finite summation gives
\(Z_{D_{j-1}}^\bullet=Z_{D_j}^\bullet
\mathbb E_{D_j}^\bullet I_j\).  Positivity of all normalizers
follows from Lemma~\ref{l:finitefacts}.  Consequently
\begin{equation}\label{e:removaltelescope}
 \log\Psi_G(\tau)-\log\Psi_{G\setminus X}(\tau)
   =\sum_{j=1}^p
      \left(\log\mathbb E_{D_j}^\tau I_j
                         -\log\mathbb E_{D_j}^0 I_j\right).
\end{equation}

We first give a uniform lower bound on each expectation in
\eqref{e:removaltelescope}.  Distinct neighbours of a square-lattice
site have distance two from one another, so no two \(u_r\)'s are
adjacent.  Each satisfies \(b_G(u_r)\ge b_G(x_j)-1\ge2\),
has no original leg, and has at most three present bonds after
\(x_j\) is removed.  Condition on every spin except these
\(k\) spins, on a conditioning event of positive probability.
Because there is no bond between these sites and none has an original
leg, the joint conditional weight is a constant times the product,
over \(r\), of its site weight and its surviving bond factors.
Summing over the \(k\) spins gives the product of their one-site
normalizers, each positive by Lemma~\ref{l:finitefacts}.
The joint conditional law is therefore the product of these
one-site laws.  Lemma~\ref{l:35} applies to each of them, with at
most three factors, so their means \(m_r\) satisfy \(|m_r|\le3/5\).
For \(k=0\) the product is the empty product.  Thus
\[
 \sum_{r<t}m_r\cdot m_t
    =\frac12\left(\left|\sum_r m_r\right|^2
                                  -\sum_r|m_r|^2\right)
    \ge-\frac{k}{2}\frac9{25}.
\]
For \(k\le3\) this proves
\(\mathbb E_{D_j}^\bullet I_j\ge1-3k/50\ge41/50>1/2\).
For \(k=4\), the pointwise bound \(|Q_j|\le3\) gives
\[
 \mathbb E_{D_j}^\bullet I_j
       \ge1-\frac6{25}-\frac3{15}
       =\frac{14}{25}>\frac12.
\]
For \(k=0\) the polynomial is one and its expectation is one.
These arguments average over positive conditioning fibres only;
zero-probability fibres are immaterial.

For each pair \(r<t\), the sites \(u_r,u_t\) survive in
\(D_j\), have distance two, and have the common removed neighbour
\(x_j\in Y_j\).  If \(z\in S(G)\), the distance function
\(b_G\) is one-Lipschitz and \(b_G(z)\le2\).  Therefore
\[
 d(u_r,z)\ge b_G(u_r)-2\ge b_G(x_j)-3,
 \qquad d(u_r,S(G))\ge b_G(x_j)-3.
\]
Property \(\mathsf P(A,\theta)\) now gives
\[
 \left|\mathbb E_{D_j}^\tau(n_{u_r}\cdot n_{u_t})
             -\mathbb E_{D_j}^0(n_{u_r}\cdot n_{u_t})\right|
       \le \mathsf F(b_G(x_j)-3).
\]
The set \(S(G)\) here is the original one throughout the
telescoping.  In particular, no distance to a newly created hole
has entered this estimate.

For positive real \(a,b\ge1/2\), integration of \(1/t\)
between them gives \(|\log a-\log b|\le2|a-b|\).
When \(j\ge2\), there are at most three pairs in
\eqref{e:removalpolynomial}, and hence the absolute value of
the \(j\)-th logarithmic difference is at most
\(2\mathsf F(b_G(x_j)-3)\).  At \(j=1\) there are six pairs,
and \(|\mathbb E_{D_1}^\tau Q_1-\mathbb E_{D_1}^0Q_1|\le6\).
The corresponding logarithmic difference is at most
\[
 2\left(2\mathsf F(b_G(x_1)-3)+\frac6{15}\right)
       =4\mathsf F(b_G(x_1)-3)+\frac45
       \le4\mathsf F(b_G(x_1)-3)+1.
\]
All inserted pair and four-spin products used here have separate
degree at most four: their endpoints have lost the bond to \(x_j\)
and each endpoint occurs once in an insertion.  No higher moment
identity is used.  Summing the bounds in
\eqref{e:removaltelescope} proves the claim, including \(p=1\).
\end{proof}

\subsection{A tapered set of physical sites}\label{s:cutgeom}

\begin{definition}\label{d:thickcut}
Fix \(l_0=64\), a nonexceptional shell with \(2\le i\le s\),
and write \(k=k_i\), \(F=F_i\), \(B=B_i\), \(G=F\setminus B\).
For \(z\in\mathbb Z^2\), set
\[
 t(z)=\big||z-y|_\infty-k\big|,\qquad
 h(z)=\operatorname{dist}_\infty(z,C_u(a)^c)\quad
        (h(z)=0\text{ outside }C_u(a)).
\]
Define
\[
 \begin{gathered}
 U=\{z\in C_u(a):t(z)\le\min(h(z)-32,s/12)\},\\
 X_0=\{z\notin U:\operatorname{dist}_\infty(z,U)=1\},\qquad
 X=\{z:d(z,X_0)\le1\}.
 \end{gathered}
\]
Set \(G_1=(U\setminus B)\setminus X\) and
\(G_2=G\setminus(U\cup X)\).
\end{definition}

\begin{lemma}\label{l:cutgeom}
There exists \(s_0\ge\max\{16,6(64+1)\}\) such that, uniformly
in the translate, \(y\), and nonexceptional integer \(2\le i\le s\),
the following hold for \(s=u^{1/4}\ge s_0\).
The set \(U\) is finite, nonempty, and star-connected, and \(U^c\)
is star-connected. We have \(B\subset U\subset F\), all neighbours
of \(B\) lie in \(U\), and no site of \(U\) is adjacent to \(F^c\).
The set \(X_0\) is nonempty and lies in \(G\). For \(x\in X_0\),
\[
 b_G(x)\ge
 \begin{cases}\max\{t(x),16\}-2,&t(x)\le s/12,\\s/12-1,&t(x)>s/12,
 \end{cases}
 \qquad b_G(x)\ge14 .
\]
For each integer \(j\ge0\), there are at most \(48\) sites of
\(X_0\) with \(t(x)=j\le s/12\).
For \(0<\vartheta<1\), put
\[
 C_{\rm near}(\vartheta)=48\vartheta^{14}
                \left(17+\frac1{1-\vartheta}\right).
\]
Then
\[
 \sum_{x\in X_0}\vartheta^{b_G(x)}
 \le C_{\rm near}(\vartheta)
                 +4(2u+1)^2\vartheta^{s/12-1}.
\]
The thickened set \(X\) lies in \(G\), has \(b_G\ge13\), and
\(G\setminus X=G_1\sqcup G_2\) satisfies all separation and
no-adjacency assumptions of Lemma~\ref{l:cut}. Moreover
\[
 \sum_{z\in X}\vartheta^{b_G(z)}
       \le5\vartheta^{-1}\sum_{x\in X_0}\vartheta^{b_G(x)} .
\]
\end{lemma}

\begin{proof}
We may take \(s_0=390\).  Translate so that \(a=0\), and put
\(C=C_u(0)\), \(T=s/12\), and \(\rho(z)=|z-y|_\infty\).
Here \(s,T\) are real, whereas \(k,\rho(z),t(z),h(z)\) are
integers at lattice sites.
The functions \(h(z)=\max\{0,u+1-|z|_\infty\}\), \(\rho\),
and \(t=|\rho-k|\) are one-Lipschitz for \(\ell^\infty\)
distance.  We shall repeatedly use the following consequences of
\(s\ge390\), \(u=s^4\), and \(2\le i\le s\):
\begin{equation}\label{e:sitecutscales}
 \begin{gathered}
 2s-1\le k\le s^2\le u-31,\qquad T\ge1,\qquad
 k-T-1>(i-1)s,\\
 s/3-2T=s/6\ge65>32,\qquad
 2u+2-2T>32,\qquad u+1-T>32.
 \end{gathered}
\end{equation}
Indeed, the third inequality follows from
\(k\ge is-1\) and \(s-T-2=11s/12-2>0\).
The inequalities involving \(u\) follow from \(u=s^4\)
and \(s\ge390\); for instance \(s^4-s^2>31\) and
\(s^4-s/12>32\).  All constants in what follows are thus uniform.

\emph{Connectivity of the tapered set.}
On all of \(\mathbb Z^2\) its definition can be written as
\begin{equation}\label{e:taperconditions}
 U=\{z:t(z)\le T,\ h(z)-t(z)\ge32\}.
\end{equation}
The second condition forces \(z\in C\).
If \(\rho(z)>k\), move each coordinate whose absolute
displacement from \(y\) is maximal one step toward the corresponding
coordinate of \(y\).  This is a star-step, decreases \(\rho\)
and \(t\) by one, and decreases \(h\) by at most one.
If \(\rho(z)<k\), move one coordinate with maximal absolute
displacement one step away from \(y\); this is a nearest-neighbour
step, increases \(\rho\) by one, and again decreases \(t\) by
one and \(h\) by at most one.  The case of zero displacement can
use either sign; in fact \(y\notin U\).  Starting in \(U\),
these steps preserve both inequalities in
\eqref{e:taperconditions}.  After exactly \(t(z)\) steps they
reach
\[
 \mathcal K=\{z:\rho(z)=k,\ h(z)\ge32\}\subset U.
\]

We verify that \(\mathcal K\) is nonempty and nearest-neighbour
connected.  Set \(m=u-31\).  Since \(|y_j|\le u+1=m+32\)
and \(32\le k\le m\), each integer interval
\[
 J_j=[y_j-k,y_j+k]\cap[-m,m]\cap\mathbb Z,
       \qquad j=1,2,
\]
is nonempty.  It contains at least one of \(y_j-k,y_j+k\):
otherwise \([-m,m]\cap\mathbb Z\) would lie strictly between
these two integers, which is impossible since \(m\ge k\).
Consequently the rectangle \(J_1\times J_2\) contains a corner
of \(y+[-k,k]^2\).  Reverse the signs of the coordinates for
this geometric argument so that the contained corner is
\(y-(k,k)\).  The intersections of this rectangle with the left
and bottom sides of the square are nonempty lattice segments
joined at that corner.  If the right side meets the rectangle,
its bottom endpoint \(y+(k,-k)\) lies on the bottom segment.
If the top side meets the rectangle, its left endpoint
\(y+(-k,k)\) lies on the left segment.  Thus all side segments
are joined by nearest-neighbour paths.  Their union is precisely
\(\mathcal K\), since \(h\ge32\) is equivalent to
\(|z|_\infty\le m\).  The radial paths above now show that
\(U\) is nonempty and star-connected.  It is finite because
it is contained in \(C\).

The exterior \(C^c\) is nearest-neighbour connected.  To see this
directly, a site with first coordinate outside \([-u,u]\) can
move vertically to height \(u+1\), remaining outside \(C\).
A site whose second coordinate is outside \([-u,u]\) can first
move horizontally to first coordinate \(u+1\), then vertically
to height \(u+1\).  Along the row at height \(u+1\) these paths
join.  We join each \(z\in C\setminus U\) to \(C^c\) while
avoiding \(U\).  There are three cases.
If \(\rho(z)>k+T\), repeatedly move a coordinate of maximal
displacement away from \(y\).  Each step increases \(\rho\),
so stays outside \(U\), and the path eventually leaves \(C\).
If \(\rho(z)<k-T\), repeatedly move a coordinate of maximal
nonzero displacement toward \(y\).  The value of \(\rho\)
does not increase, and the sum of the absolute displacements
decreases by one at every step.  The path stays outside \(U\)
and reaches \(y\notin C\), unless it leaves \(C\) earlier.
In the remaining case \(t(z)\le T\) and \(h(z)-t(z)<32\).
Here \(z\ne0\), since \(h(0)-t(0)\ge u+1-T>32\) whenever
\(t(0)\le T\).  Move a coordinate of maximal absolute value
away from zero and continue in that direction until leaving
\(C\).  At every step in \(C\), \(h\) decreases by one,
whereas \(t\) decreases by at most one.  Thus \(h-t<32\)
persists, independently of whether \(t\le T\) persists.
This path also avoids \(U\).  These cases prove that \(U^c\)
is nearest-neighbour connected, hence star-connected.

\emph{The window and the raw boundary.}
For \(z\in B\), \(t(z)=0\) and \(h(z)\ge65\), so \(z\in U\).
A nearest neighbour of such a point has \(t\le1\) and
\(h\ge64\); it too belongs to \(U\).  For \(z\in U\) we
have \(h(z)\ge32\) and \(\rho(z)\ge k-T\).  Hence \(z\in F\),
and every nearest neighbour \(w\) of \(z\) satisfies
\(h(w)\ge31>0\) and \(\rho(w)\ge k-T-1>(i-1)s\).
There is therefore no bond from \(U\) to \(F^c\).

If \(x\in X_0\), choose \(v\in U\) with
\(|x-v|_\infty=1\).  Then
\(h(x)\ge31>0\) and \(\rho(x)\ge k-T-1>(i-1)s\),
so \(x\in F\).  Since \(x\notin U\) and \(B\subset U\),
we have \(x\in G\).  The raw boundary is finite, and is
nonempty: a point of \(U\) with largest first coordinate has
its right nearest neighbour in \(X_0\).

\emph{Depths of the raw boundary.}
Fix \(x\in X_0\), put \(t=t(x)\), \(h=h(x)\), and keep the
point \(v\) chosen above.  Distance from \(x\) to \(B\) in
the \(\ell^\infty\) metric is at least \(t\), and also at
least \(65-h\), by the Lipschitz properties of \(\rho\) and
\(h\).  Its distance to \(C^c\) is \(h\).  For
\(w\in C\setminus F\),
\[
 |x-w|_\infty\ge\rho(x)-\rho(w)
       \ge k-t-(i-1)s\ge s-1-t.
\]
If this last set is empty the same lower bound is valid with
distance \(+\infty\).  As \(G^c=B\cup C^c\cup(C\setminus F)\)
and \(d\ge\operatorname{dist}_\infty\), we conclude that
\begin{equation}\label{e:sitecutdepth}
 b_G(x)\ge
    \min\{\max(t,65-h),\ h,\ s-1-t\}.
\end{equation}

Suppose first \(t\le T\).  Since \(x\notin U\) but \(v\in U\),
\[
 h<32+t,\qquad
 h\ge h(v)-1\ge32+t(v)-1\ge30+t.
\]
Thus
\begin{equation}\label{e:rawdepthrange}
                30+t\le h<32+t.
\end{equation}
In \eqref{e:sitecutdepth}, the first entry is at least
\(\max(t,33-t)\), which is at least \(\max(t,16)\).
The second entry is at least \(30+t\ge\max(t,16)-2\).
The third is at least \(11s/12-1\), which is at least
\(\max(t,16)-2\): use \(t\le T\), \(s\ge390\), and
\(s-1-2T>0\).  This proves the stated near depth.

If \(t>T\), then \(t\le T+1\), and
\(h\ge32+t(v)-1\ge30+t>T+30\).
The three entries in \eqref{e:sitecutdepth} are consequently
at least \(T\), \(T+30\), and \(11s/12-2\), respectively.
Each is at least \(T-1\), proving the far depth.
The near depth is at least \(14\), and the far depth is
at least \(s/12-1\ge390/12-1>14\).

\emph{Counting near sites, including possible coincident sides.}
Fix an integer \(0\le j\le T\).  A raw boundary site with
\(t(x)=j\) lies on a square boundary of radius \(k-j\) or
\(k+j\) about \(y\), and on a boundary \(h=h'\), where
\(h'\) is an integer in \([30+j,32+j)\).
There are at most three choices for \(h'\).  These levels
are the boundaries of squares about zero of radius \(u+1-h'\),
which is positive by \eqref{e:sitecutscales}.

We exclude coincident supporting lines of these two square
boundaries.  For each edge of \(C\), define the signed inward
depth of a parallel line.  For a vertical line with first
coordinate \(q\), the left and right depths are
\(q+u+1\) and \(u+1-q\), respectively; for horizontal lines
use the same expressions in the second coordinate.  Each side
of a radius-\(k\pm j\) square corresponds to the side with
the same orientation and sign of the radius-\(k\) square.
Their signed depths from any fixed parallel exterior half-plane
differ by at most \(j\).

The varying-coordinate intervals of both corresponding sides
meet \([-u,u]\): their centres have coordinate in
\([-u-1,u+1]\), and their radii are at least
\(k-j\ge2s-1-T\ge1\).
If the original side's fixed coordinate lies in \([-u,u]\),
the original side therefore meets \(C\).  We have
\(u\ge390^4=M_{\mathrm g}(64)\), \(2\le i\le s<i_1\), and
\(i\notin I_0\).  Thus the nonexceptional parallel-side assertion
of Lemma~\ref{l:shellgeom} applies and gives original signed depth
at least \(s/3\).  The shifted depth is at least
\[
                   s/3-j>32+j>h'.
\]
The strict inequality uses \(s/3-2j\ge s/6>32\).
If the original line lies beyond the specified edge, its signed
depth is at most zero; the shifted depth is at most
\(j<30+j\le h'\).  If it lies beyond the opposite edge,
its original depth is at least \(2u+2\); the shifted depth is
at least \(2u+2-j>32+j>h'\), by
\eqref{e:sitecutscales}.  These three cases exhaust all positions
of the original line.  None can yield shifted depth \(h'\),
which would be necessary for coincidence with a side of the
\(h'\)-square.

Without coincident supporting lines, a side of one axis-parallel
square meets the boundary of another in at most two points:
its parallel sides do not intersect it, and there are only two
perpendicular sides.  Their square boundaries thus meet in
at most eight lattice sites.  There are at most two radius
choices and at most three depth levels.  We have proved the
required bound \(2\cdot3\cdot8=48\).

For \(0<\vartheta<1\), the near sites therefore contribute at most
\[
 \begin{split}
 48\sum_{j=0}^{\infty}\vartheta^{\max(j,16)-2}
 &=48\vartheta^{14}
                    \left(17+\frac{\vartheta}{1-\vartheta}\right)\\
 &\le48\vartheta^{14}
                    \left(17+\frac1{1-\vartheta}\right).
 \end{split}
\]
The far sites contribute at most
\(|X_0|\vartheta^{T-1}\le(2u+1)^2\vartheta^{T-1}\),
which is at most the displayed far term in the statement.
This proves the raw cut sum.

\emph{Thickening and separation.}
The ambient distance function \(b_G(z)=d(z,G^c)\), defined
also outside \(G\) by zero, is one-Lipschitz for \(d\).
Every \(z\in X\) has a point \(x\in X_0\) with \(d(z,x)\le1\).
It follows that \(b_G(z)\ge b_G(x)-1\ge13\), so \(X\subset G\).
In particular \(X\) meets neither \(B\) nor the protected
original layer \(S(G)\).  Directly summing over the five sites
in each nearest-neighbour cross
gives
\[
 \sum_{z\in X}\vartheta^{b_G(z)}
    \le\sum_{x\in X_0}\sum_{d(z,x)\le1}
                              \vartheta^{b_G(z)}
    \le5\vartheta^{-1}\sum_{x\in X_0}\vartheta^{b_G(x)}.
\]

Since \(B\subset U\subset F\) and \(X\subset G\), the two
sets specified in Definition~\ref{d:thickcut} give the disjoint
union \(G\setminus X=G_1\sqcup G_2\).  If a bond joined
\(p\in G_1\) to \(q\in G_2\), then \(p\in U\), \(q\notin U\),
and \(|p-q|_\infty=1\); hence \(q\in X_0\subset X\), a
contradiction.  A point of \(G_1\subset U\) is not adjacent to
\(F^c\), as already proved.  A point adjacent to \(B\) lies
in \(U\), so cannot belong to \(G_2\).  Finally \(B\) itself
has no neighbour in \(F^c\), because \(B\subset U\) and
\(U\) has this property.  These are all the separation and
no-adjacency assumptions of Lemma~\ref{l:cut}.
\end{proof}

\begin{lemma}\label{l:Xconn}
Under Lemma~\ref{l:cutgeom}, the finite nonempty cut \(X\) is
nearest-neighbour connected. Hence it can be ordered so that every
site after the first has a previously listed nearest neighbour.
\end{lemma}

\begin{proof}
Lemma~\ref{l:cutgeom} says that \(U\) is finite, nonempty and
star-connected, and that \(U^c\) is star-connected.
Lemma~\ref{l:boundaryconn} in Appendix~\ref{s:topology} therefore makes
\(X_0=\partial_\infty U\) nonempty and star-connected.
Its finiteness follows from that of \(U\).
Lemma~\ref{l:thicken} in that appendix makes \(X\) nearest-neighbour connected,
and its definition makes it finite and nonempty.

For completeness, choose a root \(x_1\in X\) and order all
points by nondecreasing shortest-path distance from the root
in the nearest-neighbour graph induced by \(X\), breaking
ties arbitrarily.  These distances are finite nonnegative
integers.  Every nonroot point has a neighbour at distance
one smaller, obtained from a shortest path; that neighbour
has already been listed.  This also covers a singleton cut.
\end{proof}

\subsection{The uniform window estimate}

\begin{theorem}\label{t:cutMOS}
If \(\mathsf P(A,\theta)\) holds, then the family of plane frame
specifications has \(\mathsf W(64,p_0,M_\star)\), with common
\(p_0\in(0,1]\) and integer \(M_\star\ge1\) depending only on
\(A,\theta\). The constants work for every proper plane
three-colouring and every translate.
Specifically, let
\[
 E=1+4A+10A\theta^{-4}
       \left[48\theta^{14}\left(17+\frac1{1-\theta}\right)+1\right],
 \qquad p_0=e^{-2E}.
\]
There exists an integer \(M_\star\ge1\) for which (W) holds
with this displayed choice of \(p_0\).
\end{theorem}

\begin{proof}
Fix \(A\ge2\) and \(0<\theta<1\) satisfying
Definition~\ref{h:P}.  We first choose a common scale threshold.
Let \(c=-\log\theta>0\).  For \(u=s^4\ge1\),
\[
 4(2u+1)^2\theta^{s/12-1}
       \le36\theta^{-1}s^8 e^{-cs/12}
       \le K_\theta s^{-2},\qquad
 K_\theta=36\theta^{-1}\,10!\left(\frac{12}{c}\right)^{10}.
\]
The second inequality follows from
\(e^r\ge r^{10}/10!\) for \(r=cs/12>0\), by the
nonnegative power series for the exponential.
Choose, for example,
\[
 M_\star=\left\lceil\max\{390^4,K_\theta^2,1\}\right\rceil.
\]
Then every integer \(u\ge M_\star\) satisfies the threshold
of Lemma~\ref{l:cutgeom} and
\begin{equation}\label{e:windowtail}
                 4(2u+1)^2\theta^{s/12-1}\le1.
\end{equation}

Fix any proper plane colouring, any translate \(C_u(a)\),
any indicated exterior site \(y\), and any nonexceptional
integer \(2\le i\le s\).  Form \(F,B,G,U,X_0,X,G_1,G_2\)
as in Definition~\ref{d:thickcut}.  Lemmas~\ref{l:cutgeom}
and~\ref{l:Xconn} give a nonempty finite cut in \(G\), an
ordering \(x_1,\ldots,x_p\) as in Lemma~\ref{l:removal},
and \(b_G\ge13\) on it.  For every admissible exterior
assignment \(\eta\) and window assignment \(s_B\), use the
unit pins \(\eta s_B\) on the original legs of \(G\).
The removal lemma applies to these pins without any restriction
on their directions.  As \(b_G-3\ge0\), it yields
\begin{equation}\label{e:windowcutcost}
 \begin{split}
 |e(\eta,s_B)|
 &\le1+4A\theta^{b_G(x_1)-3}
                    +2A\theta^{-3}\sum_{j=2}^p\theta^{b_G(x_j)}\\
 &\le1+4A+2A\theta^{-3}\sum_{z\in X}\theta^{b_G(z)}\\
 &\le1+4A+10A\theta^{-4}
       \left[48\theta^{14}\left(17+\frac1{1-\theta}\right)
                        +4(2u+1)^2\theta^{s/12-1}\right]
 \le E.
 \end{split}
\end{equation}
The first site is kept as a seed in this estimate; its additive
cost is included in the term one.

All assumptions of Lemma~\ref{l:cut} hold by
Lemma~\ref{l:cutgeom}.  That lemma supplies one strictly
positive reference law on the actual window alphabet and the
minorization \(e^{-2E}\nu\), uniformly in \(\eta\).
Thus any two admissible boundaries for \(F_i\) give window
marginals with total variation at most \(1-e^{-2E}\).
The statement is also automatic for an empty window, whose
marginal is a point mass.  This is precisely
\(\mathsf W(64,p_0,M_\star)\) in Definition~\ref{d:W}.
All terms defining \(E\) are finite and positive, so
\(p_0=e^{-2E}\in(0,1]\).

The geometric estimates did not use a colouring, and the
removal estimate was uniform over original unit pins.
Quadrature in Lemma~\ref{l:cut} transfers only the sums over
the sites outside the window.  The reference law retains
\(W_B\) with the chosen frames and may vary with the colouring
and placement.  No equality of window laws at different
colourings or translates is needed.  The constants chosen
above consequently apply to every colouring and translate.
\end{proof}

\section{Walk expansion without scalar trace masses}\label{s:walks}

We work on a finite subset of the plane with the laws of
Definition~\ref{d:moment}. We first prove the finite expansion and
uniform response estimates, and then give the local and open-trail
contracts used to bound it.

\begin{definition}\label{d:domain}
Fix finite \(G\), \(S=S(G)\), and original unit pins \(\tau\).
A domain assigns statuses \(\mathrm U,\mathrm D,\mathrm X\);
\(\mathrm X\) sites are removed. Each nonremoved site has a kept
subset of its incident bonds inside \(G\). A bond is present if both
ends are nonremoved and both keep it. A \(\mathrm D\) site has at
most two present bonds. Protected sites in \(S\) have status
\(\mathrm U\) and keep every incident bond inside \(G\).
Use \(\lambda_{\mathrm U},\lambda_{\mathrm D}\) at the corresponding
sites and define
\[
 Z_D^\tau[f]=\sum_n f(n)
       \prod_{\text{present }\{x,y\}}(1+n_x\cdot n_y)
       \prod_{\substack{(x,z)\in\mathcal L(G)\\x\in S}}
                   (1+\tau_{xz}\cdot n_x)
       \prod_{x:\,\mathrm{st}(x)\ne\mathrm X}\lambda_{\mathrm{st}(x)}(n_x).
\]
The sum is over the product of the supports \(D\) at \(\mathrm U\)
sites and \(D\cup\{0\}\) at \(\mathrm D\) sites.
Define the free sum, normalized expectations and \(\Psi_D\) as before.
The class \(\mathcal D_{\mathrm{cons}}(G)\) consists of domains in
which every \(\mathrm U\) site keeps all its bonds, every \(\mathrm D\)
site keeps exactly two, and every bond between nonremoved sites is
kept by both ends or by neither.
An auxiliary network has the same data but permits protected
\(\mathrm U\) sites to have arbitrary kept subsets. It retains every
original leg factor. Neither domains nor auxiliary networks carry
a scalar trace mass. Only the networks occurring below, with at most
four factors at each site, are used.
\end{definition}

All normalizers in this definition are positive by
Lemma~\ref{l:finitefacts}. Original legs occur only at protected
\(\mathrm U\) sites. In particular every departure site outside
\(S\) is unpinned.
The domains \(G\setminus Y\), \(Y\subset G\setminus S\), in
Definition~\ref{d:pinmodel} are members of
\(\mathcal D_{\mathrm{cons}}(G)\), with statuses
\(\mathrm U,\mathrm X\) only.
An auxiliary network allows a stopped arrival to delete a protected
site's arrival bond while retaining all of that site's original legs.

\begin{definition}\label{d:trace}
A trail uses distinct present unoriented bonds. Its prefixes are
\(w_j=(v_0,\ldots,v_j)\), and \(\mathcal T_D(u,T)\) denotes trails
starting at \(u\) and stopping at their first hit of \(T\).
No trail departs from \(S\).
At departure from \(y=v_{j-1}\), drop \(\{y,v_j\}\) at both ends.
If \(y\) is \(\mathrm U\), change it to \(\mathrm D\), keeping its
other kept bonds. If \(y\) is \(\mathrm D\), remove it.
Call the resulting network \(D[w_j]\), with \(D[w_0]=D\).
Admissible initial starts are unpinned sites outside \(S\) of status
\(\mathrm U\) with at most three present bonds, or status \(\mathrm D\)
with at most one. The same restrictions then hold at every later
departure. Set
\[
 a_D(w)=3^{-N_{\mathrm U}(w)}5^{-N_{\mathrm D}(w)},\qquad
 R_D^\bullet(w)=Z^\bullet_{D[w]}/Z^\bullet_D,
\]
where the counts record the status at departure, not the original status.
Original legs are always unchanged, including at a terminal hit of \(S\).
\end{definition}

We record why this tracing rule is defined on every trail under
consideration. At length zero the stated bounds hold by the start
hypothesis. After an arrival, a \(\mathrm U\) site has lost the
arrival bond from at most four present bonds, so it has at most three;
a \(\mathrm D\) site has lost it from at most two, so it has at
most one. At a \(\mathrm U\) departure, deleting the out-bond
leaves at most two present bonds at the new \(\mathrm D\) site.
At a \(\mathrm D\) departure there is at most the out-bond:
deleting it and removing the site removes no other present bond.
In particular no unused present bond is destroyed by removal.
Induction on the prefix length now shows that every unused bond
of the trail is still present when it is needed, and that every
departure satisfies the asserted bounds.
Sites in \(S\) are never departed from, so none is removed.
They retain all original leg factors even when the final arrival
drops a bond at such a site.

There are at most two departures from any site during this trace:
after its first \(\mathrm U\) departure it is \(\mathrm D\), and
after a \(\mathrm D\) departure it is removed. Also
\[
 |w|=N_{\mathrm U}(w)+N_{\mathrm D}(w),\qquad
 a_D(w)=3^{-|w|}(3/5)^{N_{\mathrm D}(w)}.
\]
This identity places the entire revisit factor in the coefficient.
Partition-function ratios telescope to
\(R_D^\bullet(w)=\prod_{j=1}^{|w|}\rho_j^\bullet\), where
\(\rho_j^\bullet\) is the ratio for the \(j\)-th step.

\begin{lemma}\label{l:stepid}
For vectors \(a,u_1,\ldots,u_k\in\mathbb R^3\),
\[
 \mathcal M\!\left((a\cdot n)\prod_{i=1}^k(1+n\cdot u_i)\right)
 =\frac13\sum_{i=1}^k(a\cdot u_i)
       \lambda_{\mathrm D}\!\left(\prod_{i'\ne i}(1+n\cdot u_{i'})\right)
       \quad(0\le k\le3),
\]
and
\[
 \lambda_{\mathrm D}\!\left((a\cdot n)\prod_{i=1}^k(1+n\cdot u_i)\right)
       =\frac15\sum_{i=1}^k a\cdot u_i\quad(0\le k\le1).
\]
An empty sum is zero. The second identity removes the departure site
with no residual scalar mass.
\end{lemma}

\begin{proof}
Expand the product into subsets of \(\{1,\ldots,k\}\).
In the first identity, the inserted linear factor permits only
subsets of sizes one and three to contribute, by
Lemma~\ref{l:moments}. The value is
\[
 \frac13\sum_i a\cdot u_i+
 \frac1{15}\sum_{i<j<l}
 \bigl((a\cdot u_i)(u_j\cdot u_l)
       +(a\cdot u_j)(u_i\cdot u_l)
       +(a\cdot u_l)(u_i\cdot u_j)\bigr).
\]
For the residual product, which has at most two factors,
\[
 \lambda_{\mathrm D}\left(\prod_{i'\ne i}(1+n\cdot u_{i'})\right)
 =1+\frac15\sum_{\substack{j<l\\j,l\ne i}}u_j\cdot u_l.
\]
Multiplying by \((a\cdot u_i)/3\) and summing over \(i\)
gives exactly the preceding expression, including the empty sums
when \(k=0,1,2\).

For the second identity, \(k=0\) gives zero by the first moment.
For \(k=1\), the constant part again vanishes and the second
moment is
\((3/5)(a\cdot u_1)/3=(a\cdot u_1)/5\).
There is no residual factor depending on \(n\).
Thus the site can be deleted with the displayed coefficient and
with no change of the normalization of the remaining site laws.
All polynomials evaluated here have degree at most four.
\end{proof}

\begin{proposition}\label{p:pairwalk}
Let \(D\) be a domain or auxiliary network, and let \(u\) be an
admissible initial start. Let \(v\ne u\) be nonremoved; if it is
\(\mathrm U\), assume at most three incident factors, counting
original legs; if it is \(\mathrm D\), assume at most one present bond.
Put \(T=S\cup\{v\}\). For \(\bullet=\tau,0\),
\[
 \mathbb E_D^\bullet(n_u\cdot n_v)=
 \sum_{w\in\mathcal T_D(u,T)}a_D(w)R_D^\bullet(w)
           \mathbb E_{D[w]}^\bullet(n_{v_{|w|}}\cdot n_v).
\]
If the endpoint is \(v\), its last factor is independent of the
choice of \(\bullet\): it is \(1\) for status \(\mathrm U\),
and \(3/5\) for status \(\mathrm D\).
Every partition product with an inserted observable used in this
identity has separate degree at most four.
\end{proposition}

\begin{proof}
We prove an identity of unnormalized numerators. Write
\(n_u\cdot n_v=\sum_{\beta=1}^3n_{u,\beta}n_{v,\beta}\), and fix
\(\beta\). At any nonterminal prefix ending at \(y\), the marked
factor is \(n_{y,\beta}n_{v,\beta}\), with \(y\ne v\) and
\(y\notin S\). Sum over \(n_y\) first.
If its status is \(\mathrm U\), the first identity of
Lemma~\ref{l:stepid}, with \(a=e_\beta\) and the present
neighbouring spins as the \(u_i\), gives
\[
 Z^\bullet_E[n_{y,\beta}n_{v,\beta}]
   =\frac13\sum_{o:\,\{y,o\}\text{ present in }E}
        Z^\bullet_{E[(y,o)]}[n_{o,\beta}n_{v,\beta}].
\]
Here \(E\) denotes the current network; \(E[(y,o)]\) is its
one-step trace, so its law at \(y\) is
\(\lambda_{\mathrm D}\) and the selected bond has been dropped.
If \(y\) is \(\mathrm D\), the second identity gives instead
\[
 Z^\bullet_E[n_{y,\beta}n_{v,\beta}]
   =\frac15\sum_{o:\,\{y,o\}\text{ present in }E}
        Z^\bullet_{E[(y,o)]}[n_{o,\beta}n_{v,\beta}],
\]
where \(y\) is removed in every term.
The departure bounds needed for these identities were verified
after Definition~\ref{d:trace}. There are no leg factors at \(y\).
If there is no present bond, either displayed sum is zero.

Iterate these equalities until the marked site reaches \(T\).
Every iteration consumes a previously unused present unoriented
bond. There are only finitely many bonds initially, so the
iteration terminates after at most that many steps, even for
branches which end with an empty sum. No unused present bond is
destroyed by the trace, as proved after
Definition~\ref{d:trace}. Hence the branches that reach \(T\) are
in bijection with the trails in \(\mathcal T_D(u,T)\): a sequence
of choices determines one such trail and conversely. The product of its
coefficients is \(a_D(w)\). We have proved
\[
 Z^\bullet_D[n_{u,\beta}n_{v,\beta}]
 =\sum_{w\in\mathcal T_D(u,T)}
       a_D(w)Z^\bullet_{D[w]}
                   [n_{v_{|w|},\beta}n_{v,\beta}].
\]
Sum over \(\beta\) and divide once by the positive \(Z^\bullet_D\).
This proves the expansion.

If the endpoint is \(v\), this is its first hit, so the status of
\(v\) has not changed. For status \(\mathrm U\), its spin has
squared norm one pointwise. For status \(\mathrm D\), its sole
possible present bond in \(D\) must have been the arrival bond;
there are no original legs at a \(\mathrm D\) site. After arrival
it is isolated and has its normalized law
\(\lambda_{\mathrm D}\), independently of all other spins.
Thus its expected squared norm is \(3/5\), for either pin choice.

For completeness, before a terminal hit the fixed target \(v\)
has at most three factors if \(\mathrm U\), or at most one
present bond if \(\mathrm D\). Inserting its one coordinate
therefore gives degree at most four or two, respectively.
The moving marked site has at most three bonds if \(\mathrm U\),
or one if \(\mathrm D\), and no legs; its inserted coordinate
has the same respective bounds. At a terminal site in
\(S\setminus\{v\}\), the lost arrival bond leaves at most three
factors, counting all unchanged original legs, and the inserted
coordinate gives degree at most four. At a terminal \(v\) of
status \(\mathrm U\), at most two factors remain after arrival;
inserting \(|n_v|^2\) again gives degree at most four.
A terminal \(v\) of status \(\mathrm D\) is isolated and has
only the degree-two insertion. All unmarked sites have at most
four factors, and changing a law or deleting a bond raises no
degree. These cases cover every numerator in the iteration.
\end{proof}

\begin{definition}\label{d:ratios}
At a step \(j\), put
\(\rho_j^\bullet=Z^\bullet_{D[w_j]}/Z^\bullet_{D[w_{j-1}]}\).
At a \(\mathrm D\) departure \(\rho_j^\tau=\rho_j^0=1\).
At a \(\mathrm U\) departure from \(y\) to \(o\), let \(E'\)
be the pre-step network with \(y\) removed.
If the other present neighbours are \(f,g\), set
\[
 I=1+\tfrac13(P_{fg}+P_{fo}+P_{go}),\quad
 J=1+\tfrac15P_{fg},\quad P_{ab}=n_a\cdot n_b,\qquad
 (\rho_j^\bullet)^{-1}=\mathbb E_{E'}^\bullet I/
                                      \mathbb E_{E'}^\bullet J.
\]
With only one other neighbour \(f\), use \(I=1+P_{fo}/3,\ J=1\);
with no other neighbour, \(\rho_j^\bullet=1\).
Straight/turning refers to the incoming direction, with a designated
absent incoming bond at the first step.
\end{definition}

These formulas follow by summing over the departure spin before
performing the remaining product sum. With three neighbours, the
pre-step \(\lambda_{\mathrm U}\) sum is \(I\), by its second
moment and vanishing odd moments. The post-step
\(\lambda_{\mathrm D}\) sum over its two retained bonds is \(J\),
by its second moment \(1/5\). Thus
\[
 Z^\bullet_{D[w_{j-1}]}=Z^\bullet_{E'}\mathbb E^\bullet_{E'}I,
 \qquad
 Z^\bullet_{D[w_j]}=Z^\bullet_{E'}\mathbb E^\bullet_{E'}J.
\]
With two neighbours these same sums are \(1+P_{fo}/3\) and \(1\);
with one they are both \(1\). At a \(\mathrm D\) departure
there is at most the selected bond, and
\(\lambda_{\mathrm D}(1+n\cdot n_o)=1\), so deletion leaves
the partition function unchanged. All divisions are justified by
Lemma~\ref{l:finitefacts}. These computations also show directly
that the revisit has ratio one for both pin choices.

Incoming directions are lattice unit vectors, namely
\((\pm1,0)\) or \((0,\pm1)\).
At an admissible start, all four neighbouring sites lie in \(G\)
because the start is outside \(S\); there are at most three
present bonds. Hence an absent bond can always be designated as
the initial incoming bond. At every later step the actual incoming
bond has already been deleted at both ends.

\begin{lemma}\label{l:response}
At a \(\mathrm U\) departure \(y\notin S\), write
\(r_{ab}=\mathbb E^\tau_{E'}P_{ab}-\mathbb E^0_{E'}P_{ab}\).
With three present neighbours,
\[
 \log\rho_j^0-\log\rho_j^\tau
       =\alpha_jr_{fg}+\beta_j(r_{fo}+r_{go}),\qquad
 0<\alpha_j\le\bar\alpha=\frac{18094}{100000},\quad
 0<\beta_j\le\bar\beta=\frac{36751}{100000}.
\]
With two present neighbours the response is \(\beta_j'r_{fo}\),
where \(0<\beta_j'\le35510/100000\le\bar\beta\);
with one it is zero. Also \(\rho_j^0\le1\) and
\(\rho_j^\tau\le12/5\).
A \(\mathrm D\) departure has zero logarithmic response.
\end{lemma}

\begin{proof}
For positive \(x,y\), define
\[
 L(x,y)=\int_0^1x^t y^{1-t}\,\mathrm dt
 =\begin{cases}
 (x-y)/(\log x-\log y),&x\ne y,\\
 x,&x=y.
 \end{cases}
\]
The equality follows by integrating the exponential when \(x\ne y\).
The integral shows that \(L\) is positive, increasing in each argument,
and between the minimum and maximum of its arguments. In both cases
\(\log x-\log y=(x-y)/L(x,y)\).

First suppose that the present neighbours are \(f,g,o\).
Write \(I^\bullet=\mathbb E^\bullet_{E'}I\),
\(J^\bullet=\mathbb E^\bullet_{E'}J\), and
\(L_I=L(I^\tau,I^0)\), \(L_J=L(J^\tau,J^0)\).
Using the ratio formula in Definition~\ref{d:ratios} gives
\[
 \begin{split}
 \log\rho_j^0-\log\rho_j^\tau
 &=\frac{I^\tau-I^0}{L_I}
       -\frac{J^\tau-J^0}{L_J}\\
 &=\left(\frac1{3L_I}-\frac1{5L_J}\right)r_{fg}
       +\frac1{3L_I}(r_{fo}+r_{go}).
 \end{split}
\]
We use these displayed expressions to define \(\alpha_j,\beta_j\),
also when a response or both arguments of a logarithmic mean coincide.

Since \(y\notin S(G)\), it has distance at least three from
\(\mathbb Z^2\setminus G\). Each of its neighbours has distance
at least two from that complement and carries no original leg.
After \(y\) is removed, each has at most three present bonds.
The neighbours are pairwise nonadjacent. On every positive fibre
of the other spins in \(E'\), Lemmas~\ref{l:finitefacts}
and~\ref{l:35} therefore give independent conditional spins with
means \(m_f,m_g,m_o\) of norms at most \(3/5\), under either
pin choice. Put \(Q=\sum_{\{a,b\}\subset\{f,g,o\}}m_a\cdot m_b\).
Then
\[
 -\frac{27}{50}\le
 \frac{|\sum_a m_a|^2-\sum_a|m_a|^2}{2}=Q
 \le\sum_{\{a,b\}}|m_a||m_b|\le\frac{27}{25}.
\]
Also \(|m_f\cdot m_g|\le9/25\).
Averaging these conditional inequalities proves
\[
 \frac{41}{50}\le I^\bullet\le\frac{34}{25},\qquad
 \frac{116}{125}\le J^\bullet\le\frac{134}{125}.
\]
Every pair insertion used here is admissible for free positivity:
removal of \(y\) leaves at most three factors at each inserted
endpoint, and the pair adds one degree there.
Thus Lemma~\ref{l:finitefacts} gives \(I^0\ge1\), \(J^0\ge1\).
The monotonicity of \(L\) now gives
\[
 L_I\ge L(41/50,1),\quad L_I\le34/25,\qquad
 116/125\le L_J\le134/125.
\]
In particular
\[
 \alpha_j\ge\frac{25}{102}-\frac{25}{116}
       =\frac{175}{5916}>0,\qquad \beta_j>0.
\]

Here are exact rational upper bounds for the logarithms needed in
the upper estimates. For \(0<s<1\), set \(t=(1-s)/(1+s)\).
The substitution giving
\(-\log s=2\int_0^t(1-u^2)^{-1}\,\mathrm du\), together with
\[
 \frac1{1-u^2}=1+u^2+\frac{u^4}{1-u^2},\qquad
 \frac1{1-u^2}\le\frac1{1-t^2}\quad(0\le u\le t),
\]
proves
\[
 \frac{-\log s}{1-s}\le
 B(s):=\frac2{1-s}
       \left(t+\frac{t^3}{3}+\frac{t^5}{5(1-t^2)}\right).
\]
Substitution and rational simplification give
\[
 \begin{array}{c|c|c}
 s&B(s)/3&\text{rational upper bound}\\ \hline
 41/50&340634561/926892330&367502/10^6\\
 22/25&24332081/68523180&35510/100000
 \end{array}
\]
The middle entries are strictly smaller than the right entries;
all denominators are positive. For example their cross-product
differences are respectively \(224059660\) and \(50021800\).
Consequently
\[
 \beta_j\le\frac1{3L(41/50,1)}
       \le\frac{367502}{10^6}<\bar\beta,\qquad
 \alpha_j\le\frac{367502}{10^6}-\frac{25}{134}
       =\frac{6061317}{33500000}<\bar\alpha.
\]
The last strict inequality has difference
\(173/33500000>0\).

With just the neighbours \(f,o\), \(J=1\) and
\[
 I^\bullet\in[22/25,28/25],\qquad I^0\ge1.
\]
The same logarithmic identity gives
\(\beta_j'=1/(3L(I^\tau,I^0))>0\) and response
\(\beta_j'r_{fo}\). The second row of the table proves
\(\beta_j'\le35510/100000\le\bar\beta\).
With only \(o\), both partition ratios are one and the response
is zero.

For the free ratio with three neighbours, free positivity also gives
\[
 I^0-J^0
 =\frac2{15}\mathbb E^0_{E'}P_{fg}
      +\frac13\mathbb E^0_{E'}(P_{fo}+P_{go})\ge0.
\]
For two neighbours use \(I^0-1=\mathbb E^0_{E'}P_{fo}/3\ge0\);
for one neighbour the ratio is one. Hence \(\rho_j^0\le1\)
in all cases. Pointwise, for three neighbours of norm at most one,
\[
 \sum_{\{a,b\}}n_a\cdot n_b
   =\frac{|\sum_a n_a|^2-\sum_a|n_a|^2}{2}\ge-\frac32,
 \qquad I\ge\frac12,\qquad J\le\frac65.
\]
It follows that \(\rho_j^\tau=J^\tau/I^\tau\le12/5\).
For two neighbours, \(I\ge2/3\), \(J=1\), so the bound is
\(3/2\le12/5\); for one it is one.
Finally a \(\mathrm D\) departure has both ratios equal to one
by Definition~\ref{d:ratios}, proving its zero logarithmic response.
\end{proof}

\begin{lemma}\label{l:exposure}
For an admissibly started stopped trail, let \(\Phi\) be its
\(\mathrm U\) departures, and suppose \(\rho_j^\tau\le x_j\),
\(x_j\ge1\), for \(j\in\Phi\).
For every ordering \(\sigma\) of \(\Phi\),
\[
 |R_D^\tau(w)-R_D^0(w)|
 \le\sum_{k=1}^{|\Phi|}
       |\log\rho_{\sigma(k)}^0-\log\rho_{\sigma(k)}^\tau|
                         \prod_{l\le k}x_{\sigma(l)}.
\]
Multiplication by \(a_D(w)\) gives the corresponding bound on the
coefficient in the pair expansion.
\end{lemma}

\begin{proof}
All step ratios are positive. The \(\mathrm D\) ratios are one
for both pin choices, so
\[
 R_D^\tau(w)=\prod_{j\in\Phi}b_j,\qquad
 R_D^0(w)=\prod_{j\in\Phi}c_j,\qquad
 b_j=\rho_j^\tau,\quad c_j=\rho_j^0.
\]
For \(m=|\Phi|\), telescoping in the given order gives the exact
identity
\[
 \prod_{j\in\Phi}b_j-\prod_{j\in\Phi}c_j
 =\sum_{k=1}^m
   \left(\prod_{l<k}b_{\sigma(l)}\right)
   (b_{\sigma(k)}-c_{\sigma(k)})
   \left(\prod_{l>k}c_{\sigma(l)}\right).
\]
For positive \(b,c\), integration of the derivative of
\(\exp t\) between \(\log b\) and \(\log c\) yields
\[
 |b-c|\le\max\{b,c\}\,|\log b-\log c|.
\]
Lemma~\ref{l:response} gives \(c_j\le1\), and the hypothesis
gives \(b_j\le x_j\), with \(x_j\ge1\).
Thus \(\max\{b_j,c_j\}\le x_j\), and the absolute value of
the \(k\)-th telescoping term is at most
\[
 |\log\rho_{\sigma(k)}^0-\log\rho_{\sigma(k)}^\tau|
          \prod_{l\le k}x_{\sigma(l)}.
\]
Sum these inequalities. If \(m=0\), both products are one and
the sum is empty, so the same assertion holds.
Finally \(a_D(w)>0\), and multiplying the inequality by it
gives the stated coefficient bound.
\end{proof}

\begin{definition}\label{d:pattern}
Let \(\mathsf T=\{z:d(z,0)\le2\}\), with passage site \(0\),
arrival site \(\mathsf s=(0,-1)\), other neighbours
\(\mathsf e=(1,0),\mathsf N=(0,1),\mathsf W=(-1,0)\),
four corners \((\pm1,\pm1)\), and tips \((\pm2,0),(0,\pm2)\).
The external slots are the two external bonds at each corner and
the three at each tip, twenty in total.
A pattern chooses an out-neighbour \(o\in\{\mathsf e,\mathsf N,\mathsf W\}\)
and statuses/kept sets with the following rules.
The passage site is \(\mathrm U\) and keeps precisely its three
non-arrival bonds. The arrival site is removed or is \(\mathrm D\)
keeping two of its three non-passage bonds.
Each other neighbour is \(\mathrm U,\mathrm X\), or \(\mathrm D\)
keeping its passage bond and one other bond; \(o\) is not removed.
Each corner is \(\mathrm U,\mathrm X\), or \(\mathrm D\) keeping
two of its four bonds. Every \(\mathrm U\) site other than the
passage site keeps all four bonds.
Between nonremoved non-tip sites, a bond is kept by both ends or
by neither. Tips are covering \(\mathrm U\) sites; their bond to
a neighbour \(z\) is present exactly when \(z\) is nonremoved and
keeps that bond. No consistency condition with a global tip kept
set is imposed.
For \(m\in(\mathbb S^2)^{20}\), include every external slot factor
at every nonremoved corner and tip, including both corner slots
when the corresponding actual external bonds are missing.
Let \(Z_{\rm cur}(\pi;m)\) be the finite product functional of
all present template bonds and these slots, using the status laws.
Let \(Z_{\rm next}(\pi;m)\) change the passage law to
\(\lambda_{\mathrm D}\) and drop its out-bond.
The resulting class is \(\mathcal P_{\rm ext}\); it is straight
when \(o=\mathsf N\), and turning otherwise.
\end{definition}

Both local sums are strictly positive for every unit slot choice.
Indeed a corner has at most two internal factors and two slots,
a tip has at most one internal factor and three slots, a neighbour
has at most its four internal factors, and the passage site has
at most three before the step and two after it. Removed sites
contribute no variable or slot. Lemma~\ref{l:finitefacts}
therefore applies to each sum. Every slot vector occurs once,
affinely; this property also applies when a slot at a removed site
has no effect.

\begin{definition}[(H1)]\label{h:H1}
For \(c_{\mathrm s},c_{\mathrm t}\in(0,1]\), (H1) means that
for every \(\pi\in\mathcal P_{\rm ext}\) and every
\(m\in(\mathbb S^2)^{20}\),
\[
 Z_{\rm cur}(\pi;m)\ge c_{\rm type(\pi)}Z_{\rm next}(\pi;m).
\]
In the paper the fixed thresholds are
\(c_{\mathrm s}=961/1000\), \(c_{\mathrm t}=187/200\).
No minimum, minimizing pins, or optimal inversion is asserted.
\end{definition}

We use (H1) through its pointwise inequality on the full covered
pattern class of Definition~\ref{d:pattern}.

\begin{lemma}\label{l:local}
At any \(\mathrm U\) departure outside \(S\) whose normalized
non-tip template data belong to \(\mathcal P_{\rm ext}\), (H1)
implies \(\rho_j^\tau\le1/c_{\rm type}\).
The conclusion is uniform in all original unit pins, all outside
spins of norm at most one, all missing external factors, and all
actual tip statuses and kept sets allowed by the ordinary-domain rules.
\end{lemma}

\begin{proof}
Let \(E\) and \(E^+\) be the pre-step and post-step networks,
and let \(\mathsf T(y)=y+\mathsf T\) be the template.
Because \(y\notin S(G)\), this template lies in \(G\).
Apply a lattice symmetry to put the incoming direction and the
out-neighbour in the positions of Definition~\ref{d:pattern}.
This just relabels the sites and their spins.

Fix the outside spins at \(\omega\). Split off all site weights
and factors involving only sites outside the template; their
product, denoted \(H(\omega)\), is nonnegative and identical
for \(E\) and \(E^+\). Let \(C(\omega)\), \(N(\omega)\) be the
remaining template sums for \(E\), \(E^+\), respectively.
There is a factor \(1+n_z\cdot b\) for each actual external
bond or original leg of a nonremoved template site \(z\), with
\(|b|\le1\). Original pins have \(|b|=1\).
Each external direction carries at most one such factor.
No external factor is incident to the passage site or its
neighbours, since these sites are at distances at least three
and two, respectively, from the original complement.

We first replace the actual tips. A tip \(t\) has precisely one
possible neighbour \(z\) inside the template; its other three
directions are external slots. If \(t\) is nonremoved, write
\[
 W_t(n)=\prod_{\text{actual external factors at }t}(1+n\cdot b),
 \quad A_t=\lambda_{\mathrm{st}(t)}(W_t)>0,\quad
 b_t=\frac{\lambda_{\mathrm{st}(t)}(nW_t)}{A_t}.
\]
There are at most three external factors, so \(A_t>0\) and
\(|b_t|\le3/5\), by the proof of Lemma~\ref{l:35}.
If its internal bond is present, summing over this tip gives
\(A_t(1+b_t\cdot n_z)\); if the bond is absent it gives \(A_t\).
A removed tip contributes the constant one.
These sums and constants are identical in the current and next
template, since the step changes only the passage site and its
out-bond.

In the covering pattern the bond from this tip to \(z\) is
required exactly when \(z\) is nonremoved and keeps it.
If that bond is required but the actual tip has no present bond,
use the desired effective mean zero and the actual constant
as \(A_t\), taking \(A_t=1\) for a removed tip.
If it is required and actually present, use the effective mean
\(b_t\) just computed. In either case Lemma~\ref{l:35}
supplies three unit external fields for a covering
\(\mathrm U\) tip with that exact mean. Let \(A'_t>0\) be
its external three-factor normalizer. Its integrated factor
is the desired factor multiplied by \(A'_t/A_t\).
If the covering bond is not required, the actual bond is
absent as well; choose any three unit fields for the covering
tip and use the same positive normalizer adjustment.
Thus replacing all four tips gives a covered-tip template whose
current and next sums are multiplied by the same strictly
positive scalar \(\prod_t A'_t/A_t\).
The zero mean and the full ball of means used here are attained
by the explicit moment formula in Lemma~\ref{l:35}.

We next fill and resolve the corner slots. For each external
direction at a nonremoved corner, use its actual field
\(b\) if present, and use \(b=0\) if absent.
For any \(|b|\le1\), choose a unit \(e\) parallel to \(b\)
when \(b\ne0\), and any unit \(e\) when \(b=0\). Then
\[
 1+n_z\cdot b
 =\frac{1+|b|}{2}(1+n_z\cdot e)
       +\frac{1-|b|}{2}(1-n_z\cdot e).
\]
This is also the filling identity for an absent factor.
Expanding independently over the at most eight corner slots
expresses the current and next sums as convex combinations
with identical coefficients of sums with unit vectors in every
corner slot. Both slots are filled at every nonremoved corner,
regardless of its actual external kept bonds.
At removed corners the slot values can be assigned arbitrarily.

By the premise, the non-tip statuses and internal bonds are
exactly those of a pattern \(\pi\in\mathcal P_{\mathrm{ext}}\).
After the two replacements, every term is therefore
\(Z_{\rm cur}(\pi;m)\) or \(Z_{\rm next}(\pi;m)\) for the same
unit slot assignment \(m\); its type is the type of the given
step. Apply (H1) term by term, average with the nonnegative
corner coefficients, and undo the common positive tip scalar.
This proves
\(C(\omega)\ge c_{\rm type}N(\omega)\)
for every outside configuration \(\omega\).
Multiply by \(H(\omega)\) and sum over the outside spins.
The result is
\[
 Z_E^\tau\ge c_{\rm type}Z_{E^+}^\tau.
\]
Both partition functions are positive by
Lemma~\ref{l:finitefacts}, so division proves
\(\rho_j^\tau=Z_{E^+}^\tau/Z_E^\tau\le1/c_{\rm type}\).
The construction allowed each fixed outside field to range
over the whole closed unit ball, and every missing slot and
allowed tip status was treated explicitly. This proves all
the asserted uniformity.
\end{proof}

\begin{definition}\label{d:weights}
Put \(x_{\mathrm s}=1/c_{\mathrm s}\), \(x_{\mathrm t}=1/c_{\mathrm t}\).
A lattice trail \(w=(v_0,\ldots,v_n)\) has distinct unoriented
nearest-neighbour bonds. Give it a designated unit incoming direction
\(a\), and forbid \(v_1-v_0=-a\) when \(n\ge1\); the virtual
incoming bond is not itself a used bond of the trail.
For departure \(i\), \(0\le i<n\), give weight \(1/5\) if \(v_i\)
has been departed from earlier, and otherwise weight \(x_{\mathrm s}/3\)
for continuing in the incoming direction and \(x_{\mathrm t}/3\)
for a turn. Let \(\mathsf W(w)\) be the product, with
\(\mathsf W(\varnothing)=1\).
\end{definition}

At a later departure the incoming direction is
\(v_i-v_{i-1}\); at \(i=0\) it is \(a\).
Distinctness of the unoriented bonds forbids immediate reversal
at every later step. The first-step rule is an additional
restriction specified by the incoming direction.
In a trace, a site which has been departed from earlier cannot
again have status \(\mathrm U\). At any subsequent departure
from that site it is \(\mathrm D\), with coefficient \(1/5\)
and ratio one. A site initially of status \(\mathrm D\) may
instead have its first lattice departure weighted as a first
visit; since \(c_{\rm type}\le1\),
\[
 \frac15\le\frac13\le\frac{x_{\rm type}}3.
\]
This is the comparison for initially \(\mathrm D\) sites.

\begin{definition}[(H2)]\label{h:H2}
For \(C_V\ge1\) and \(0<q<1\), (H2) means that for every
integer \(n\ge0\), every start and every designated incoming direction,
\[
 \sum_{\substack{|w|=n\\\text{initial reversal forbidden}}}
                      \mathsf W(w)\le C_Vq^n .
\]
The sum is over plane lattice trails with those specified initial data.
\end{definition}


\begin{lemma}\label{l:onepoint}
Assume (H1) and (H2). Start a trace at an admissible \(x\notin S\)
with incoming direction \(a\) such that \(\{x-a,x\}\) is absent.
Suppose each stopped trail from \(x\) to \(S\) has at most twelve
\(\mathrm U\) departures whose templates are not covered patterns.
Then
\[
 \sum_{w\in\mathcal T_D(x,S)}a_D(w)
                  \max\{R_D^\tau(w),R_D^0(w)\}
 \le (12/5)^{12}\frac{C_V}{1-q}\,q^{d(x,S)}.
\]
Pointwise the summand is at most \((12/5)^{12}\mathsf W(w)\).
\end{lemma}

\begin{proof}
The absent bond \(\{x-a,x\}\) ensures that every trail of the
network from \(x\) has the forbidden initial reversal required
in Definition~\ref{d:weights}. Subsequent steps have distinct
unoriented bonds, so every stopped trail in the left-hand side
is among the plane trails in (H2) with these same initial data.

Fix such a stopped trail \(w\), and let \(\omega_j\) be its
\(j\)-th lattice weight in \(\mathsf W(w)\).
We bound its coefficient for both choices of \(\bullet\)
by the same product of step majorants.
At a \(\mathrm D\) departure the product of the departure
coefficient and \(\rho_j^\bullet\) is exactly \(1/5\).
If this site was departed from earlier in \(w\), this is
exactly \(\omega_j\). If it was not, then
\(\omega_j=x_{\rm type}/3\ge1/3\), so it still bounds \(1/5\).
This includes an initially \(\mathrm D\) start, whose
admissibility ensures at most one present bond.

A \(\mathrm U\) departure is necessarily the first departure
from its site in \(w\). If its pattern is covered,
Lemma~\ref{l:local} gives \(\rho_j^\tau\le x_{\rm type}\)
and Lemma~\ref{l:response} gives
\(\rho_j^0\le1\le x_{\rm type}\). Thus, for both choices,
\[
 \frac13\rho_j^\bullet\le\frac{x_{\rm type}}3=\omega_j.
\]
At an uncovered \(\mathrm U\) departure, the same uniform
response lemma gives \(\rho_j^\tau\le12/5\),
\(\rho_j^0\le1\le12/5\). Therefore
\[
 \frac13\rho_j^\bullet\le\frac45
          \le(12/5)\,\omega_j,
\]
because \(\omega_j\ge1/3\).
There are at most twelve such departures, wherever they occur
along \(w\). Multiplying the step estimates and telescoping
the ratios proves, for both pin choices with the same bound,
\[
 a_D(w)R_D^\bullet(w)
   \le(12/5)^{12}\prod_j\omega_j
   =(12/5)^{12}\mathsf W(w).
\]
Taking the maximum preserves this pointwise inequality.

The length of a nearest-neighbour trail from \(x\) to \(S\)
is at least \(d_0=d(x,S)\). If there are no stopped trails,
the asserted sum is zero. Otherwise, group the finite domain
sum by length and enlarge each group to all plane trails
of that length with the fixed start and incoming direction.
Using (H2), nonnegativity and \(0<q<1\), we obtain
\[
 \begin{split}
 \sum_{w\in\mathcal T_D(x,S)}
       a_D(w)\max\{R_D^\tau(w),R_D^0(w)\}
 &\le(12/5)^{12}\sum_{n\ge d_0}
              \sum_{\substack{|w|=n\\
                    \text{initial reversal forbidden}}}\mathsf W(w)\\
 &\le(12/5)^{12}C_V\sum_{n\ge d_0}q^n\\
 &=(12/5)^{12}\frac{C_V}{1-q}\,q^{d_0}.
 \end{split}
\]
Here \(d_0\) is a finite integer: the existence of \(x\) makes
\(G\) nonempty, and every nonempty finite \(G\subset\mathbb Z^2\)
has a site adjacent to its complement, hence \(S\ne\varnothing\).
This completes the proof.
\end{proof}


\section{Pair closure, a positive-vector witness, and renewal}
\label{s:recursion}

All networks in this section are finite plane networks with the site
laws of Definition~\ref{d:moment} and the original legs of
Definition~\ref{d:pinmodel}.  The boundary used in every distance is
the original protected set \(S=S(G)\).  Ordinary domains have the
meaning of Definition~\ref{d:domain}.  We first close the response
class under tracing, and then derive a nonnegative shift recursion.
The quantitative hypotheses are (H1) and (H2) of
Definitions~\ref{h:H1} and~\ref{h:H2}, and the positive-vector
assertion (H3\(_h\)) defined below.

The trace operations, the stopped trail class, and the coefficients
\(a_E(w),R_E^\bullet(w)\) are those of
Definition~\ref{d:trace}.  The ratios of
Definition~\ref{d:ratios} telescope to
\(R_E^\bullet(w)=\prod_j\rho_j^\bullet\), with
\(\rho_j^\tau=\rho_j^0=1\) at every admissible
\(\mathrm D\) departure.  Proposition~\ref{p:pairwalk} supplies
the finite expansion, its target constants, and its separate-degree
bound.  We use it only with its prescribed stopping set
\(S\cup\{\text{target}\}\) and check its start and target
hypotheses at both expansions below.

The local comparison used here is Lemma~\ref{l:local}, including
its uniformity in outside spins, original unit pins, missing
external factors, and actual tip data allowed by the ordinary-domain
rules.  The ordered product estimate is Lemma~\ref{l:exposure};
the common pointwise majorant and its summed consequence are
Lemma~\ref{l:onepoint}.  These are results of the preceding walk
section.  In particular the last lemma bounds the sum of pointwise
maxima for the two pin choices and permits its twelve exceptional
departures at arbitrary trail positions.

\begin{definition}\label{d:pairenv}
For \(D\in\mathcal D_{\rm cons}(G)\), let \(P_B(D)\) and \(P_C(D)\)
be the ordered nonremoved pairs \((u,v)\) at distance two,
collinear or diagonal respectively, with a designated common neighbour
\(c\) of status \(\mathrm X\). Each endpoint is \(\mathrm U\),
or \(\mathrm D\) with at most one present bond.
For integer \(d\ge0\), let \(F_X(d)\) be the supremum of
\[
 \left|\mathbb E_D^\tau(n_u\cdot n_v)
                    -\mathbb E_D^0(n_u\cdot n_v)\right|
\]
over every finite \(G\), every original unit-pin assignment,
every \(D\in\mathcal D_{\rm cons}(G)\), and every pair in \(P_X(D)\)
with \(\min\{d(u,S(G)),d(v,S(G))\}\ge d\).
Use supremum zero for an empty class, and set \(F_X(d)=F_X(0)\)
for negative integers. Thus \(0\le F_X\le2\), and \(F_X\)
is nonincreasing.
\end{definition}

Indeed, both expectations are expectations under positive finite
probability laws by Lemma~\ref{l:finitefacts}, and
\(|n_u\cdot n_v|\le1\).  Increasing \(d\) restricts the class over
which the supremum is taken.  In particular the asserted bounds and
monotonicity also hold when that class is empty.

\begin{lemma}\label{l:closure}
Let \(D\in\mathcal D_{\rm cons}(G)\), \((u,v)\in P_X(D)\) have
centre \(c\), and \(u\notin S\). Before tracing, if \(u\) is
\(\mathrm U\), drop its kept bond to the removed \(c\); this changes
no partition function. Trace any trail from \(u\) without visiting
\(S\) before its last site; it need not avoid \(v\).
After every prefix ending at the live endpoint \(y\):
all departed sites are \(\mathrm D\) or \(\mathrm X\);
consistency between nonremoved sites holds; all nonlive
\(\mathrm U\) sites have full kept sets and all nonlive \(\mathrm D\)
sites have exactly two kept bonds.
After a positive-length prefix a live \(\mathrm U\) site has lost
only its arrival bond, and a live \(\mathrm D\) site keeps one bond;
the arrival site's other endpoint is \(\mathrm D\) or \(\mathrm X\).
At length zero a live \(\mathrm U\) site has its three non-\(c\)
kept bonds, while a live \(\mathrm D\) site retains its two kept
bonds but has at most one present bond.
Whenever \(y\notin S\), removing \(y\) yields an ordinary member
of \(\mathcal D_{\rm cons}(G)\). Every \(\mathrm U\) departure
at \(y\) before stopping has a covered pattern. In the ordinary
domain obtained by removing the live site \(y\) from the pre-step
network, every ordered pair of distinct nonremoved non-arrival
neighbors of \(y\) belongs to \(P_B\) or \(P_C\), with designated
center \(y\).
If the endpoint is \(e\in S\), it is \(\mathrm U\) with just its
arrival bond lost. This is the sole kept-set defect in the terminal
auxiliary network; all original legs remain. No removal of \(e\)
or ordinary-domain membership after removing \(e\) is asserted.
\end{lemma}

\begin{proof}
Write the trail as \((v_0,\ldots,v_m)\), with \(v_0=u\).
All departures, arrivals, and removals in the following induction
are the operations of Definition~\ref{d:trace}.
Every neighbour of a site outside \(S\) is in \(G\), since such a
site has distance at least three from \(\mathbb Z^2\setminus G\).
Thus the three non-\(c\) kept bonds mentioned at length zero really
exist as bonds inside \(G\).  A bond incident to a removed site is
absent regardless of its kept flag at the other end.  Dropping the
flag at \(u\) for \(\{u,c\}\) therefore changes no factor, pinned
or free, and changes no normalized expectation.  The initial live
site is now the only possible failure of the ordinary kept-set
conditions.  Consistency is unchanged because \(c\) is removed.

We prove the assertions by induction on the prefix length.  At a
departure from a live \(\mathrm U\) site, the incoming bond is absent
and its three other bonds are kept.  The outgoing bond is present and
differs from the incoming bond.  Dropping the outgoing bond and
changing the law to \(\lambda_{\mathrm D}\) leaves exactly two kept
bonds at the departed site.  If the initial site is \(\mathrm D\),
it has at most one present bond, which must be the chosen outgoing
bond.  Its removal therefore leaves no other present bond incident
to it.  At any later live \(\mathrm D\) site there is just one kept
bond, and an actual departure uses that bond and removes the site.
These are all possible departure cases.  In particular a departed
site never becomes \(\mathrm U\) again.

On arrival the new live site cannot be removed, since the traversed
bond was present.  If it is \(\mathrm U\), it was not previously
departed from; before this arrival its kept set was full.  Dropping
the arrival bond leaves precisely the stated live kept set.  If it
is \(\mathrm D\), its nonlive kept set had two elements and contained
the present arrival bond; dropping that bond leaves one element.
All traversed bonds are dropped by both ends.  Removing a
\(\mathrm D\) departure site destroys only bonds to a removed site,
and so cannot create a consistency failure between nonremoved sites.
Every other kept flag is unchanged.  This proves the induction,
including the assertion about the other end of the arrival bond.

Before stopping, no protected site has been visited.  Removing the
live endpoint when it is outside \(S\) removes the only exceptional
kept set; all protected sites still have their original full kept
sets.  The remaining \(\mathrm U\) kept sets are full, the remaining
\(\mathrm D\) kept sets have size two, and consistency holds.
Consequently the resulting network is an ordinary member of
\(\mathcal D_{\rm cons}(G)\).  At a terminal hit \(e\in S\),
the site \(e\) has not previously been visited.  It still has its
original \(\mathrm U\) law and has dropped just the arrival bond.
All other sites satisfy the ordinary kept-set conditions.  Tracing
never changes an original leg factor, including those at \(e\).

It remains to verify the local pattern and pair assertions.  At a
\(\mathrm U\) departure \(y\), the designated arrival site is
\(c\) at the initial departure, and is the preceding trail site
otherwise.  It is removed, or is \(\mathrm D\) with its two kept
bonds among the three bonds other than its bond to \(y\).
Every nonremoved non-arrival neighbour of \(y\) keeps its bond to
\(y\): the live \(\mathrm U\) site keeps that bond, so consistency
forces the other end to keep it.  Such a neighbour of status
\(\mathrm D\) has exactly one additional kept bond.  Every corner
of status \(\mathrm D\) has two kept bonds, and every other
\(\mathrm U\) site has its full kept set.  The outgoing neighbour
is nonremoved because its bond is present.  These are exactly the
non-tip rules of Definition~\ref{d:pattern}.  The diamond of radius
two about \(y\) lies in \(G\).  Original legs incident to its
corners or tips occupy external slots.  Actual missing external
factors, both slots at each corner, and actual tip data are treated
by the covering procedure of Lemma~\ref{l:local}; that procedure
does not require consistency with an actual tip kept set.

Finally let \(p\) be a pair of distinct nonremoved non-arrival
neighbours of \(y\).  The two sites are at distance two and are
collinear or diagonal.  Removing \(y\) supplies their designated
removed centre.  Each endpoint of status \(\mathrm D\) previously
had its bond to \(y\) present, so after removal it has at most one
present bond.  The endpoints therefore meet
Definition~\ref{d:pairenv} in the ordinary domain just obtained.
None of this reasoning used avoidance of \(v\).
\end{proof}


\begin{lemma}\label{l:terminalcover}
In the first expansion of a pair of Definition~\ref{d:pairenv},
assume also that \(v\notin S\), and
suppose the trail avoids its target \(v\) and stops at \(e\in S\).
In the terminal auxiliary network, start a second expansion from \(v\)
to \(S\), again dropping its kept bond to \(c\) if it is \(\mathrm U\).
The starts of both expansions have their designated incoming bond
absent and are admissible in Definition~\ref{d:trace}.
Before the second trail hits \(S\), a \(\mathrm U\) departure
outside \(\{y':0<d(y',e)\le2\}\) has a covered pattern.
There are at most twelve exceptional \(\mathrm U\) departures
on each second trail. Both expansions satisfy Lemma~\ref{l:onepoint}.
\end{lemma}

\begin{proof}
The first expansion has an admissible start outside \(S\) by
Proposition~\ref{p:pairwalk}.  Its incoming bond is \(\{c,u\}\),
which is absent because \(c\) is removed.  At \(u\), a
\(\mathrm U\) law has at most three present bonds; a
\(\mathrm D\) law has at most one by the pair definition.

By Lemma~\ref{l:closure}, the terminal network has precisely the
fixed defect at \(e\) in addition to the ordinary data elsewhere.
The first trail never visits \(v\), so its status and kept set
have not changed.  Its number of present bonds can only have
decreased.  Since \(v\notin S\), it is an admissible second start:
it is \(\mathrm U\) with at most three present bonds, or
\(\mathrm D\) with at most one.  Its designated incoming bond is
\(\{c,v\}\), again absent.  Dropping that kept flag at a
\(\mathrm U\) start leaves every partition function unchanged.
Thus neither expansion can initially reverse across its designated
incoming bond.

For completeness the second pair expansion, with target \(e\),
also meets the target degree hypothesis in
Proposition~\ref{p:pairwalk}.  Before its first arrival, the
protected site \(e\) had at most four factors in total, including
original legs.  The first arrival dropped one present bond, leaving
at most three factors.  Its law remains \(\lambda_{\mathrm U}\).
The second start is distinct from \(e\) because \(v\notin S\).

Repeat the kept-set induction of Lemma~\ref{l:closure} during the
second trace, retaining \(e\) as a fixed defective site.  Before a
hit of \(S\), no second departure is at \(e\), and no other
protected site changes its kept set.  Outside \(e\) and the live
endpoint, the full \(\mathrm U\) and two-element \(\mathrm D\)
kept-set rules continue to hold.  Consistency holds for bonds not
incident to \(e\).  This follows in each departure and arrival
case from the same two flag changes as in the preceding induction;
no operation changes a leg factor.

If \(d(y',e)>2\), the defective site is outside the template
about a \(\mathrm U\) departure \(y'\).  Its non-tip data
therefore satisfy every rule in Definition~\ref{d:pattern}.
Any absent external factor is allowed by Lemma~\ref{l:local},
including a factor whose absence is caused by the defect outside
the template.  The actual tips are also distinct from \(e\),
so their statuses and kept sets satisfy that lemma's ordinary-domain
tip rules.  Every outside spin has norm at most one; the original
leg factors are retained.  No ordinary-domain assumption is imposed
on the whole terminal auxiliary network.  Since no departure is
at \(e\), every possible
exception lies in
\(\{y':0<d(y',e)\le2\}\).  There are four sites at distance
one and eight at distance two in the square lattice.  A site can
be departed from with status \(\mathrm U\) at most once, since
that departure changes its status to \(\mathrm D\).  Hence there
are at most twelve exceptions, at arbitrary positions along the
second trail.

The first trace has no exceptions by Lemma~\ref{l:closure},
even if its stopping rule is enlarged to all first hits of \(S\).
Both traces thus have the absent incoming bond, admissible start,
and exception bound required by Lemma~\ref{l:onepoint}.  Whenever
(H1) and (H2) are assumed, its estimate applies to each trace.
\end{proof}

\begin{remark}
The additional hypothesis \(v\notin S\) in
Lemma~\ref{l:terminalcover} is necessary.  For example, take
\(G=\{-3,\ldots,3\}^2\), so
\(S=\{x\in G:\|x\|_\infty\ge2\}\).  Give \(c=(1,0)\)
status \(\mathrm X\) and every other site status \(\mathrm U\)
with its full kept set.  The pair \(u=(0,0),v=(2,0)\) belongs to
\(P_B\), and the first trail
\(((0,0),(0,1),(0,2))\) avoids \(v\) and stops in \(S\).
Nevertheless \(v\in S\), so it cannot be the start of a second
trace of Definition~\ref{d:trace}.  In the recursion below both
endpoints have distance at least one from \(S\), as required.
\end{remark}

\begin{lemma}\label{l:log-mean-bounds}
Let \(\mathcal L(x,y)=(x-y)/(\log x-\log y)\) for \(x\ne y\)
and \(\mathcal L(x,x)=x\), for positive \(x,y\).
Put \(f(s)=1/\mathcal L(s,1)\), so \(f(1)=1\).
For \(3/5\le c\le1\), the function
\[
 G_c(x,y)=\frac1{3\mathcal L(cx,y)}-\frac1{5\mathcal L(x,y)}
\]
is nonincreasing in each positive variable. In particular,
for \(x\ge a=116/125\), \(y\ge1\),
\[
 G_c(x,y)\le f(ac)/3-f(a)/5.
\]
For rational \(0<s<1\), put \(t=(1-s)/(1+s)\) and
\[
 f_-(s)=\frac2{1-s}\sum_{k=0}^{7}\frac{t^{2k+1}}{2k+1},
 \qquad
 f_+(s)=f_-(s)+\frac{2t^{17}}{17(1-t^2)(1-s)}.
\]
Then \(f_-(s)\le f(s)\le f_+(s)\).
For the two rows of the following table, all numbers are exact rationals:
\[
 \begin{array}{c|c|ccc}
 \tau&c_\tau&a_\tau&b_\tau&b'_\tau\\ \hline
 {\rm s}&961/1000&145217/10^6&352782/10^6&340008/10^6\\
 {\rm t}&187/200&150000/10^6&357549/10^6&344661/10^6
 \end{array}
\]
They satisfy
\[
 \frac{f_+(ac_\tau)}3-\frac{f_-(a)}5\le a_\tau,\qquad
 \frac{f_+(ac_\tau)}3\le b_\tau,\qquad
 \frac{f_+(c_\tau)}3\le b'_\tau,\qquad
 0<b'_\tau\le b_\tau .
\]
\end{lemma}

\begin{proof}
For positive \(r,s\), direct integration, also when \(r=s\), gives
\begin{equation}\label{e:rec-log-integral}
 \frac1{\mathcal L(r,s)}
       =\int_0^1\frac{\mathrm dt}{(1-t)r+ts}.
\end{equation}
In particular the reciprocal logarithmic mean is decreasing in each
argument, and \(\mathcal L\) is increasing in each argument.
For fixed \(t\), put \(A=(1-t)x\) and \(B=ty\).
The integrand for \(G_c\) is
\(1/(3(cA+B))-1/(5(A+B))\).  Its derivatives with respect to
\(x\) and \(y\) are respectively
\[
 (1-t)\left(\frac1{5(A+B)^2}
                   -\frac{c}{3(cA+B)^2}\right),\qquad
 t\left(\frac1{5(A+B)^2}
                   -\frac1{3(cA+B)^2}\right).
\]
Both are nonpositive.  Indeed \(cA+B\le A+B\), \(5c\ge3\),
and hence
\[
 5c(A+B)^2\ge3(cA+B)^2,\qquad
 5(A+B)^2\ge3(cA+B)^2.
\]
On every compact positive rectangle the denominators are bounded
away from zero, justifying differentiation under the integral in
\eqref{e:rec-log-integral}.  This proves both monotonicities and
the value at the lower endpoints \((a,1)\).

For the logarithm enclosure, \(0<t<1\) and
\(-\log s=\log((1+t)/(1-t))\).  Integrating the identity
\[
 \frac1{1-z^2}=\sum_{k=0}^7z^{2k}+\frac{z^{16}}{1-z^2}
       \qquad(0\le z\le t)
\]
from zero to \(t\) shows that
\[
 0\le-\log s-2\sum_{k=0}^7\frac{t^{2k+1}}{2k+1}
   =2\int_0^t\frac{z^{16}}{1-z^2}\,\mathrm dz
   \le\frac{2t^{17}}{17(1-t^2)}.
\]
Division by \(1-s>0\) proves the asserted enclosure, with no
approximate logarithms.

Here are explicit rational checks for all the numerical endpoints.
Let \(K=10^9\).  The integers in the last two columns satisfy
\(\ell_s/K\le f_-(s)\le f_+(s)\le u_s/K\):
\begin{equation}\label{e:rec-log-table}
\begin{array}{c|c|r|r}
 s&t&\ell_s&u_s\\ \hline
 116/125&9/241&1037827030&1037827031\\
 27869/31250&3381/59119&1058344574&1058344575\\
 5423/6250&827/11673&1072644315&1072644316\\
 961/1000&39/1961&1020022307&1020022308\\
 187/200&13/387&1033980764&1033980765
\end{array}
\end{equation}
For clarity, the table is checked solely by the following integer
formula.  If \(s=u/v\), \(t=p/q\) are the fractions in a row, put
\[
 \begin{gathered}
 D_*=45045=\operatorname{lcm}(1,3,5,7,9,11,13,15),\\
 P(p,q)=\sum_{k=0}^7\frac{D_*}{2k+1}p^{2k+1}q^{14-2k},
 \qquad Q=(v-u)D_*q^{15}.
 \end{gathered}
\]
All these denominators are positive integers, and substitution gives
\[
 f_-(s)=\frac{2vP(p,q)}{Q},\qquad
 f_+(s)=\frac{2v\{17(q^2-p^2)P(p,q)+D_*p^{17}\}}
                    {17(q^2-p^2)Q}.
\]
Thus the two checks in each row are the explicit integer inequalities
\[
 \ell_sQ\le2vKP(p,q),\qquad
 2vK\{17(q^2-p^2)P(p,q)+D_*p^{17}\}
                \le17u_s(q^2-p^2)Q.
\]
Inserting the five displayed pairs \((u/v,p/q)\) verifies each
inequality by sums and products of nonnegative integers.
To check the requested coefficient bounds, write
\(A_\tau=10^6a_\tau\), \(B_\tau=10^6b_\tau\), and
\(B'_\tau=10^6b'_\tau\).  After using
\eqref{e:rec-log-table}, their positive residual numerators are
\[
\begin{array}{c|r|r|r}
 \tau&15000A_\tau-5u_{ac_\tau}+3\ell_a
       &3000B_\tau-u_{ac_\tau}&3000B'_\tau-u_{c_\tau}\\ \hline
 {\rm s}&13215&1425&1692\\
 {\rm t}&259510&2684&2235
\end{array}
\]
The respective positive denominators are \(15K,3K,3K\).
The two values of \(b'_\tau\) are positive and at most the
corresponding \(b_\tau\) by direct comparison of their integer
numerators.  This proves every table inequality.
\end{proof}

\begin{lemma}\label{l:response-typed}
Fix \(c_{\rm s}=961/1000\), \(c_{\rm t}=187/200\) and assume (H1).
At a covered \(\mathrm U\) departure of type \(\tau\in\{{\rm s},{\rm t}\}\),
the response coefficients in Lemma~\ref{l:response} satisfy
\[
 0<\alpha_j\le a_\tau,\qquad 0<\beta_j\le b_\tau
       \quad\hbox{with three present neighbours},\qquad
 0<\beta'_j\le b'_\tau\quad\hbox{with two}.
\]
Here the exact constants are those of Lemma~\ref{l:log-mean-bounds};
in particular \(a_{\rm t}=150000/10^6\).
With \(a=116/125\), their analytic bounds are
\[
 \beta_j\le f(ac_\tau)/3,\qquad
 \beta'_j\le f(c_\tau)/3,\qquad
 \alpha_j\le f(ac_\tau)/3-f(a)/5.
\]
These follow using \(I^\tau\ge c_\tau J^\tau\),
\(I^0\ge J^0\ge1\), and \(J^\tau\ge a\).
Also \(a_\tau\le18094/100000\) and
\(b_\tau,b'_\tau\le36751/100000\).
\end{lemma}

\begin{proof}
Use \(E'\), \(I,J\) from Definition~\ref{d:ratios}, and abbreviate
\(I^\bullet=\mathbb E_{E'}^\bullet I\),
\(J^\bullet=\mathbb E_{E'}^\bullet J\).  A covered departure is
outside \(S\), as in Lemma~\ref{l:local}.  That lemma and (H1)
give \(I^\tau/J^\tau\ge c_\tau\); its stated uniformity covers
the actual tips and all missing external factors, including both
corner slots.  We derive the coefficient identities directly from
Definition~\ref{d:ratios}.  For positive \(A,B\), including
\(A=B\),
\(\log A-\log B=(A-B)/\mathcal L(A,B)\).  Since
\(\rho_j^\bullet=J^\bullet/I^\bullet\), with three present
neighbours this gives
\[
 \begin{split}
 \log\rho_j^0-\log\rho_j^\tau
 &=\frac{I^\tau-I^0}{\mathcal L(I^\tau,I^0)}
       -\frac{J^\tau-J^0}{\mathcal L(J^\tau,J^0)}\\
 &=\left(\frac1{3\mathcal L(I^\tau,I^0)}
            -\frac1{5\mathcal L(J^\tau,J^0)}\right)r_{fg}
       +\frac{r_{fo}+r_{go}}{3\mathcal L(I^\tau,I^0)}.
 \end{split}
\]
Thus the response of Lemma~\ref{l:response} has the canonical
coefficients
\begin{equation}\label{e:rec-typed-coefficients}
 \beta_j=\frac1{3\mathcal L(I^\tau,I^0)},\qquad
 \alpha_j=\frac1{3\mathcal L(I^\tau,I^0)}
                     -\frac1{5\mathcal L(J^\tau,J^0)}.
\end{equation}
This fixes the coefficient choice used below, also when responses
vanish or logarithmic-mean arguments coincide.

With three present neighbours, free positivity gives
\[
 J^0\ge1,\qquad
 I^0-J^0=\tfrac2{15}\mathbb E^0_{E'}P_{fg}
             +\tfrac13\mathbb E^0_{E'}(P_{fo}+P_{go})\ge0.
\]
These uses of Lemma~\ref{l:finitefacts} satisfy its degree
condition: after \(y\) is removed, each neighbour has at most
three incident factors and no original leg, so insertion of one
scalar product has separate degree at most four.  The sites
\(f,g,o\) are pairwise nonadjacent.  On every positive
conditioning fibre of the other spins, their spins are independent
and their conditional means have norm at most \(3/5\) by
Lemmas~\ref{l:finitefacts} and~\ref{l:35}.  Hence
\(\mathbb E^\tau_{E'}P_{fg}\ge-9/25\), and
\[
 J^\tau=1+\tfrac15\mathbb E^\tau_{E'}P_{fg}
          \ge1-\tfrac9{125}=a.
\]
The same conditioning under either pin choice yields
\[
 \frac{41}{50}\le I^\bullet\le\frac{34}{25},\qquad
 \frac{116}{125}\le J^\bullet\le\frac{134}{125}.
\]
Indeed, for the three conditional means \(m_f,m_g,m_o\),
\[
 -\frac{27}{50}\le
 \sum_{\{b,b'\}\subset\{f,g,o\}}m_b\cdot m_{b'}
 =\frac{|m_f+m_g+m_o|^2-|m_f|^2-|m_g|^2-|m_o|^2}{2}
 \le\frac{27}{25}.
\]
Together with \(|m_f\cdot m_g|\le9/25\), averaging proves
both intervals.  Formula~\eqref{e:rec-log-integral} places
\(\mathcal L\) between its positive arguments, so the very same
canonical coefficients satisfy
\[
 \alpha_j\ge\frac{25}{102}-\frac{25}{116}
              =\frac{175}{5916}>0,\qquad \beta_j>0.
\]
All four arguments of the logarithmic means are strictly positive.
By their monotonicity and \(I^\tau\ge c_\tau J^\tau\),
\(I^0\ge J^0\), formula~\eqref{e:rec-typed-coefficients} yields
\[
 \beta_j\le\frac1{3\mathcal L(c_\tau a,1)},\qquad
 \alpha_j\le G_{c_\tau}(J^\tau,J^0)
        \le\frac{f(ac_\tau)}3-\frac{f(a)}5.
\]
The last step uses \(3/5\le c_\tau\le1\) and
Lemma~\ref{l:log-mean-bounds}.

With two present neighbours, \(J=1\) and
\(I=1+P_{fo}/3\).  Free positivity gives \(I^0\ge1\),
whereas local (H1) gives \(I^\tau\ge c_\tau\).
The same divided-difference calculation now gives the response
\(\beta'_jr_{fo}\), where
\(\beta'_j=1/(3\mathcal L(I^\tau,I^0))>0\), so
\(\beta'_j\le f(c_\tau)/3\).  With just one present neighbour
both ratios are one by Definition~\ref{d:ratios}, and there is
no response coefficient.  At a \(\mathrm D\) departure that
definition likewise makes both ratios one, so its response is zero.
The rational enclosures of Lemma~\ref{l:log-mean-bounds} prove
the stated bounds.  Finally the largest \(a_\tau\) is
\(150000/10^6<18094/100000\), and each \(b_\tau,b'_\tau\)
is at most \(357549/10^6<36751/100000\).
\end{proof}

\begin{definition}\label{d:kernel}
For \(X=B,C\), fix \(c=0\), \(u=(-1,0)\), and \(v=(1,0)\)
for \(B\), \(v=(0,1)\) for \(C\).
Let \(\mathcal T_n^X\) be the length-\(n\) plane trails avoiding \(c\),
from \(u\) to their first hit of \(v\), with incoming direction \(u-c\).
A departure is a lattice revisit when its site was departed from earlier.
At a lattice first visit \(i\), put \(x_i=x_{\tau_i}\) according to
its straight/turn type. Of the three neighbours other than the arrival
site, write \(o\) for the outgoing neighbour and \(f,g\) for the others.
A neighbour is dead if it is \(c\) or was departed from twice before \(i\).
The outgoing neighbour of a true trail is not dead.
If \(f,g\) are both live, assign coefficient \(a_{\tau_i}\) to
\(\{f,g\}\), and \(b_{\tau_i}\) to each of \(\{f,o\},\{g,o\}\).
If exactly one of \(f,g\) is live, assign \(b'_{\tau_i}\) to its
pair with \(o\); if neither is live there is no response pair.
The constants are those of Lemma~\ref{l:log-mean-bounds}.
A live pair \(p\) has type \(Y=B,C\) and shift
\[
 s(p)=\max_{z\in p}\min\{d(z,u),d(z,v)\}.
\]
Let \(\kappa_{Y,s}(i)\) sum its assigned coefficients.
For any fixed, \(\theta\)-independent orderings \(\sigma_w\) of lattice
first visits, define
\[
 \begin{split}
 M_{XY,s}
 &=\sum_{n\ge1}\sum_{w\in\mathcal T_n^X}
     3^{-n}(3/5)^{\operatorname{rev}(w)}
     \sum_k\kappa_{Y,s}(\sigma_w(k))
                   \prod_{l\le k}x_{\sigma_w(l)},\\
 M_\theta[X,Y]&=\sum_{s\ge0}M_{XY,s}\theta^{-s}.
 \end{split}
\]
These are nonnegative series. The identity
\(3^{-n}(3/5)^r=3^{-(n-r)}5^{-r}\) retains the mass-free
trace coefficients. This \(M\) is the type-dependent response kernel.
\end{definition}

Until finiteness is proved, each series is understood in
\([0,+\infty]\) as the supremum of its finite subsums.
Regrouping these nonnegative sums therefore requires no prior
real-valued summability assertion.

To justify the assertion about the outgoing neighbour, a site other
than \(u\) that has already been departed from twice has used three
distinct incident bonds: two outgoing bonds and an intervening
arrival bond.  Entering it again uses its fourth bond, leaving no
possible subsequent departure.  It cannot be the final target,
since the target has no earlier departure.  The start \(u\) has
only three available bonds because \(c\) is excluded; after two
departures all three have been used.  Neither case permits an
outgoing step to a dead site on a trail in \(\mathcal T_n^X\).
The incoming direction also forbids the initial reversal, since that
step would enter \(c\).

\begin{definition}\label{d:kernel-order}
Fix
\[
 h=(1,h_C),\quad h_C=\frac{5404319552844595}{4503599627370496},
 \qquad \rho_0=\frac{993015}{10^6},\qquad L=12 .
\]
For a first visit set \(\kappa_h(i)=\sum_{Y,s}\kappa_{Y,s}(i)h_Y\).
For each \(X\in\{B,C\}\) and \(w\in\mathcal T_n^X\), define
\(\sigma_w\) as follows. If \(n\le L\), order its lattice
first-departure indices by increasing
\((x_i-1)/(x_i\kappa_h(i))\), assigning value \(+\infty\) when
\(\kappa_h(i)=0\), and break ties by increasing path index.
If \(n>L\), list straight first departures in path order and
then turning first departures in path order.
Hereafter \(M_{XY,s}\) and \(M_\theta\) use exactly these orderings.
They do not depend on \(\theta\).
\end{definition}

Both \(x_{\rm s}\) and \(x_{\rm t}\) exceed one, and \(h_C>1\).
Thus every finite ordering key has a positive denominator.  The
infinite-key rule and the index tie-break give a unique ordering on
each finite trail.  All data in these comparisons are rational.

\begin{definition}[(H3\(_h\))]\label{h:H3}
For the type-dependent kernel and fixed orderings of
Definitions~\ref{d:kernel}--\ref{d:kernel-order}, (H3\(_h\)) is the
following mathematical assertion, with the specified \(h,\rho_0\):
for \(X=B,C\) and every integer \(R_0\ge0\),
\[
 \sum_{n=1}^{R_0}\sum_{w\in\mathcal T_n^X}
 3^{-n}(3/5)^{\operatorname{rev}(w)}
 \sum_k\kappa_h(\sigma_w(k))\prod_{l\le k}x_{\sigma_w(l)}
       \le\rho_0h_X .
\]
Equivalently, the full nonnegative sums satisfy \(M_1h\le\rho_0h\).
The assertion is proved using the stationary lattice Bellman
certificates, with no periodic spatial walk or remainder.
\end{definition}

The equivalence follows directly by taking increasing finite
length cutoffs.  For each length there are finitely many trails,
and each trail has finitely many first visits and response pairs.
Rearrangement by \((Y,s)\) preserves a nonnegative sum.  Consequently
the supremum of the displayed finite sums is exactly
\(\sum_{Y,s}M_{XY,s}h_Y\).  A bound for the full sum bounds each
cutoff, and bounds for every cutoff bound their supremum.

\begin{lemma}\label{l:Mtheta}
Assume (H2) and (H3\(_h\)) for the type-dependent kernel.
For every \(q^2<\theta\le1\), all entries of \(M_\theta\) are
finite. There exist \(q^2<\theta<1\) and \(\rho_0<\rho<1\)
such that the same vector \(h\) satisfies \(M_\theta h\le\rho h\).
Every shift of a length-\(n\) trail is at most \(n/2+1\), and
each matrix entry is bounded by
\[
 \sum_{n\ge1} n(\bar\alpha+2\bar\beta)
                      C_V q^n\theta^{-n/2-1}.
\]
\end{lemma}

\begin{proof}
Let \(w=(v_0,\ldots,v_n)\in\mathcal T_n^X\).  At departure
index \(i\),
\[
 d(v_i,u)\le i,\qquad d(v_i,v)\le n-i.
\]
The minimum of these two distances is at most \(n/2\).  Each
site in a response pair is a neighbour of \(v_i\); the triangle
inequality therefore gives \(s(p)\le n/2+1\).
Put \(C=\bar\alpha+2\bar\beta>0\).  At each first visit, the
sum of all its response coefficients is at most \(C\), since
the rational coefficient table has \(a_\tau\le\bar\alpha\) and
\(b_\tau,b'_\tau\le\bar\beta\).  There are at most \(n\) first
visits.  Since every \(x_i\ge1\), an initial product in any
ordering is at most the product over all first visits.  Hence
the length-\(n\) contribution to a fixed entry of \(M_\theta\)
is at most
\[
 nC\theta^{-n/2-1}
       \sum_{w\in\mathcal T_n^X}\mathsf W(w)
       \le nCC_Vq^n\theta^{-n/2-1}.
\]
The weights of Definition~\ref{d:weights} agree because
\(3^{-n}(3/5)^{\operatorname{rev}(w)}\prod_i x_i
 =\mathsf W(w)\).  These trails form a subset of the open trails
with the same initial data in (H2).  Since
\(q/\sqrt\theta<1\), the displayed series converges; explicitly
it equals
\(CC_V\theta^{-1}(q/\sqrt\theta)/(1-q/\sqrt\theta)^2\).
This part uses only (H2).

We now give an explicit choice transporting the witness.  Set
\[
 H=\frac{\max h}{\min h},\quad
 b=\frac{1+q}{2},\quad r=\frac qb,\quad
 B_*=2HCC_Vb^{-2}\frac{r}{(1-r)^3},
\]
and define
\begin{equation}\label{e:rec-shift-choice}
 \varepsilon=\min\left\{\frac{1-b}{2},
                         \frac{1-\rho_0}{2B_*}\right\},
 \qquad \theta=1-\varepsilon,\qquad
 \rho=\frac{1+\rho_0}{2}.
\end{equation}
Here \(0<q<b<1\), \(0<r<1\), \(B_*>0\), and
\(0<\varepsilon\le(1-b)/2\).  Thus
\(q^2<q<b<\theta<1\), and \(\rho_0<\rho<1\).
For every integer shift \(0\le s\le n/2+1\le n+1\),
the mean value theorem on \([\theta,1]\) gives
\[
 0\le\theta^{-s}-1
     \le(1-\theta)s b^{-s-1}
     \le(1-\theta)(n+1)b^{-n-2}.
\]
This includes \(s=0\), whose left side is zero.
For a fixed row \(X\), multiply each coefficient by
\(h_Y/h_X\le H\), use the same length bound, and sum.  All
summands are nonnegative and finite by the preceding estimate, so
\[
 \frac{((M_\theta-M_1)h)_X}{h_X}
 \le (1-\theta)HCC_Vb^{-2}
                  \sum_{n\ge1}n(n+1)r^n
 = (1-\theta)B_*\le\frac{1-\rho_0}{2}.
\]
The series identity follows by differentiating the geometric series
twice, or by multiplying its partial sums by \((1-r)^3\) and
taking the limit.  Adding \((M_1h)_X/h_X\le\rho_0\) proves
\(M_\theta h\le\rho h\) for the same \(h\) and orderings.
\end{proof}

\begin{lemma}\label{l:recineq}
Fix the two H1 thresholds above. Assume (H1) and (H2), use
the bounds of Lemma~\ref{l:response-typed}, and use the
type-dependent kernel with the fixed orderings of
Definition~\ref{d:kernel-order}.
For \(X=B,C\) and every integer \(d\ge1\),
\[
 F_X(d)\le\sum_{Y=B,C}\sum_{s\ge0}M_{XY,s}F_Y(d-s)+C_Tq^{2d},
 \qquad
 C_T=2(12/5)^{24}\left(\frac{C_V}{1-q}\right)^2 .
\]
\end{lemma}

\begin{proof}
Fix a finite \(G\), unit pins \(\tau\), an ordinary consistent
domain \(D\), and \((u,v)\in P_X(D)\) with designated centre
\(c\), such that
\[
 d_0=\min\{d(u,S),d(v,S)\}\ge d\ge1.
\]
In particular \(u,v\notin S\).  Normalize the kept flag from
\(u\) to \(c\) when \(u\) is \(\mathrm U\), as in
Lemma~\ref{l:closure}.  We continue to denote this normalized
starting network by \(D\); all its partition functions and
expectations agree with those of the original ordinary domain.
Both endpoints meet the restrictions of Proposition~\ref{p:pairwalk}:
a \(\mathrm U\) endpoint has no original leg and at most three
present bonds because \(c\) is removed, and a \(\mathrm D\)
endpoint has at most one present bond by the pair definition.

Fix once and for all a plane lattice symmetry \(\varphi\) sending
the designated triple \((c,u,v)\) to the canonical triple for
type \(X\) in Definition~\ref{d:kernel}.  Explicitly, for \(B\)
choose a signed coordinate permutation \(Q\) with
\(Q(u-c)=(-1,0)\); then \(v-c=-(u-c)\) implies
\(Q(v-c)=(1,0)\).  For \(C\), the vectors \(u-c,v-c\) are
perpendicular lattice unit vectors, so choose the signed coordinate
permutation with \(Q(u-c)=(-1,0)\) and \(Q(v-c)=(0,1)\).
In both cases set \(\varphi(z)=Q(z-c)\).  This choice depends
only on the original ordered pair and its designated centre,
and is used for every trail in both sums below.  In particular
\(w\mapsto\varphi(w)\) is injective.

Use that proposition for both choices \(\bullet=\tau,0\), with
stopping set \(S\cup\{v\}\).  The set of trails and their
coefficients \(a_D(w)\) depend only on the status and bond data,
and hence are common to the two expansions.  Split them into
\(\mathcal W_v\), ending at \(v\), and \(\mathcal W_S\),
ending at a point of \(S\) while avoiding \(v\).  The terminal
factor for \(w\in\mathcal W_v\) is a common
\(\eta_v\in\{1,3/5\}\), specifically \(1\) for a
\(\mathrm U\) target and \(3/5\) for a \(\mathrm D\) target,
by Proposition~\ref{p:pairwalk}.  For \(w\in\mathcal W_S\) ending
at \(e\), put
\(X^\bullet(w)=\mathbb E_{D[w]}^\bullet(n_e\cdot n_v)\).
The response equals \(T_1+T_2\), where
\begin{align*}
 T_1&=\sum_{w\in\mathcal W_v}a_D(w)
                 (R_D^\tau(w)-R_D^0(w))\eta_v,\\
 T_2&=\sum_{w\in\mathcal W_S}a_D(w)
              \{R_D^\tau(w)X^\tau(w)-R_D^0(w)X^0(w)\}.
\end{align*}
These sums are finite: every trail uses distinct bonds from a
fixed finite network.

For \(T_1\), fix a trail \(w=(v_0,\ldots,v_n)\in\mathcal W_v\).
Let \(\Phi_U(w)\) be its departure indices whose current status
is \(\mathrm U\).  At each such index Lemmas~\ref{l:closure}
and~\ref{l:local} give a covered pattern and
\(\rho_i^\tau\le x_{\tau_i}\).  Its lattice type is computed
using the same incoming direction as the trace, including
\(u-c\) at the initial departure.  Order \(\Phi_U(w)\) by
restricting \(\sigma_w\), after identifying the pair with the
canonical pair by the fixed symmetry \(\varphi\); thus here
\(\sigma_w\) means the ordering of \(\varphi(w)\).
The trail is admissibly started with absent incoming bond
\(\{c,u\}\), and \(x_{\tau_i}\ge1\).  These are precisely the
start and product-bound hypotheses of Lemma~\ref{l:exposure}.
Lemma~\ref{l:exposure}, followed by Lemma~\ref{l:response-typed},
then gives
\begin{equation}\label{e:rec-exposed}
 a_D(w)|R_D^\tau(w)-R_D^0(w)|
 \le a_D(w)\sum_{i\in\Phi_U(w)}
       \left(\sum_{p\in\mathcal Q_i}\gamma_{i,p}|r_{i,p}|\right)
       \prod_{\substack{j\in\Phi_U(w)\\j\preceq_{\sigma_w}i}}
                           x_{\tau_j}.
\end{equation}
Here \(\mathcal Q_i\) is the set of nonremoved pairs among
the non-arrival neighbours of the departure site \(y=v_i\),
and \(r_{i,p}\) is its pinned-minus-free scalar-product
expectation in the network with \(y\) removed.  With three
present neighbours, \(\gamma_{i,\{f,g\}}\le a_{\tau_i}\) and
\(\gamma_{i,\{f,o\}},\gamma_{i,\{g,o\}}\le b_{\tau_i}\).
With two present neighbours there is one pair and its coefficient
is at most \(b'_{\tau_i}\).  With one present neighbour the
inner sum is zero.  Departures with status \(\mathrm D\) have
ratio one by Definition~\ref{d:ratios} and zero logarithmic
response, so contribute no term.

By Lemma~\ref{l:closure}, removing \(y\) gives an ordinary
domain \(E'_i\), and each \(p\in\mathcal Q_i\) belongs to
\(P_Y(E'_i)\) for its type \(Y\).  For each \(z\in p\)
choose the nearer endpoint of the original pair.  The triangle
inequality gives
\[
 d(z,S)\ge d_0-\min\{d(z,u),d(z,v)\},\qquad
 \min_{z\in p}d(z,S)\ge d_0-s(p)\ge d-s(p).
\]
If \(d-s(p)<0\), use the constant extension of \(F_Y\).
In all cases the definition of the ordinary-domain supremum yields
\begin{equation}\label{e:rec-pair-envelope}
 |r_{i,p}|\le F_Y(d-s(p)).
\end{equation}

We verify every part of the comparison with the lattice kernel.
The removed site \(c\) cannot be visited, and \(w\) first hits
\(v\) at its endpoint, so the lattice symmetry maps it into
\(\mathcal T_n^X\).  A true \(\mathrm U\) departure is a
lattice first visit, since a previous departure would already have
changed its law to \(\mathrm D\) or removed it.  Every lattice
revisit is a \(\mathrm D\) departure.  Some additional lattice
first visits may also be \(\mathrm D\) departures because the
site was already \(\mathrm D\) in the initial domain.  Thus
\[
 N_{\mathrm D}(w)\ge\operatorname{rev}(w),\qquad
 a_D(w)=3^{-n}(3/5)^{N_{\mathrm D}(w)}
             \le3^{-n}(3/5)^{\operatorname{rev}(w)}.
\]
The coefficient identity here is Definition~\ref{d:trace} with
\(N_{\mathrm U}(w)+N_{\mathrm D}(w)=n\); no partition-ratio
factor has been included in \(a_D(w)\).
Passing from the restricted ordering in \eqref{e:rec-exposed}
to the full ordering only inserts factors \(x_{\tau_j}\ge1\)
before a retained index.  It can also insert additional nonnegative
response terms.

A kernel-dead site is removed in the actual trace: this holds for
\(c\), and two departures change any initially nonremoved site
to \(\mathrm X\) at latest.  Consequently every actual response
pair is kernel-live.  At an actual \(\mathrm U\) departure, all
nonremoved non-arrival neighbours have present bonds by the
consistency assertion of Lemma~\ref{l:closure}.  The following
cases exhaust the coefficient comparison.  If both non-outgoing
neighbours are present, they are both kernel-live, and the kernel
assigns exactly the majorants \(a_{\tau_i},b_{\tau_i},b_{\tau_i}\).
If just one is present, its pair with \(o\) has coefficient
at most \(b'_{\tau_i}\).  The kernel assigns that same majorant
when it knows the other neighbour is dead, and assigns
\(b_{\tau_i}\ge b'_{\tau_i}\) if it still regards the other
neighbour as live.  If neither is present there is no response
term.  The outgoing neighbour is always present.  This also
explains why an additional removal cannot increase the required
kernel coefficient.

Using these comparisons in \eqref{e:rec-exposed} and
\eqref{e:rec-pair-envelope}, summing over \(\mathcal W_v\),
enlarging to all kernel trails, and using \(\eta_v\le1\), gives
\begin{equation}\label{e:rec-T1}
 |T_1|\le\sum_{Y=B,C}\sum_{s\ge0}M_{XY,s}F_Y(d-s).
\end{equation}
Spatial lattice symmetries preserve all distances, straight/turn
types, and collinear/diagonal pair types; relabeling sites and pins
preserves the finite product functionals.  Thus the canonical
kernel applies uniformly to the original pair.  Its unshifted
row sums are finite by the (H2)-only summability estimate in the
proof of Lemma~\ref{l:Mtheta}; (H3\(_h\)) is not needed here.

For \(T_2\), Lemma~\ref{l:terminalcover} makes the second
expansion in \(D[w]\), from \(v\) with target \(e\) and
stopping set \(S=S\cup\{e\}\), admissible.  Thus its stopping
set is exactly the one prescribed by Proposition~\ref{p:pairwalk},
and the start is distinct from its nonremoved target.
That proposition
and \(|n_z\cdot n_e|\le1\) give, for either choice of pins,
\[
 |X^\bullet(w)|
 \le\sum_{w'\in\mathcal T_{D[w]}(v,S)}
                a_{D[w]}(w')R_{D[w]}^\bullet(w')
 \le K_*q^{d(v,S)},\qquad
 K_*=(12/5)^{12}\frac{C_V}{1-q}.
\]
The last inequality is the common majorant of
Lemma~\ref{l:onepoint}, so is uniform in the first trail \(w\).
Its hypotheses hold: (H1) and (H2) are assumed, the incoming
bond \(\{c,v\}\) is absent, the start has the required status
and degree, and there are at most twelve exceptional
\(\mathrm U\) departures by Lemma~\ref{l:terminalcover}.
The bound uses the original \(S\), including when the terminal
auxiliary network has extra removed sites.  Both sets of original
legs remain unchanged.  The degree restriction at target \(e\)
was checked in Lemma~\ref{l:terminalcover}; all inserted
observables in this second expansion therefore have separate
degree at most four.

The ratios are positive.  Hence, pointwise in the first trail,
\[
 |R_D^\tau(w)X^\tau(w)-R_D^0(w)X^0(w)|
    \le2\max\{R_D^\tau(w),R_D^0(w)\}K_*q^{d(v,S)}.
\]
The first-trail sum can be enlarged from \(\mathcal W_S\)
to \(\mathcal T_D(u,S)\), dropping its avoidance of \(v\).
The closure lemma applies to that larger class too.  Applying
Lemma~\ref{l:onepoint} is legitimate because its start has absent
incoming bond \(\{c,u\}\), and all its trails have zero
exceptional \(\mathrm U\) departures by Lemma~\ref{l:closure}.
In particular its pointwise assertion for this larger class is
\[
 a_D(w)\max\{R_D^\tau(w),R_D^0(w)\}
          \le(12/5)^{12}\mathsf W(w).
\]
Taking the maximum therefore precedes summing.  Its summed
assertion gives
\[
 |T_2|\le2K_*q^{d(v,S)}
       \sum_{w\in\mathcal T_D(u,S)}
           a_D(w)\max\{R_D^\tau(w),R_D^0(w)\}
 \le2K_*^2q^{d(u,S)+d(v,S)}\le C_Tq^{2d}.
\]
Combining this with \eqref{e:rec-T1} and taking the supremum
over the defining class of \(F_X(d)\) proves the result.  If
that class is empty its supremum is zero and the same inequality
holds by nonnegativity of the right side.
\end{proof}

In the next lemma \(X,Y\) range over \(\{B,C\}\), and every
shift index \(s\) ranges over the nonnegative integers.

\begin{lemma}\label{l:renewal}
Let \(m_{XY,s}\ge0\), \(h_X>0\), \(0<q^2<\theta<1\),
\(0<\rho<1\), and
\(\sum_{Y,s}m_{XY,s}h_Y\theta^{-s}\le\rho h_X\).
Suppose \(0\le F_X(d)\le2\) on \(\mathbb Z\), with
\(F_X(d)=F_X(0)\) for \(d<0\), and, for integers \(d\ge1\),
\[
 F_X(d)\le\sum_{Y,s}m_{XY,s}F_Y(d-s)+C_Tq^{2d},
 \qquad C_T\ge0.
\]
Writing \(H=\max_Xh_X/\min_Xh_X\), one has, for \(d\ge0\),
\[
 F_X(d)\le A\theta^d,\qquad
 A=\frac{H\max\{2,C_T\}}{1-\rho}\ge2 .
\]
\end{lemma}

\begin{proof}
For a nonnegative function \(f\) on \(\{B,C\}\times\mathbb Z\),
define the nonnegative shift operator
\[
 (Kf)_X(d)=\sum_{Y,s}m_{XY,s}f_Y(d-s).
\]
For an arbitrary nonnegative \(f\), this denotes the extended
nonnegative sum; every application used below is finite by the
unshifted or shifted row bound, as specified in the proof.
Put \(h_{\min}=\min_Xh_X>0\), \(C_*=\max\{2,C_T\}\), and
\[
 g_X(d)=
 \begin{cases}2,&d\le0,\\ C_Tq^{2d},&d\ge1.\end{cases}
\]
The assumed inequality and \(F_X(d)\le2\) give \(F\le KF+g\)
on all integers.  Also
\(\sum_{Y,s}m_{XY,s}h_Y\le\rho h_X\), since
\(\theta^{-s}\ge1\).  In particular all row sums of \(K\)
are finite, so it acts on bounded functions and preserves
nonnegativity and inequalities.  Iterating a finite number of
times yields, for every integer \(k\ge1\),
\begin{equation}\label{e:rec-renewal-iterate}
 F\le K^kF+\sum_{j=0}^{k-1}K^jg.
\end{equation}

View \(h_X\) as constant in \(d\).  The unshifted row inequality
and induction show \(K^kh\le\rho^kh\).  Since
\(F_X(d)\le2h_X/h_{\min}\) for every integer \(d\),
\begin{equation}\label{e:rec-renewal-remainder}
 (K^kF)_X(d)\le\frac{2}{h_{\min}}\rho^kh_X.
\end{equation}
This is an unweighted bound, valid even when the iterated shifts
reach arbitrarily negative integers.

For the source, if \(d\le0\) then \(\theta^d\ge1\), and if
\(d\ge1\) then \(q^{2d}\le\theta^d\).  Thus for every integer
\(d\),
\[
 g_X(d)\le\frac{C_*}{h_{\min}}h_X\theta^d.
\]
Writing \(e_X(d)=h_X\theta^d\), the shifted row inequality is
exactly \(Ke\le\rho e\).  Its sums are finite even for
negative \(d\), by the stated weighted summability hypothesis.
Induction and positivity therefore imply
\[
 (K^jg)_X(d)\le\frac{C_*}{h_{\min}}\rho^jh_X\theta^d.
\]
All rearrangements used here involve nonnegative summands; they
can equivalently be proved first for finite shift sums and then
by increasing limits.  Combining these bounds with
\eqref{e:rec-renewal-iterate}, for fixed \(d\), gives
\[
 F_X(d)\le\frac{2h_X}{h_{\min}}\rho^k
    +\frac{C_*h_X}{h_{\min}}\theta^d
                       \frac{1-\rho^k}{1-\rho}.
\]
Let \(k\) tend to infinity.  The first term tends to zero
because \(0<\rho<1\).  Since \(h_X/h_{\min}\le H\),
the desired bound follows for every \(d\ge0\).
Finally \(H\ge1\), \(C_*\ge2\), and \(0<1-\rho<1\),
so the displayed \(A\) is at least two.
\end{proof}

\begin{theorem}\label{t:rec}
Fix \((c_{\rm s},c_{\rm t})=(961/1000,187/200)\).
Assume (H1), (H2) with any \(C_V\ge1\), \(0<q<1\), and
(H3\(_h\)) for the type-dependent kernel and fixed orderings above.
Then there exist \(A\ge2\) and \(0<\theta<1\) such that
\(\mathsf P(A,\theta)\) holds.
One may take the \(\theta,\rho,h\) supplied by the shifted-witness
lemma and
\[
 A=\frac{\max h}{\min h}\,
       \frac{\max\{2,\,2(12/5)^{24}(C_V/(1-q))^2\}}{1-\rho}.
\]
\end{theorem}

\begin{proof}
Lemma~\ref{l:Mtheta} supplies \(q^2<\theta<1\) and
\(\rho_0<\rho<1\), with the same fixed positive vector \(h\),
such that \(M_\theta h\le\rho h\).  Set
\(m_{XY,s}=M_{XY,s}\).  The envelope bounds and constant
extension are part of Definition~\ref{d:pairenv};
Lemma~\ref{l:recineq} supplies the recursion on all positive
integers with finite
\(C_T=2(12/5)^{24}(C_V/(1-q))^2\).
Lemma~\ref{l:renewal} therefore gives
\(F_X(d)\le A\theta^d\) for both pair types and all integers
\(d\ge0\), with exactly the stated finite \(A\ge2\).

To obtain the property in Definition~\ref{h:P}, fix its data
\(G,\tau,Y,u,v\).  Set all sites of \(Y\) to \(\mathrm X\)
and all other sites to \(\mathrm U\) with their full kept sets.
This is an ordinary consistent domain: \(Y\cap S=\varnothing\),
and every bond between nonremoved sites is kept by both ends.
It has precisely the partition sums of
Definition~\ref{d:pinmodel}, with the original legs and no new
pins at removed sites.  Designate the given common neighbour in
\(Y\) as centre.  A distance-two lattice pair is either collinear
or diagonal, so it lies in \(P_B\) or \(P_C\).
The bound for the corresponding envelope, at
\(d=\min\{d(u,S),d(v,S)\}\), is the required inequality.
If there are no such pairs, the assertion is vacuous.  In the
nonvacuous case \(G\ne\varnothing\) is finite.  A site with
maximal first coordinate has a neighbour outside \(G\), so
\(S\ne\varnothing\) and this minimum is a nonnegative integer.
Thus all quantifiers
of \(\mathsf P(A,\theta)\) are satisfied.
\end{proof}


\section{Exact computational inputs}

The computational inputs have three finite certificates. The local
input tests integer Hermitian-form matrices reconstructed from the
covered patterns. The open-trail input tests a positive integer
vector for a memory-six one-step matrix. The kernel input
tests five stationary lattice supersolutions and their exact
length-twelve prefixes. The certificates consist of the finite
integer equalities, inequalities, and divisibility tests specified
below and in Appendix C, with positive rational denominators cleared.
The accompanying formalization uses proofs checked by the Lean
kernel to verify these assertions. Each check whose kernel proof is
pending is marked \textbf{[certificate check pending]}. The soundness
arguments prove deterministic implications from these finite assertions.

For the local input we use the archived representation and inversion
choices, and the division-free Gram/residual format wherever its
witness is available. For a block without a supplied nonzero Gram
factor which is not directly diagonally dominant, we use the
Bareiss fallback permitted by Lemma~\ref{l:integercheck}. The local
certificate consists of the coverage, reconstruction, and positivity
checks on these data specified below.

\subsection*{The finite local verification}
The pattern data are \texttt{pext7962.json}, under
\texttt{programs/data/pext}, and
\texttt{certificates\_covered\_operator\_only.jsonl}, under
\texttt{programs/h1-local/out}, relative to the archive root.
The two supplied nonzero factors are
\texttt{h1\_chol\_\allowbreak5420\_\allowbreak131.json} and
\texttt{h1\_chol\_\allowbreak7697\_\allowbreak131.json} in \texttt{data}.
The finite assertions to be checked from these data are as follows.

First enumerate the patterns from Definition~\ref{d:pattern},
rather than assuming that the supplied pattern list is exhaustive.
Encode a kept set by a four-bit mask in direction order E,N,W,S.
Each of the three non-arrival neighbours has five choices: U
keeping all four bonds, X, or D keeping its passage bond and one
of its three other bonds. The arrival site has four choices: X,
or D keeping two of its three non-passage bonds. Each corner has
eight choices: U, X, or one of the six two-bond D kept sets.
The passage site's kept set and status, and all covering U tips,
are fixed. There are therefore
\(4\cdot5^3\cdot8^4=2048000\) assignments before consistency is
tested. Retain exactly those satisfying the nonremoved non-tip
bond consistency rule. Choose every nonremoved out-neighbour
among E,N,W. The finite enumeration assertions are
\[
 \begin{array}{c|r}
 \text{consistent status/kept-set assignments}&106101\\
 \text{such assignments with a marked out-neighbour}&184914\\
 \text{marked assignments with out-neighbour N or E}&123339.
 \end{array}
\]
\certpending{h1-pattern-enumeration}
In forming an effective record retain the present internal bonds
and the status of every non-tip site incident to such a bond.
Discard isolated factors as justified below. Identify records under
reflection \((x_1,x_2)\mapsto(-x_1,x_2)\), which fixes the arrival
site and interchanges the two turning directions. The finite
coverage test requires that their canonical keys are precisely
the \(7962\) distinct keys in the supplied pattern file. Every
entry in the operator-choice file must refer to the corresponding
pattern, use the covered model and tip radius \(3/5\), have the
correct out-neighbour, and use threshold \(961/1000\) for a
straight pattern and \(187/200\) for a turn. Each of the \(7962\)
pattern indices must occur exactly once in that choice file.
\certpending{h1-pattern-coverage}

For every record, reconstruct its chosen reduced or unreduced
representation and its vector groups using Lemma~\ref{l:leaf}.
The group's order is the ring-site order
\[
 (1,0),(1,1),(0,1),(-1,1),(-1,0),(-1,-1),(0,-1),(1,-1),
\]
omitting empty groups. A neighbour contributes its single bonded-tip
vector when that bond is present. A corner contributes its single
leaf vector in the reduced one-bond case, and its two slot vectors
otherwise. The inversion field lists zero-based vector indices;
check that it is a union of whole groups and that all its indices
are in range. Reconstruct every coefficient from the finite moments
of Lemma~\ref{l:moments}, the passage identity of
Lemma~\ref{l:Iw}, and these leaf replacements. The positive
denominator \(D_0\) is the product of the positive denominators in
this reconstruction, not an unverified matrix scale.
Construct the inverted Pauli operators of
\(3D_0\widehat Z_{\rm cur}\) and
\(5D_0\widehat Z_{\rm next}\), and their Hermitian forms
\(A_\epsilon,A'_\epsilon\) on every supported central exchange
sector, as in Lemma~\ref{l:integercheck}. In particular these are
forms in the unnormalized integer-coordinate sector bases.
At threshold \(c=p/q\), reconstruct
\[
 M_\epsilon=5qA_\epsilon-3pA'_\epsilon .
\]
A supplied matrix or a claimed list of sector sizes is not an
input to this reconstruction. Zero-dimensional sectors impose no
row tests. The finite construction yields \(12173\) supported
nonzero-dimensional blocks, all of size at most \(179\).
\certpending{h1-sector-reconstruction}

Of these blocks, \(8439\) satisfy the direct integer inequalities
\[
 (M_\epsilon)_{ii}>
       \sum_{j\ne i}|(M_\epsilon)_{ij}|\quad\text{for every }i.
\]
\certpending{h1-direct-dominance}
They use the Gram/residual certificate with \(L=0\) and
\(g=d=S=1\). The two archived nonzero factors match the regenerated
\(131\)-dimensional central sectors of their respective patterns.
Their supplied matrix arrays are compared entry by entry with the
regenerated matrix divided by their checked positive scalar
\(1600\); they are not trusted as matrices of the polynomial.
For the actual regenerated matrices, their certificate parameters
and minimum residual margins are
\[
 \begin{array}{c|r|rrr|r}
 \text{pattern}&\dim M&g&d&S&
       \min_i(E_{ii}-\sum_{j\ne i}|E_{ij}|)\\ \hline
 5420&131&1600&1&1&6820166897600\\
 7697&131&1600&1&1&27361669345600
 \end{array}
 \qquad E=dS^2M-gLL^{\mathsf T}.
\]
Check each supplied \(L\) to be an integer lower-triangular matrix
of the correct size and check every residual row, not just the
displayed minimum. These are the only two nonzero factors used.
\certpending{h1-gram-residual}

For each of the remaining \(3732\) blocks use the following exact
fallback. One may remove a positive integer common scalar
\(a_\epsilon\), checking every entry of
\(M_\epsilon=a_\epsilon B^{(0)}\) in the integers.
Set \(b_{-1}=1\). For \(0\le j<t=\dim M_\epsilon\), require
\(b_j=B^{(j)}_{jj}>0\), and, for every \(i,l>j\), require
\[
 \begin{gathered}
 b_{j-1}\mid
   \bigl(b_jB^{(j)}_{il}-B^{(j)}_{ij}B^{(j)}_{jl}\bigr),\\
 B^{(j+1)}_{il}
   =\frac{b_jB^{(j)}_{il}-B^{(j)}_{ij}B^{(j)}_{jl}}{b_{j-1}}.
 \end{gathered}
\]
Only the upper triangle need be stored: symmetry is preserved by
the recursion, so each lower-triangle numerator and quotient equals
its transposed upper-triangle entry.
It is enough to compute the upper triangle and reflect it, because
these matrices are symmetric. At the last stage there are no
trailing entries to update. These requirements are finite integer
divisibility and strict positivity assertions over matrices
regenerated from the archived choices. No stored pivot or floating
eigenvalue replaces any requirement. As an additional finite check,
all \(3734\) blocks failing the direct dominance test, including
the two with supplied nonzero factors, pass this fallback too.
\certpending{h1-bareiss}

\begin{theorem}\label{t:H1}
For every \(\pi\in\mathcal P_{\rm ext}\) and every
\(m\in(\mathbb S^2)^{20}\),
\[
 Z_{\rm cur}(\pi;m)\ge
 \begin{cases}
 (961/1000)Z_{\rm next}(\pi;m),&\pi\text{ straight},\\
 (187/200)Z_{\rm next}(\pi;m),&\pi\text{ turning}.
 \end{cases}
\]
\end{theorem}

\begin{proof}
We first justify the finite reduction and the implication of its
arithmetic tests. The finite assertions specified above constitute
the local certificate, verified by the Lean kernel in the accompanying
formalization, with pending checks marked at their specifications.

The pattern enumeration is exhaustive because every free status
and kept-set choice in Definition~\ref{d:pattern} appears in its
indicated finite option list. Each complete assignment appears
exactly once in their Cartesian product. The consistency predicate
is precisely the requirement that both nonremoved non-tip ends
make the same keep decision. A permitted out-neighbour is exactly
a nonremoved E,N or W neighbour: it keeps its passage bond, and
the passage site keeps that bond, so the out-bond is present.
These observations prove that filtering and marking produce
exactly all raw patterns, without a geometric or boundary assumption.

A nonremoved isolated tip or corner contributes a positive factor
depending only on its own slots, by Lemma~\ref{l:leaf}. It is
the same factor in the current and next partition polynomials,
because the only change is at the passage site and its out-bond.
An isolated arrival site has no slot factors, and its probability
functional applied to \(1\) is \(1\). The other nonremoved
non-arrival neighbours are incident to their passage bonds and
are not isolated. Removed sites contribute no factor. Thus all
discarded factors are common and positive. Every remaining factor
is determined by the present internal bonds, the remaining statuses
and the prescribed external slots. In particular, two kept sets
which give the same effective record cannot differ in a remaining
factor: both corner slots are included independently of whether
their original external bonds were kept.

The reflection used for effective keys only renames template sites
and external slot indices. Relabeling the finite product variables
is a bijection preserving their U and D laws and every scalar
product factor. It also preserves straight versus turning type.
Consequently the checked effective record covers its entire
reflection class, for arbitrary original unit slots, with the
common positive factors just described. This proves the coverage
implication of the finite key comparison, rather than inferring it
from a list of pattern IDs.

All moment contractions used in the reconstruction have separate
degree at most four. A bonded tip has its three slots and its one
bond. A corner has its two slots and at most two internal bonds.
After passage integration, a neighbouring site has at most its
three remaining factors and one inserted coordinate from a passage
pair; the two members of such a pair are distinct. The arrival
site has at most two present bonds and no passage insertion.
Leaf replacements are exactly those of Lemma~\ref{l:leaf}.
Thus Lemmas~\ref{l:moments} and \ref{l:Iw} justify the rational
coefficient tensors, and Lemma~\ref{l:integercheck} identifies each
reconstructed \(M_\epsilon\) with \(15D_0q\) times its inverted
Pauli Hermitian-form matrix. Every scale here is positive.

For a directly dominant block the first set of inequalities is
Lemma~\ref{l:gram-residual} with \(L=0\). For either supplied
nonzero factor, the checked residual inequalities and positive
scales give the same conclusion by that lemma. In both cases
\(M_\epsilon\) is positive definite.

We justify also the optional removal of a common scalar in a
fallback block. Put \(a=a_\epsilon>0\), and compare the displayed
Bareiss recursion on \(B^{(0)}=M_\epsilon/a\) with the recursion
of Lemma~\ref{l:integercheck} on \(A^{(0)}=M_\epsilon\).
Induction on \(j\) gives, on the trailing entries,
\[
 A^{(j)}=a^{j+1}B^{(j)},\qquad
 p_j=a^{j+1}b_j>0.
\]
For the induction step the numerator scales by \(a^{2j+2}\)
and the preceding pivot scales by \(a^j\); the quotient is
\(a^{j+2}B^{(j+1)}_{il}\). At \(j=0\) the preceding pivot
is \(p_{-1}=b_{-1}=1\), so this calculation includes the initial
step. Every normalized division is exact by its checked
divisibility assertion, and therefore every unnormalized division
is also exact and gives the displayed integer quotient. All
unnormalized pivots are positive. Lemma~\ref{l:integercheck}
therefore proves positivity of the actual \(M_\epsilon\).
The trailing matrices remain symmetric: the two products in each
numerator are unchanged when \(i,l\) are interchanged. Hence a
checker may store and test their upper triangles and recover the
lower triangles by symmetry, without omitting a distinct division.
This argument requires neither a numerical factorization nor a
claim that a supplied matrix has the required origin.

For each pattern, all supported central exchange-sector forms are
now positive definite. By the positive scale in
Lemma~\ref{l:integercheck}, the same holds for the inverted Pauli
operator on each central-weight space. Lemma~\ref{l:Ostruct}
gives its Hermitian symmetry, the required commutators and the
exchange-sector decomposition. Lemma~\ref{l:sector} extends
positivity from each central-weight space to its full supported
sector. The operator vanishes on discarded antisymmetric U-pair
subspaces, by Lemma~\ref{l:Ostruct}; those zero subspaces do not
require strict positivity. Hence the full inverted operator is
positive semidefinite.

Lemma~\ref{l:inversion} yields
\(\widehat F_c\ge0\) on the product of closed unit balls, for
the chosen groupwise inversion. Corollary~\ref{c:reduce} transfers
this to \(F_c=Z_{\rm cur}-cZ_{\rm next}\ge0\) on the original
twenty unit slots. Equivalently, on an assignment to the raw
pattern, the common positive reduction factor multiplies a
nonnegative reduced value. The preceding coverage argument proves
this for every raw pattern and every unit slot assignment.
Choosing its correct straight or turning threshold proves the
two displayed inequalities.
\end{proof}

\begin{theorem}\label{t:inputs}
At \(c_{\rm s}=961/1000\), \(c_{\rm t}=187/200\), (H1) holds
by the operator-only certificates of Appendix D, using the supplied
Gram/residual checks and its Bareiss fallback.
Property (H2) holds with
\[
 q=\frac{1119948287353}{1120437614225}<\frac{2499}{2500}<1,
 \qquad C_V=10 .
\]
For the type-dependent coefficients of Lemma~\ref{l:log-mean-bounds},
the orderings and vector \(h\) of Definition~\ref{d:kernel-order},
and \(\rho_0=993015/10^6\), property (H3\(_h\)) holds.
The latter assertion is supplied by five stationary lattice Bellman
certificates with memory \(m=10\), exact prefix \(L=12\ge m+2\),
radius \(R=70\), value scale \(2^{16}\), and coefficient scale
\(2^{48}\), using the checked 3741-state automaton quotient.
\end{theorem}

\begin{proof}
Theorem~\ref{t:H1} is precisely Definition~\ref{h:H1} at the two
thresholds, with its pointwise quantifiers. The direct Gram checks,
the two supplied nonzero factors and the fallback cover every
supported sector as established above.
Proposition~\ref{p:H2-test} gives a positive vector with
\(\mathsf A_6v\le q_*v\) and \(0<q_*<2499/2500<1\).
Lemma~\ref{l:alm} converts that finite inequality to the all-length
open-trail bound, and Lemma~\ref{l:H2-CV} permits \(C_V=10\).
Thus Definition~\ref{h:H2} holds uniformly in the start and
designated incoming direction, with initial reversal forbidden.

Proposition~\ref{p:automaton10} checks the labelled memory-ten
quotient. Proposition~\ref{p:stationary-check} gives the five
finite supersolution and exact-prefix assertions at
\(L=12\ge10+2\), \(R=70\) and the stated two scales.
Lemma~\ref{l:stationary-assembly} proves
\(M_1h\le(993015/10^6)h\), uniformly for every finite
true-trail cutoff. This is Definition~\ref{h:H3} for the stated
type-dependent coefficients and fixed orderings.
\end{proof}

\subsection*{Certification scope}
The mathematical reductions prove deterministic implications from
the specified finite integer assertions to the pointwise local
input and the two all-length trail inputs. The H1 calculation uses
the frozen representation and inversion records. Its residual
tests require all rows, and its fallback requires every division
exact and every pivot positive. The fallback discharges exactly
the blocks without supplied nonzero Gram factors, using the fixed
representation and inversion records.

Appendix C tests the true H2 revisit coefficient \(1/5\) and
the stated rate \(q_*\). Its H3 assertion concerns the kernel of
Definition~\ref{d:kernel}, with turning coefficient
\(a_{\rm t}=150000/10^6\), the exact dyadic vector and the
fixed orderings of Definition~\ref{d:kernel-order}. The stationary
certificates retain the initial bond and initial-site rules,
centre exclusion and first-hit target stopping at every suffix
time. No periodic spatial walk, long-time remainder or parity
improvement enters these input theorems.

The finite assertions require regenerated mathematical index sets
and coefficients. Supplied matrices, sector counts, numerical
eigenvalues, file hashes, and native exit codes do not replace
the prescribed comparisons. Each certificate is a proof of the
specified finite assertions checked by the Lean kernel in the
accompanying formalization. The placeholders identify the pending
kernel checks. None of these arguments assumes a physical window
estimate, a frame finite-energy bound, or a spectral-gap conclusion.


\appendix
\section{Finite probability and operator tools}
\label{a:finite}

Throughout this appendix, \(V\) is finite, each \(S_x\) is a
nonempty finite set, and \(\Omega=\prod_{x\in V}S_x\).
An empty product has its one empty configuration.  For a strictly
positive probability \(\mu\) on \(\Omega\), use
\(\langle f,h\rangle_\mu=\sum_\omega\mu(\omega)
\overline{f(\omega)}h(\omega)\).  If \(U\subset V\), write
\(\mu_{-U}\) for the marginal on \(V\setminus U\), and define
\[
 (E_Uf)(a,b)=\sum_{t\in\prod_{x\in U}S_x}
       \frac{\mu(t,b)}{\mu_{-U}(b)}f(t,b),\qquad
 \operatorname{Var}_U f=E_U|f-E_Uf|^2.
\]
All denominators here are positive.  We use the single-site Dirichlet
form and Poincare-gap convention of Definition~\ref{d:framemodel}.
More generally, a Poincare gap at least \(c\) for an operator \(H\)
means \(\langle f,Hf\rangle_\mu\ge c\operatorname{Var}_\mu f\)
for every \(f\).  This convention includes one-point spaces.

\begin{lemma}\label{l:CE}
For a positive probability \(\mu\) on a finite product
\(\Omega=\prod_{x\in V}S_x\), with nonempty finite alphabets,
\(E_U\) is a self-adjoint orthogonal projection preserving constants
and \(\mu\). For real or complex \(f\),
\[
 \langle f,(I-E_U)f\rangle_\mu
   =\mu(\operatorname{Var}_U f)=\|f-E_Uf\|_\mu^2,\qquad
 \operatorname{Var}_\mu f=\mu(\operatorname{Var}_U f)
                          +\operatorname{Var}_\mu(E_Uf).
\]
A finite sum \(\sum_a(I-E_{B_a})\), allowing repeated blocks,
has quadratic form \(\sum_a\mu(\operatorname{Var}_{B_a}f)\).
\end{lemma}

\begin{proof}
Fix a complement configuration \(b\) and put
\(p_b(t)=\mu(t,b)/\mu_{-U}(b)\).
On this fibre, \(E_Uf\) is the constant \(p_b(f)\).  Consequently
\(E_U^2=E_U\), \(E_U1=1\), and grouping the finite sum by \(b\)
gives
\[
 \langle f,E_Uh\rangle_\mu
   =\sum_b\mu_{-U}(b)\overline{p_b(f)}p_b(h)
   =\langle E_Uf,h\rangle_\mu,
 \qquad \mu(E_Uf)=\mu(f).
\]
Thus \(E_U\) is the orthogonal projection onto the functions of
the complement of \(U\).  Orthogonality of \(f-E_Uf\) and \(E_Uf\)
proves
\[
 \langle f,(I-E_U)f\rangle_\mu
   =\|f-E_Uf\|_\mu^2
   =\sum_b\mu_{-U}(b)p_b(|f-p_b(f)|^2).
\]
This is the first asserted identity.  Since \(\mu(E_Uf)=\mu(f)\),
the decomposition
\(f-\mu(f)=(f-E_Uf)+(E_Uf-\mu(f))\) is also orthogonal.
Its squared norm is the total-variance identity.  Summing the first
identity over block labels gives the last assertion, including repeated
blocks.  The argument uses complex conjugates as displayed and therefore
applies to complex functions.  For \(U=\varnothing\) it gives \(E_U=I\);
for \(U=V\) it gives the projection onto constants, also when \(V\)
is empty.
\end{proof}

\begin{lemma}\label{l:finitecoupling}
For probabilities \(p,p'\) on a finite set,
\[
 d_{\rm TV}(p,p')=\tfrac12\sum_z|p(z)-p'(z)|
          =1-\sum_z\min\{p(z),p'(z)\}.
\]
There exists a coupling with disagreement probability equal to this
quantity; every coupling has at least that disagreement probability.
Marginalization cannot increase total variation.
For laws on configurations in a finite \(V\), a marginal coupling
on \(Q\subset V\) can be completed to a coupling of the full laws.
If its disagreement probability on \(Q\) is at most \(\varepsilon\),
the completion can have expected Hamming cost at most
\(|V\setminus Q|+|V|\varepsilon\).
If both laws dominate \(p_0\nu\) for a common probability \(\nu\)
and \(0\le p_0\le1\), their total variation is at most \(1-p_0\).
\end{lemma}

\begin{proof}
For each state \(z\) put \(a(z)=\min\{p(z),p'(z)\}\) and
\(t=\sum_z a(z)\).  The equality
\(|p(z)-p'(z)|=p(z)+p'(z)-2a(z)\), summed over \(z\), proves
the formula for total variation.  If \(t<1\), define a joint law by
\[
 \gamma(z,z')=a(z)\mathbf1_{\{z=z'\}}
       +\frac{(p(z)-a(z))(p'(z')-a(z'))}{1-t}.
\]
Its two marginals are \(p,p'\).  At each \(z\), at least one of
\(p(z)-a(z)\) and \(p'(z)-a(z)\) is zero, so the second term
has no diagonal mass.  Hence its disagreement probability is \(1-t\).
This formula includes \(t=0\); if \(t=1\), use only the diagonal
law \(\gamma(z,z)=p(z)=p'(z)\).  In any coupling, the mass at
\((z,z)\) is at most \(a(z)\), proving optimality.

For a map \(T\) to a finite set, summing over its fibres gives
\[
 \sum_b\left|\sum_{z:Tz=b}(p(z)-p'(z))\right|
       \le\sum_z|p(z)-p'(z)|.
\]
Taking \(T\) to be a coordinate restriction proves the marginalization
claim.

Let \(\gamma_Q\) be any coupling of the two \(Q\) marginals.
After sampling \((a,a')\) from \(\gamma_Q\), sample the remaining
coordinates independently with laws
\(p(\,\cdot\mid n_Q=a)\) and \(p'(\,\cdot\mid n'_Q=a')\).
A zero-probability conditioning value has zero probability of being
sampled; its conditional law may be chosen arbitrarily.  Summing out
the other component verifies both full marginals.  Pointwise in this
completion,
\[
 d_H(n,n')\le |V\setminus Q|
                      +|V|\mathbf1_{\{n_Q\ne n'_Q\}}.
\]
Taking expectations proves the requested bound.  Empty coordinate sets
and one-point configuration spaces are included in this construction.
Finally, domination by \(p_0\nu\) implies
\[
 \sum_z\min\{p(z),p'(z)\}\ge p_0\sum_z\nu(z)=p_0.
\]
The total-variation formula gives the last assertion, also for
\(p_0=0\) or \(p_0=1\).
\end{proof}

\begin{lemma}\label{l:surgery}
Let \(\nu\) couple two positive finite-volume specification laws.
Fix a set \(F\) in the volume. For every input pair, resample \(F\)
from any coupling whose first and second marginals are the two
respective conditional laws given their own complements, and keep
the complements fixed. The resulting law still couples the original
two laws. The chosen coupling may depend on the entire input pair.
A finite composition of such kernels has the same property.
\end{lemma}

\begin{proof}
Write the two original laws as \(\mu,\mu'\), and an input pair as
\((\omega,\omega')\).  For any test function \(f\) of the first
output alone, its conditional expectation after the operation is
\((E_F^\mu f)(\omega)\): by the stipulated marginal of the chosen
coupling this value does not depend on \(\omega'\), even when the
joint coupling does.  Integrating against \(\nu\) therefore gives
\[
 \sum_{\omega,\omega'}\nu(\omega,\omega')
                  (E_F^\mu f)(\omega)
       =\mu(E_F^\mu f)=\mu(f),
\]
where the last equality is Lemma~\ref{l:CE}.  The same calculation
with a function of the second output gives its law \(\mu'\).
Consistency identifies these conditional laws with the specification
laws on \(F\) with their respective complements fixed.  Applying the
calculation at each stage proves the assertion about finite compositions.
If \(F\) is empty, its unique resampling is the identity, so this case
is included.
\end{proof}

\begin{lemma}\label{l:density}
Let \(\pi\) be a positive product probability on a finite product,
and \(d\mu=\varrho\,d\pi\), with \(0<m\le\varrho\le M<\infty\).
Then for every \(f\),
\[
 \mathcal D_\mu(f)\ge m\mathcal D_\pi(f),\qquad
 \operatorname{Var}_\mu f\le M\operatorname{Var}_\pi f,\qquad
 \mathcal D_\pi(f)\ge\operatorname{Var}_\pi f .
\]
Thus \(\lambda(\mu)\ge m/M\).
In a conditional frame block, absorbing site weights and every
fixed boundary factor into \(\pi\) gives
\(\lambda(\mu)\ge r_b^{|E(Q)|}\).
\end{lemma}

\begin{proof}
For any positive probability \(\nu\), Lemma~\ref{l:CE} gives the
following variational formulas, with the infima taken over complex
functions or constants when \(f\) is complex:
\[
 \nu(\operatorname{Var}_x f)
   =\min_{h:\,h\text{ depends only on }V\setminus\{x\}}
                      \nu(|f-h|^2),\qquad
 \operatorname{Var}_\nu f=\min_{z\in\mathbb C}\nu(|f-z|^2).
\]
Indeed the two minimizers are \(E_x^\nu f\) and \(\nu(f)\),
respectively; subtracting either minimizer splits the squared norm
orthogonally.  For every admissible \(h\), the lower density bound
implies \(\mu(|f-h|^2)\ge m\pi(|f-h|^2)\).
Taking the first minimum, and then summing over sites, gives
\(\mathcal D_\mu(f)\ge m\mathcal D_\pi(f)\).
Using the constant \(z=\pi(f)\) in the second minimum gives
\[
 \operatorname{Var}_\mu f\le\mu(|f-\pi(f)|^2)
                         \le M\operatorname{Var}_\pi f.
\]

For completeness, the product-law inequality follows from finite
orthogonal projections.  Enumerate \(V=\{x_1,\ldots,x_v\}\) and let
\(R_j=E_{x_j}^\pi\).  These projections commute because \(\pi\) is
a product probability.  Put \(f_0=f\) and
\(f_j=R_1\cdots R_jf\).  Then \(f_v=\pi(f)\).
The differences \(f_{j-1}-f_j\) are mutually orthogonal: for \(j<k\)
the former is killed by \(R_j\), whereas the latter is fixed by
\(R_j\).  Also
\[
 \|f_{j-1}-f_j\|_\pi
   =\|R_1\cdots R_{j-1}(I-R_j)f\|_\pi
   \le\|(I-R_j)f\|_\pi.
\]
Each projection has norm at most one by Lemma~\ref{l:CE}.
Adding the squared bounds proves
\(\operatorname{Var}_\pi f\le\mathcal D_\pi(f)\).
For \(v=0\) both sides vanish.  Combining all three inequalities
yields
\(\mathcal D_\mu(f)\ge(m/M)\operatorname{Var}_\mu f\).

We spell out the frame application.  By
Definitions~\ref{d:framemodel} and~\ref{d:frame}, and
Lemma~\ref{l:threeframes}, each conditional block law has the form
\[
 \mu(n)=\frac1{Z_Q}
       \prod_{x\in Q}h_x(n_x)
       \prod_{e=\{x,z\}\in E(Q)}(1+\varepsilon_e n_x\cdot n_z),
 \qquad h_x(t)>0.
\]
Here \(h_x\) is the site weight times all fixed boundary factors
incident to \(x\); the signs are \(+1\) for the plane law and
\(-1\) for the antiferromagnetic law.  Define
\(\pi_x(t)=h_x(t)/\sum_{a\in S_x}h_x(a)\) and
\(\pi=\bigotimes_{x\in Q}\pi_x\).
Proper colouring ensures that both endpoints of every internal or
boundary bond use distinct frames.  Thus all boundary factors are
positive, and every internal factor is between \(b\) and \(2-b\).
With \(e_Q=|E(Q)|\), set
\[
 \psi(n)=\prod_{e=\{x,z\}\in E(Q)}
                     (1+\varepsilon_e n_x\cdot n_z).
\]
Then \(d\mu/d\pi=\psi/\pi(\psi)\), and
\[
 \frac{b^{e_Q}}{\pi(\psi)}
      \le\frac{d\mu}{d\pi}
      \le\frac{(2-b)^{e_Q}}{\pi(\psi)}.
\]
The ratio of these bounds is \(r_b^{e_Q}\), proving the frame claim.
For an empty block the empty product is one and the Poincare inequality
is automatic.  Only internal bonds have entered this density ratio.
\end{proof}

\begin{lemma}\label{l:comparison}
Let \(\mu>0\) be on a finite product, and let \((B_a)\) be a
nonempty finite labelled family of nonempty blocks.
Suppose the single-site heat bath of every conditional law on each
\(B_a\) has Poincare gap at least \(g>0\), uniformly in its exterior fibre,
and \(\sum_a(I-E_{B_a})\) has Poincare gap at least \(c>0\).
If \(t_{\max}=\max_x\#\{a:x\in B_a\}\), then
\[
 \operatorname{Var}_\mu f\le
             \frac{t_{\max}}{cg}\mathcal D_\mu(f).
\]
Repeated blocks are counted with multiplicity.
\end{lemma}

\begin{proof}
Fix a block label \(a\) and a complement configuration \(b\).
Apply its assumed Poincare inequality to
\(f_b(t)=f(t,b)\) in the conditional probability
\(\mu(\,\cdot\mid n_{V\setminus B_a}=b)\).
If \(x\in B_a\), further conditioning on \(B_a\setminus\{x\}\)
fixes precisely all coordinates other than \(x\).  The one-site
conditional variance in this fibre is therefore the restriction of
\(\operatorname{Var}_x f\) for the original law.  Consequently,
\[
 \operatorname{Var}_{B_a}f
      \le\frac1g\sum_{x\in B_a}E_{B_a}
                                      (\operatorname{Var}_x f).
\]
Integrate and sum over labels.  Lemma~\ref{l:CE} and the finiteness
of all sums give
\[
 \begin{split}
 c\operatorname{Var}_\mu f
 &\le\sum_a\mu(\operatorname{Var}_{B_a}f)\\
 &\le\frac1g\sum_x\#\{a:x\in B_a\}
                         \mu(\operatorname{Var}_x f)
 \le\frac{t_{\max}}g\mathcal D_\mu(f).
 \end{split}
\]
Division by \(c>0\) proves the assertion.  Every sum is over labels,
so repetitions contribute their full multiplicity.  Conditional laws
on one-point spaces cause no exception because their variances vanish.
\end{proof}

\begin{lemma}\label{l:contraction}
Let \(P=A^{-1}\sum_aE_{B_a}\), where \(A\ge1\) is the number
of labels. Suppose \(0<c\le A\) and that for every two input
configurations at Hamming distance one, their \(P\)-output laws
have a coupling of expected Hamming distance at most \(1-c/A\).
Then \(\sum_a(I-E_{B_a})=A(I-P)\) has Poincare gap at least \(c\).
This applies to site-dependent finite alphabets and also to a
one-point product space.
\end{lemma}

\begin{proof}
If \(\Omega\) has one point, every variance is zero and the conclusion
holds.  Otherwise, for a real function \(f\) define
\[
 \operatorname{Lip}(f)=
   \max_{d_H(\omega,\omega')=1}|f(\omega)-f(\omega')|.
\]
Any pair can be joined by changing its disagreeing coordinates one at
a time.  This path has length \(d_H(\omega,\omega')\), so
\[
 |f(\omega)-f(\omega')|
            \le\operatorname{Lip}(f)d_H(\omega,\omega').
\]
In particular a nonconstant function has positive Lipschitz seminorm,
even if some site alphabets are singletons.  For an input pair at
distance one, apply the assumed coupling to this inequality and take
expectations.  With \(r=1-c/A\in[0,1)\), it gives
\[
 |Pf(\omega)-Pf(\omega')|\le r\operatorname{Lip}(f),
 \qquad \operatorname{Lip}(Pf)\le r\operatorname{Lip}(f).
\]

Lemma~\ref{l:CE} shows that \(P\) is real, self-adjoint, preserves
constants and \(\mu\), and is a positive contraction on \(L^2(\mu)\).
Indeed it is an average of orthogonal projections.  The real subspace
\(\{f:\mu(f)=0\}\) is invariant.  On it choose a real orthonormal
eigenbasis: conjugation by
\(\operatorname{diag}(\sqrt{\mu(\omega)})\) makes the matrix real
symmetric, so finite-dimensional orthogonal diagonalization applies.
If \(Pf=\lambda f\) is a nonzero vector of that subspace, it is
nonconstant and hence
\[
 |\lambda|\operatorname{Lip}(f)
       =\operatorname{Lip}(Pf)\le r\operatorname{Lip}(f).
\]
It follows that \(A(1-\lambda)\ge A(1-r)=c\).
Expanding \(f-\mu(f)\) in this eigenbasis proves
\[
 \langle f,A(I-P)f\rangle_\mu\ge c\operatorname{Var}_\mu f.
\]
Constants are annihilated by \(A(I-P)\), and applying the real
inequality to the real and imaginary parts proves the complex version.
The argument also covers \(c=A\), when \(r=0\).
\end{proof}


\section{Star-connected boundaries and connected thickening}
\label{s:topology}

Two sites are nearest neighbours when their ambient \(\ell^1\)
distance is one, and star-neighbours when their ambient
\(\ell^\infty\) distance is one.  A path in either graph is
a finite sequence of vertices with consecutive vertices adjacent
in that graph.  A one-point set is connected in both senses.

\begin{lemma}\label{l:boundaryconn}
If \(A\subset\mathbb Z^2\) is finite, nonempty and star-connected,
and \(A^c\) is star-connected, then
\(\partial_\infty A=\{z\notin A:\operatorname{dist}_\infty(z,A)=1\}\)
is nonempty and star-connected.
Here star-neighbours differ by \(\ell^\infty\) distance one.
\end{lemma}

\begin{proof}
Let \(\mathcal K\) be the graph whose vertices are \(\mathbb Z^2\)
and whose edges join star-neighbours.  Finite edge sets form
a vector space over \(\mathbb F_2\), with addition given by
symmetric difference.  An edge set is called even if each
vertex has even degree in it.  A triangle consists of three
pairwise star-neighbouring vertices and its boundary is its
three-edge set.  Such a triangle lies in a square
\((i,j)+\{0,1\}^2\).

We first prove that every finite even edge set is a finite
sum of triangle boundaries.  Let \(Z\) be such an edge set.
For each diagonal edge \(e=\{p,p+(1,\epsilon)\}\) of \(Z\),
where \(\epsilon\in\{-1,1\}\), add the boundary of the triangle
\(\{p,p+(1,0),p+(1,\epsilon)\}\).  Its only diagonal edge
is \(e\).  The resulting edge set \(Z'\) is finite, even,
and has horizontal and vertical edges only.

Write \(H(i,j)\in\mathbb F_2\) for the indicator that
\(\{(i,j),(i+1,j)\}\) belongs to \(Z'\), and
\(V(i,j)\) for the analogous indicator of
\(\{(i,j),(i,j+1)\}\).  Choose an integer \(R\ge1\) such that
every endpoint of \(Z'\) lies in \([-R,R]^2\); when \(Z'\) is
empty take \(R=1\).  In particular \(H(i,r)=0\) for
\(|r|>R\), and \(V(i,R)=0\).  Define the finite sum
\[
 n(i,j)=\sum_{\substack{-R\le r\le R\\r>j}}H(i,r)
                    \quad\hbox{in }\mathbb F_2.
\]
For each fixed \(i\), \(\sum_{r=-R}^R H(i,r)=0\).  Indeed this sum
is the parity of the number of edges crossing the line from
the half-plane \(x_1\le i\) to its complement.  It equals
the sum of the degrees of \(Z'\) at vertices of that
half-plane, because edges inside it contribute twice.
This sum is zero because \(Z'\) is even.  All degree sums
are finite after omitting the zero terms, since \(Z'\) is
finite.  It follows that \(n(i,j)\) is nonzero for only
finitely many pairs \((i,j)\): only finitely many columns
contain horizontal edges, and in each such column the sum
vanishes both above its highest horizontal edge and below
its lowest one.

Let \(P_{ij}=(i,j)+\{0,1\}^2\), and let \(\partial P_{ij}\)
be its four nearest-neighbour edges.  We claim that
\[
                    Z'=\sum_{i,j}n(i,j)\partial P_{ij}.
\]
The horizontal edge with indicator \(H(i,j)\) occurs on
the right with coefficient
\(n(i,j-1)+n(i,j)=H(i,j)\).
The vertical edge with indicator \(V(i,j)\) occurs with
coefficient
\[
 n(i-1,j)+n(i,j)
       =\sum_{\substack{-R\le r\le R\\r>j}}
                          \bigl(H(i-1,r)+H(i,r)\bigr).
\]
For \(j<R\), sum the even-degree identity
\(H(i-1,r)+H(i,r)+V(i,r-1)+V(i,r)=0\)
over the finite integer interval \(j+1\le r\le R\).
Terms with \(r<-R\) are zero by the endpoint bound.
All vertical terms cancel in pairs except \(V(i,j)\) and
\(V(i,R)=0\), proving that this coefficient is \(V(i,j)\).
For \(j\ge R\), both the sum and \(V(i,j)\) are zero.
The claim follows.  Finally, each \(\partial P_{ij}\) is
the sum of the two triangle boundaries obtained by inserting
one of the diagonals of \(P_{ij}\).  Together with the
initial diagonal eliminations, this proves the assertion
about even edge sets.

Let \(\mathcal B\) be the set of edges of \(\mathcal K\)
with one endpoint in \(A\) and one in \(A^c\).  It is finite,
and it is nonempty because a point of \(A\) of largest first
coordinate has its right nearest neighbour outside \(A\).
Say that two edges of \(\mathcal B\) are linked if some
triangle boundary contains both.  We prove that any two
edges of \(\mathcal B\) can be joined by a finite sequence
of linked edges.  Otherwise choose a connected component
\(\mathcal B_1\) for this relation and put
\(\mathcal B_2=\mathcal B\setminus\mathcal B_1\).
Both parts are nonempty, and no linked pair crosses between
them.  For finite edge sets define the linear functional
\[
              \varphi(Z)=|Z\cap\mathcal B_1|\pmod2.
\]
A triangle has either zero or two edges in \(\mathcal B\),
according to whether all its vertices lie on the same side
of \(A\), or they do not.  In the latter case the two edges
are linked and belong to the same \(\mathcal B_r\).
Thus \(\varphi\) vanishes on every triangle boundary, and
the preceding decomposition implies that it vanishes on
every finite even edge set.

Choose \(e_r=\{a_r,b_r\}\in\mathcal B_r\), with
\(a_r\in A\) and \(b_r\in A^c\), for \(r=1,2\).
There is a finite star-path from \(a_1\) to \(a_2\) in
\(A\), and a finite star-path from \(b_2\) to \(b_1\)
in \(A^c\), by the two connectivity assumptions.
Concatenate these paths with \(e_2\) and \(e_1\) to form
a closed walk.  The set \(Z\) of edges traversed an odd
number of times by this walk is finite and even: at every
vertex the number of entrances equals the number of exits
modulo two.  Its intersection with \(\mathcal B\) consists
of exactly the distinct edges \(e_1,e_2\), since both
connecting paths stay on one side.  This gives
\(\varphi(Z)=1\), a contradiction.  The linked-edge relation
is therefore connected.

Every edge of \(\mathcal B\) has a unique endpoint in
\(A^c\), and these outer endpoints are precisely the
vertices of \(\partial_\infty A\).  Outer endpoints of
two linked edges lie in the same triangle, so they are
equal or star-neighbours.  A sequence of linked edges
thus gives a star-path between their outer endpoints,
after deleting consecutive repetitions.  Surjectivity of
the outer-endpoint map proves star-connectivity of
\(\partial_\infty A\); nonemptiness follows from that of
\(\mathcal B\).
\end{proof}

\begin{lemma}\label{l:thicken}
If \(X_0\) is finite, nonempty and star-connected, then
\(X=\{z:d(z,X_0)\le1\}\) is nearest-neighbour connected.
For any finite \(G\) containing \(X\), every \(0<\vartheta<1\)
satisfies
\[
 \sum_{z\in X}\vartheta^{d(z,G^c)}
       \le5\vartheta^{-1}
                      \sum_{x\in X_0}\vartheta^{d(x,G^c)}.
\]
\end{lemma}

\begin{proof}
For each \(x\) put
\(C(x)=\{x,x+(1,0),x-(1,0),x+(0,1),x-(0,1)\}\).
This five-point cross is nearest-neighbour connected,
and \(X=\bigcup_{x\in X_0}C(x)\).
For star-neighbours \(x,x'\), if their difference has
one zero coordinate then they are nearest neighbours
and \(C(x)\cap C(x')\) contains both centres.
Otherwise \(x'-x=(\epsilon_1,\epsilon_2)\), where
\(\epsilon_1,\epsilon_2\in\{-1,1\}\), and their crosses
share the two points \(x+(\epsilon_1,0)\) and
\(x+(0,\epsilon_2)\).  These are all possibilities for
star-neighbours.  A star-path between centres therefore
gives a chain of intersecting connected crosses.  Any
two points of their union can be joined by concatenating
paths inside the crosses through these intersection
points.  This proves nearest-neighbour connectedness
of \(X\), including when \(X_0\) is a singleton.

For finite \(G\supset X\), its complement is nonempty
and the ambient distance \(b(z)=d(z,G^c)\) is
one-Lipschitz.  Thus for \(z\in C(x)\),
\(b(z)\ge b(x)-1\).  Since \(0<\vartheta<1\),
\(\vartheta^{b(z)}\le\vartheta^{-1}\vartheta^{b(x)}\).
Every point of \(X\) is in at least one of the crosses,
so summing nonnegative terms gives
\[
 \sum_{z\in X}\vartheta^{b(z)}
    \le\sum_{x\in X_0}\sum_{z\in C(x)}\vartheta^{b(z)}
    \le5\vartheta^{-1}\sum_{x\in X_0}\vartheta^{b(x)}.
\]
This proves the claimed estimate without requiring
the crosses to be disjoint.
\end{proof}


\section{Finite trail and stationary Bellman certificates}

Throughout this appendix the thresholds are
\(c_{\rm s}=961/1000\) and \(c_{\rm t}=187/200\), so that
\(x_{\rm s}=1000/961\) and \(x_{\rm t}=200/187\).
The trail weights, response pairs, orderings and vector \(h\) are
those of Definitions~\ref{d:weights}, \ref{d:kernel} and
\ref{d:kernel-order}. In particular, a revisit means an earlier
\emph{departure} from the same site. All positions below are points
of \(\mathbb Z^2\).

The finite assertions in Propositions~\ref{p:H2-test},
\ref{p:automaton10} and \ref{p:stationary-check} specify exact checks
of the archived witnesses. The arguments below prove the soundness
of those checks. The accompanying formalization verifies these finite
assertions by proofs checked by the Lean kernel; pending checks are
marked at their full arithmetic specifications. An executable's exit
code or a floating approximation is not an additional hypothesis of
any soundness lemma.

\begin{definition}\label{d:memory-automaton}
Let \(\mathcal E=\{{\rm E}=(1,0),{\rm N}=(0,1),
{\rm W}=(-1,0),{\rm S}=(0,-1)\}\).
For \(m\ge1\), \(\Omega_m\) consists of sequences
\(\omega=(d_1,\ldots,d_m)\in\mathcal E^m\) whose path from zero
uses distinct unoriented bonds. Write \(z_0=0\),
\(z_j=\sum_{i\le j}d_i\), and \(a(\omega)=d_m\).
A direction \(d\) is admissible if \(d\ne-d_m\) and
\(\{z_m,z_m+d\}\) is not a used bond of this window.
Its successor is \(\omega d=(d_2,\ldots,d_m,d)\).
Its kind \(k(\omega,d)\) is \({\rm r}\) if
\(z_m\in\{z_0,\ldots,z_{m-1}\}\), otherwise \({\rm s}\)
if \(d=d_m\), and \({\rm t}\) otherwise.
Direction labels remain absolute lattice directions.
\end{definition}

The successor belongs to \(\Omega_m\): the last \(m-1\) old bonds
and the new bond are a subset of the distinct bonds of the extended
window, and translation to a new initial point preserves bond equality.
The explicit nonreversal condition is redundant with the used-bond
test, since the last bond is used, but fixes the incoming-direction
convention.

\begin{definition}\label{d:memorygrowth}
On \(\Omega_m\), define the nonnegative one-step matrix
\[
 \mathsf A_m(\omega,\omega')=
 \sum_{\substack{d\ {\rm admissible}\\\omega d=\omega'}}
 \gamma(\omega,d),\qquad
 \gamma(\omega,d)=
 \begin{cases}1/5,&k={\rm r},\\
 x_{\rm s}/3,&k={\rm s},\\x_{\rm t}/3,&k={\rm t}.
 \end{cases}
\]
For a prescribed start and incoming direction with initial reversal
forbidden, let \(A_n(\omega)\) be the true total \(\mathsf W\)-weight
of length-\(n\) trails whose last \(m\) steps form \(\omega\),
for \(n\ge m\), and let \(W_n=\sum_{|w|=n}\mathsf W(w)\).
The adopted finite test is \(m=6\) and extension length \(1\).
\end{definition}

There is no virtual incoming bond in the used-bond set. Only the
initial reversal is forbidden. Thus the first departure has the same
straight or turning convention as every later first departure.

\begin{lemma}\label{l:alm}
For \(n\ge m\),
\[
 A_{n+1}(\omega')\le\sum_\omega A_n(\omega)\mathsf A_m(\omega,\omega').
\]
If \(v_\omega>0\), \(0<q<1\), and \(\mathsf A_mv\le qv\),
then (H2) holds with
\[
 C_V=\max\left\{
      \max_{0\le r<m}W_rq^{-r},\
      \frac{\max v}{\min v}W_mq^{-m}\right\}.
\]
Any larger finite constant is also valid.
\end{lemma}

\begin{proof}
Fix the initial data and a true trail \(w\) of length \(n+1\).
Its length-\(n\) prefix is a true trail with the same initial data.
If its last window is \(\omega\), its next direction is admissible
for \(\omega\): a used bond of this window is a used bond of the
true trail, and reversal would reuse the last bond. The last window
of the extension is exactly \(\omega d\). The map from \(w\) to its
prefix and final direction is injective.

Suppose first that the window detects a revisit. Its current endpoint
then equals an earlier window vertex from which a window bond departs.
This is also an earlier true departure. The new true and window
weights are both \(1/5\). If the true departure is a revisit not
detected by the window, its weight is \(1/5\), whereas the window
uses \(x_{\rm s}/3\) or \(x_{\rm t}/3\), each at least \(1/3\).
Otherwise both departures are first departures; their incoming
direction is the last window direction, so their type and weight
agree. These cases exhaust all departures. Multiplication by the
nonnegative prefix weight and summation over the injective map give
the asserted componentwise inequality.

Write \(A_n\) as a row vector and put
\(v_- =\min_\omega v_\omega>0\),
\(v_+=\max_\omega v_\omega\). Nonnegativity of the matrix implies,
by induction on \(j\),
\[
 A_{m+j}\le A_m\mathsf A_m^j,\qquad
 \mathsf A_m^jv\le q^jv .
\]
Consequently
\[
 W_{m+j}\le v_-^{-1}A_m\mathsf A_m^jv
       \le (v_+/v_-)q^jW_m .
\]
This proves the claimed bound for all lengths at least \(m\).
For lengths less than \(m\) it is the defining finite maximum.
The empty trail contributes \(W_0=1\), so \(C_V\ge1\).
Translations and the eight symmetries of the square give
weight-preserving bijections between the trail sets for any two
prescribed starts and incoming directions. Thus all the base sums
\(W_r\) in this formula are uniform over those initial data.
Increasing \(C_V\) preserves the inequality, which proves
\ref{h:H2} with the stated constants.
\end{proof}

\begin{lemma}\label{l:H2-dihedral}
The dihedral group of the square acts on \(\Omega_m\) by applying
the same lattice symmetry to every direction. This action preserves
admissibility, kind, successor and the weights of \(\mathsf A_m\);
thus \(\mathsf A_m(g\omega,g\omega')=\mathsf A_m(\omega,\omega')\).
For any positive vector constant on these orbits, a row inequality
\((\mathsf A_mv)_\omega\le qv_\omega\) checked at one representative
of every orbit holds on every raw window.
For \(m=6\), \(|\Omega_6|=916\) and there are \(115\) orbits.
\end{lemma}

\begin{proof}
A square symmetry is an invertible linear map carrying the four
unit directions to themselves. It carries vertices and unoriented
bonds bijectively to vertices and bonds. It therefore preserves
bond distinctness, occurrence of the endpoint among earlier
departure vertices, equality of the outgoing and incoming directions,
and the reversal relation. It also commutes with deleting the first
direction and adjoining the last. This proves every asserted
transition property, and hence the matrix identity, by the bijection
\(d\mapsto gd\) in its defining sum. If \(v\) is constant on orbits,
then
\[
 (\mathsf A_mv)_{g\omega}
 =\sum_{\omega'}\mathsf A_m(g\omega,g\omega')v_{g\omega'}
 =(\mathsf A_mv)_\omega ,\qquad v_{g\omega}=v_\omega .
\]
This proves the row reduction.

Here is a direct count for memory six. Let \(N_j=|\Omega_j|\),
with \(N_0=1\). There are \(N_1=4\) initial directions and
three nonreversing choices at each later step before any additional
bond obstruction. A nonreversing square-lattice path cannot reuse a
bond before step five: such reuse would require a return to a previous
vertex, and the shortest nonreversing cycle has length four.
Thus \(N_2=12,N_3=36,N_4=108\).
At step five the only additional obstruction is a four-step square
starting at the initial vertex, followed by its first edge. There
are four choices of first direction and two orientations of that
square. Exactly these eight extensions must be deleted, giving
\(N_5=3\cdot108-8=316\).

For a valid length-five prefix, an additionally forbidden sixth
step must start at its vertex of time one. Indeed, a repeat of an
earlier vertex at time five has an even time difference at least
four, and only time one qualifies; a difference two would already
reverse a bond. Hence steps two through five form a square based at
that vertex. There are initially \(4\cdot3\cdot2=24\) choices
of the first direction, the square's first direction and its
orientation. Of these, exactly eight have the square's last edge
equal to the first bond of the entire path: for each first
direction, the square's first direction must be one of its two
perpendicular directions, and precisely one orientation then ends
along that first bond. These eight are invalid length-five
prefixes. Each of the remaining sixteen prefixes has precisely two
used edges at the current vertex besides the reverse of the most
recent edge: the initial bond and the first edge of the square.
Both give distinct forbidden sixth directions. This gives
\(N_6=3\cdot316-2\cdot16=916\).

For completeness, count the orbits without a supplied representative
list. The number of orbits is
\(\frac18\sum_g|\{\omega:g\omega=\omega\}|\): to see this,
sum the number of stabilizers over each orbit and use that the orbit
size times its stabilizer size is eight. A nonidentity rotation, or
a reflection in a diagonal, fixes no unit direction and hence no
window. Reflection in either coordinate axis fixes two directions.
A window using just those two opposite directions must use a single
one throughout, since a change of sign immediately reverses a bond.
Each such reflection therefore fixes exactly two windows. The orbit
count is \((916+2+2)/8=115\).
\end{proof}

\begin{proposition}\label{p:H2-test}
At the two fixed H1 thresholds there exists a positive integer
vector \(v\) on \(\Omega_6\), constant on dihedral orbits, with
\[
 \frac{\max v}{\min v}=\frac{250000000000}{46594269069}<6,\qquad
 \max_\omega\frac{(\mathsf A_6v)_\omega}{v_\omega}
       =q_*:=\frac{1119948287353}{1120437614225}<\frac{2499}{2500}.
\]
The weights have common denominator \(D_2=2695605\), with
numerators \(935000,961000,539121\) for \({\rm s},{\rm t},{\rm r}\).
For \(q_*=q_{\rm num}/q_{\rm den}\), the required row test is
\[
 q_{\rm den}\sum_{d\ {\rm admissible}}n_{k(\omega,d)}v_{\omega d}
          \le q_{\rm num}D_2v_\omega .
\]
It suffices to check the \(115\) orbit rows after proving
Lemma~\ref{l:H2-dihedral}; checking all \(916\) raw rows is also valid.
\end{proposition}

\begin{proof}[Finite verification contract]
Number the directions E,N,W,S by \(0,1,2,3\), and for this
certificate use the least-significant-first code
\[
 e_6(d_1,\ldots,d_6)=\sum_{i=1}^6d_i4^{i-1}.
\]
The canonical code is the minimum of \(e_6\) over the eight images
\(d_i\mapsto \varepsilon d_i+r\pmod4\),
\(\varepsilon\in\{1,-1\}\), \(r\in\{0,1,2,3\}\).
Use the integer field \texttt{h} indexed by the field \texttt{code}
in \texttt{h2\_exact\_m6k1\_r5.json}; pull it back by this
canonical-code map to obtain \(v\). The decimal vector files are
not used to round or reconstruct these integers.

Generate windows by starting with the empty word at zero and,
at each stage, adjoining every direction whose unoriented bond has
not yet been used. Each generated word is a trail. Conversely,
deleting the last step of a trail gives a previously generated word,
so induction on its length proves that every window is generated
exactly once. Lemma~\ref{l:H2-dihedral} proves that the canonical
codes give exactly the needed orbits. The finite data checks are:
the stored codes are precisely those \(115\) canonical codes,
without repetition, and their integer values are positive.
\certpending{h2-vector-indexing}
Every regenerated row satisfies the displayed integer inequality.
\certpending{h2-row-inequalities}
In the three cases the weights are exactly
\[
 \frac{1000}{2883}=\frac{935000}{2695605},\qquad
 \frac{200}{561}=\frac{961000}{2695605},\qquad
 \frac15=\frac{539121}{2695605}.
\]
Thus the checked sums are the rows of the mathematically defined
matrix, with no matrix supplied as an input.

To specify also the equalities in the proposition, the integer
extremum checks are
\[
 \min v=186377076276\quad(\text{canonical code }441),\qquad
 \max v=1000000000000\quad(\text{canonical code }0).
\]
\certpending{h2-vector-extrema}
At canonical code \(428\), the regenerated row has
\[
 v_{428}=224087522845,\qquad
 \sum_d n_{k(\omega,d)}v_{\omega d}=603787640626036713.
\]
Exact cancellation gives
\[
 \frac{603787640626036713}{2695605\cdot224087522845}
 =\frac{1119948287353}{1120437614225}=q_*.
\]
\certpending{h2-maximizing-row}
The upper row checks and this equality establish the asserted
maximum; the two extremum checks establish its asserted vector
ratio. The remaining strict comparisons follow by multiplying
positive integer denominators. These are finite integer checks on
the archived witness, rather than assumptions about a computed
eigenvalue or a numerical power iteration.
\end{proof}

\begin{lemma}\label{l:H2-CV}
With \(m=6\), \(q=q_*\), and the vector in Proposition~\ref{p:H2-test},
the constant supplied by Lemma~\ref{l:alm} is at most \(10\).
Indeed, uniformly over the prescribed start and incoming direction,
\[
 W_r\le x_{\rm t}^r\quad(0\le r\le6),\qquad
 \max\left\{\max_{r<6}W_rq^{-r},\
                 \frac{\max v}{\min v}W_6q^{-6}\right\}
 \le6\left(\frac{200}{187q}\right)^6<10.
\]
Consequently (H2) holds with the explicit rational constants
\(q=q_*\) and \(C_V=10\).
\end{lemma}

\begin{proof}
The first step has at most three choices because the designated
initial reversal is forbidden. At every subsequent step the last
bond is used, again leaving at most three choices. Every allowed
weight is at most \(x_{\rm t}/3\), since
\(1/5\le1/3\le x_{\rm s}/3\le x_{\rm t}/3\).
For each prefix, its total extension weight is therefore at most
\(x_{\rm t}\) times its weight. Induction, starting with \(W_0=1\),
gives \(W_r\le x_{\rm t}^r\), in particular for \(0\le r\le6\).

Put \(b=x_{\rm t}/q>1\). For \(r<6\),
\(W_rq^{-r}\le b^r\le b^6\), and the second term in
Lemma~\ref{l:alm} is at most \(6b^6\) by
Proposition~\ref{p:H2-test}. This proves the non-strict bound
in the statement. To verify its last inequality by small integer
comparisons, cross multiplication gives
\(q_*>999/1000\), and hence
\[
 \frac{200}{187q_*}<\frac{200000}{186813}<\frac{15}{14}.
\]
The second inequality is \(2800000<2802195\).
Finally
\(6\cdot15^6=68343750<75295360=10\cdot14^6\).
Thus \(6b^6<10\), and Lemma~\ref{l:alm} permits \(C_V=10\).
\end{proof}

\begin{lemma}\label{l:automaton-quotient}
Let \(\pi:\Omega_m\to Q\) be a surjection onto a finite set.
Suppose that windows with the same image have, for each absolute
direction, identical admissibility and kind, and their admissible
successors have the same image. Suppose also that two such windows
have the same final direction whenever their class has an admissible
non-revisit transition.
Then admissibility, kind and successor define a well-defined automaton
on \(Q\). Every phase-and-position transition whose coefficients
depend only on the absolute direction, kind, position, and final
direction on non-revisit edges factors through this automaton.
For every function \(G(\varphi,q,x)\), its pullback by \(\pi\)
intertwines the raw and quotient backward operators. This
intertwining also holds for each finite stopped continuation sum,
provided that the raw terminal indicator and terminal reward are
pullbacks of their quotient counterparts.
\end{lemma}

\begin{proof}
For \(q\in Q\), choose any \(\omega\) with \(\pi(\omega)=q\).
Declare direction \(d\) admissible at \(q\) exactly when it is
admissible at \(\omega\), give it kind \(k(\omega,d)\), and give
it successor \(\pi(\omega d)\). The hypotheses make all three
choices independent of \(\omega\). When the coefficient of a
non-revisit edge needs the incoming direction, use \(a(\omega)\).
That direction is independent of the representative in every class
with such an edge. In a class with only revisit edges no coefficient
needs it. This proves well-definedness even when such a class
contains several final directions.

Here is the operator assertion including possible parallel edges.
At a raw state \((\varphi,\omega,x)\), write its backward row as
the finite sum over admissible directions \(d\) and allowed next
phases \(\psi\), with coefficient
\(t(\varphi,\omega,x;d,\psi)\). The quotient row has the same
summation indices, coefficient and new position \(x+d\), and its
new window is \(\pi(\omega d)\). If
\(\pi^*G(\varphi,\omega,x)=G(\varphi,\pi(\omega),x)\), termwise
comparison gives
\[
 T_{\rm raw}(\pi^*G)=\pi^*(T_QG).
\]
There is no division by a class size: this is an equality of
backward rows, retaining every direction separately.

For stopped sums the terminal reward and stopping rule must also
factor through the displayed state, as they do for all the terminal
sets used here. The zero-step reward then pulls back, and the
intertwining identity inductively gives
\(T_{\rm raw}^jg_{\rm raw}=\pi^*(T_Q^jg_Q)\) for each \(j\).
Summing finitely many such identities proves the stopped-sum
assertion. No property of a partition-minimization algorithm is used.
\end{proof}

\begin{proposition}\label{p:automaton10}
There are \(60092\) raw windows in \(\Omega_{10}\).
The supplied map to \(Q=\{0,\ldots,3740\}\), with its supplied
representative per class, is surjective and satisfies every
hypothesis of Lemma~\ref{l:automaton-quotient}.
Its admissible labelled edges consist of \(2516\) straight,
\(5032\) turning and \(1872\) revisit edges.
\end{proposition}

\begin{proof}[Finite reconstruction contract]
For this file, in contrast to the H2 orbit file, the direction code
is most-significant-first:
\[
 e_m(d_1,\ldots,d_m)=\sum_{i=1}^m d_i4^{m-i}.
\]
The header of \texttt{windows\_m10\_min.txt} is
\((60092,3741)\). It is followed by \(60092\) pairs
\((e_{10}(\omega),\pi(\omega))\), and then one row per class.
A class row consists of its representative code, four successor
class numbers and four kinds, in absolute order E,N,W,S.
An inadmissible direction has successor and kind both \(-1\);
the kind numbers \(0,1,2\) mean \({\rm s},{\rm t},{\rm r}\).

The generating recursion in the proof of
Proposition~\ref{p:H2-test}, now continued to depth ten and ordered
lexicographically, defines the index set independently of the file.
Its finite cardinality checks, including useful intermediate counts,
are
\[
 \begin{array}{c|rrrrrrrrrr}
 j&1&2&3&4&5&6&7&8&9&10\\ \hline
 |\Omega_j|&4&12&36&108&316&916&2628&7500&21268&60092.
 \end{array}
\]
\certpending{h3-window-enumeration}
As proved for that recursion, no word is missing or duplicated.
The remaining finite assertions are these: the map rows list exactly
these words in that order; all class numbers lie in
\(\{0,\ldots,3740\}\); each representative is a generated word
and maps to its own class; and the file has exactly the stated
number of rows.
\certpending{h3-quotient-indexing}

For every raw word, regenerate its endpoint, its used bonds, and its
earlier departure vertices. For each of its four directions, compare
admissibility and kind with its representative, compare the image
of its successor with the stored successor, and, for every admitted
non-revisit direction, compare final directions. The finite check
requires equality in each comparison. Thus it verifies precisely
the hypotheses of Lemma~\ref{l:automaton-quotient}, including the
final-direction condition. Representatives mapping to their classes
prove surjectivity. Counting the admitted directions in the class
rows gives the integer assertions \(2516,5032,1872\), respectively.
\certpending{h3-quotient-transitions}
This specifies the finite automaton check completely and makes no
minimality or spatial-symmetry assertion.
\end{proof}

\begin{definition}\label{d:phaseweights}
The phases are \({\rm A},{\rm B},{\rm C},{\rm D}\). Their
diagonal transition weights, in kind order \({\rm s},{\rm t},{\rm r}\), are
\[
 \begin{array}{c|ccc}
 {\rm A}&x_{\rm s}/3&x_{\rm t}/3&1/5\\
 {\rm B}&x_{\rm s}/3&1/3&1/5\\
 {\rm C}&x_{\rm s}/3&1/3&1/5\\
 {\rm D}&1/3&1/3&1/5
 \end{array}
\]
At a first visit with reward \(\kappa_h\), additionally allow
\({\rm A}\to{\rm C}\) with weight \(x_{\rm t}\kappa_h/3\) on a turn,
or \({\rm B}\to{\rm D}\) with weight \(x_{\rm s}\kappa_h/3\)
on a straight passage. No other phase change is allowed.
At a revisit there is no phase change.
Start with weight one in each of \({\rm A},{\rm B}\);
the terminal reward is \(g=(0,0,1,1)\).
For a true trail, use its actual revisit kinds and its actual
type-dependent live-pair reward from Definition~\ref{d:kernel}.
\end{definition}

Here \(\kappa_h=\sum_{Y,s}\kappa_{Y,s}h_Y\). In the free
three-neighbour rule its two possible values are
\[
 \kappa_{h,\rm s}^{\circ}=a_{\rm s}+2b_{\rm s}h_C,
 \qquad
 \kappa_{h,\rm t}^{\circ}=b_{\rm t}+(a_{\rm t}+b_{\rm t})h_C.
\]
Indeed, the pair of non-outgoing neighbours is opposite on a
straight passage and perpendicular on a turn; among the two pairs
with the outgoing neighbour, both are perpendicular in the straight
case, and one is opposite in the turning case. A sole surviving
pair instead gives \(b'_\tau h_Y\) with its actual pair type \(Y\).

\begin{lemma}\label{l:phase-expansion}
Let \(w\in\mathcal T_n^X\) with \(n>12\), and let \(\sigma_w\)
be the ordering of Definition~\ref{d:kernel-order}: straight first
departures in increasing path order, followed by turning first
departures in increasing path order. Then its \(h\)-weighted
contribution to \(M_1\) equals the sum over phase histories starting in
\({\rm A}\) or \({\rm B}\) and ending in \({\rm C}\) or \({\rm D}\)
of the products of the weights in Definition~\ref{d:phaseweights}.
Equivalently this sum is
\[
 3^{-|w|}(3/5)^{\operatorname{rev}(w)}
       \sum_k\kappa_h(\sigma_w(k))\prod_{l\le k}x_{\sigma_w(l)}.
\]
Splitting such a phase history at any prefix length gives the
product of its exact prefix weight and its corresponding suffix weight.
\end{lemma}

\begin{proof}
Put \(n=|w|\). An allowed phase history which starts in A and
has nonzero terminal reward must change from A to C exactly once.
There is no other way to enter C, and no transition leaves C.
Its change occurs at a unique turning first departure \(j\).
Before that departure the diagonal factors are
\(x_{\rm s}/3,x_{\rm t}/3,1/5\) on straight, turning and revisit
departures. At \(j\) the factor is
\(x_{\rm t}\kappa_h(j)/3\). After it the factors are
\(x_{\rm s}/3,1/3,1/5\). Their product is therefore
\[
 3^{-n}(3/5)^{\operatorname{rev}(w)}\kappa_h(j)
 x_{\rm s}^{\#\{\text{straight first departures}\}}
 x_{\rm t}^{\#\{\text{turning first departures at indices }i\le j\}}.
\]
These are precisely the first departures at or before \(j\) in
the straight-first ordering, including all straight first departures
even if their path indices exceed \(j\).

A history starting in B with nonzero terminal reward changes from
B to D exactly once, at a straight first departure \(j\). Its
factors before, at and after that departure similarly give
\[
 3^{-n}(3/5)^{\operatorname{rev}(w)}\kappa_h(j)
 x_{\rm s}^{\#\{\text{straight first departures at indices }i\le j\}}.
\]
These are exactly the first departures at or before that \(j\) in
the same ordering. No revisit can be the changing departure.
There is thus one allowed history for each candidate responding
first departure; if its reward is zero both corresponding terms
are zero. Summing over all such departures gives the displayed
identity. Since \(n>12\), this is the ordering fixed in
Definition~\ref{d:kernel-order}, hence the claimed contribution
to \(M_1h\).

Finally split a history after \(r\) departures, \(0\le r\le n\).
Its edge-weight product is the product of the factors with indices
less than \(r\) and those with indices at least \(r\).
Grouping by the phase at the split proves the factorization.
When used for a true suffix, its weights retain the prefix history
in deciding revisits and dead neighbours. The assertion requires
no independence or loss of that history.
\end{proof}

\begin{definition}\label{d:stationary-walk}
Fix \(X=B,C\), \(d_0\in\{{\rm N},{\rm W},{\rm S}\}\),
the sites \(u,c,v\) of Definition~\ref{d:kernel}, and \(a=u+d_0\).
The state space is
\(\{{\rm A},{\rm B},{\rm C},{\rm D}\}\times Q\times\mathbb Z^2\).
At positions \(c,v\) there are no outgoing transitions.
At any other position \(x\), an admissible automaton direction \(d\)
goes to position \(x+d\) and the quotient successor window.
At \(x=u\) or \(a\), forbid reuse of \(\{u,a\}\) in either
direction and kill every memory-detected revisit departure;
every other allowed departure there has weight \(1/5\),
in the same phase and with no response.
At other positions use the automaton kind and the phase weights.
For a non-revisit response, declare only \(c\) dead among the three
neighbours other than the arrival site, and use the type-dependent
three-neighbour or two-neighbour reward rule of Definition~\ref{d:kernel}.
A transition into \(c\) has zero contribution.
Let \(T_{X,d_0}\) be this nonnegative backward operator, with zero
terminal rows, and \(g(\varphi,q,x)=\mathbf1_{\{x=v\}}g_\varphi\).
The free operator \(\mathsf W\) omits positions and special-site
rules, uses the same quotient automaton and phase weights, and
treats all three response neighbours as live.
\end{definition}

Directions remain separate summation indices even if two automaton
edges share a successor class. The position displacement of an edge
is its actual direction. Thus, for a function \(H\), the operator
means the finite sum
\[
 (T_{X,d_0}H)(\varphi,q,x)
   =\sum_{d,\psi}t_{X,d_0}(\varphi,q,x;d,\psi)
                         H(\psi,q d,x+d),
\]
where forbidden transitions have coefficient zero. One may delete
edges into the centre because its continuation value is zero.
The quotient operator is well-defined by
Lemma~\ref{l:automaton-quotient}.

\begin{lemma}\label{l:stationary-domination}
Let \(P=(v_0,\ldots,v_L)\) be a true prefix avoiding \(c,v\),
with \(v_0=u\), \(v_1=u+d_0\), and \(L\ge m+2\), \(m=10\).
Every true continuation to its first hit of \(v\), together with
each phase history, is dominated by the corresponding path of
\(T_{X,d_0}\) from its prefix phase, last-window class and endpoint.
The comparison holds for every finite suffix horizon and hence
for their nonnegative suprema.
In this comparison:
forgotten true revisits have weight \(1/5\) no greater than any
relaxed first-visit diagonal weight; forgetting dead non-outgoing
neighbours only increases the reward, using \(b'_\tau\le b_\tau\);
and the known-site killing rules cannot kill a true continuation.
Indeed, the known initial departures at times \(0,1\) are outside
the memory at every suffix time, and a memory-detected additional
departure at \(u\) or \(u+d_0\) would be a forbidden third departure.
\end{lemma}

\begin{proof}
Fix a true completed trail extending \(P\), of length \(n>L\),
and one of its phase histories. At a suffix departure time
\(j\), \(L\le j<n\), its last \(m\) directions form a raw
window. The next bond is not a used window bond, and reversing
would reuse its last bond. Its direction is therefore admitted by
the raw automaton and, by Lemma~\ref{l:automaton-quotient}, by the
quotient. The new window class is the quotient successor.
The centre is never visited. The target is not visited before time
\(n\); its terminal reward is exactly the phase reward. Thus the
centre and target stopping rules discard none of these histories.

We verify the permanent initial-site rules explicitly. The bond
\(\{u,a\}\) was used at time zero and cannot be reused in either
direction by a true trail. Both \(u\) and \(a\) have already
been departed from, at times zero and one respectively; these
departures exist because \(L\ge12\), and neither site is the
target. Any suffix departure from either site is consequently a
true revisit and has diagonal weight \(1/5\) and no response.
At a suffix time \(j\), the remembered departure indices are
\(j-m,\ldots,j-1\). Since \(j-m\ge L-m\ge2\), these do not
include times zero or one. A detected revisit at one of the two
known sites would therefore certify another departure there in
addition to its initial departure and its current departure.

Such a third departure is impossible by distinctness of the used
bonds. At \(u\), the first departure consumes one bond, and each
later departure requires a distinct arrival bond and a distinct
departure bond. Three departures would consume at least five
bonds; only three are available when \(c\) is excluded.
Here \(a\) is one of \((-1,1),(-2,0),(-1,-1)\), so it is distinct
from \(u,c,v\) in both rows. At \(a\), the initial arrival and
departure consume two
bonds, and two further departures would each consume two additional
bonds, requiring six among four. Consequently the detected-revisit
killing rule cannot kill the chosen history. An allowed departure
from a known site has exactly its true weight \(1/5\).

Now suppose the departure site is neither known site. If the
window detects a revisit, it contains an earlier departure from
that site; this is a true revisit, so the diagonal factors agree
and no phase change occurs. If the true departure is a revisit
which the window has forgotten, keep the same phase in the relaxed
history. Every relaxed first-visit diagonal factor is at least
\(1/3\), and hence at least the true \(1/5\). No responding
departure is lost in this case, because a true revisit never responds.

At a true first departure the window also has a first departure,
with the same incoming and outgoing directions. Its diagonal factors
agree. For a responding departure compare the rewards as follows.
The centre is truly dead, so every neighbour declared dead in the
relaxation is truly dead. If both true non-outgoing neighbours are
live, the relaxed and true pair rules coincide. If precisely one is
live and the relaxation retains the other as dead, the rules again
coincide. If the relaxation forgets that dead neighbour, its
coefficient on the existing pair increases from \(b'_\tau\) to
\(b_\tau\), and its other added pair coefficients are nonnegative.
If both true non-outgoing neighbours are dead, the true reward is
zero. These cases, together with \(0<b'_\tau\le b_\tau\) and
\(h_B,h_C>0\), prove that the relaxed response factor is at least
the true response factor. The outgoing neighbour cannot be dead
on a true continuation, and a true transition never enters \(c\).

Multiplying the comparisons proves domination of the chosen
history. For a fixed prefix and a specified phase at the split,
its lattice suffix and its suffix phase history are determined by
the successive positions and phases of its relaxed path:
consecutive positions determine the absolute direction. Thus the
map from true suffix histories to relaxed histories is injective even
though raw windows have been quotiented. Summing the comparisons
over histories and lengths up to any finite suffix horizon proves
the finite-horizon assertion. Taking suprema of these nonnegative
finite sums proves the final assertion. No rule is released at a
later time.
\end{proof}

\begin{lemma}\label{l:Bellman-soundness}
Let \(T\) be a nonnegative, finitely branching weighted transition
operator on a set of states, with zero rows on a terminal set.
Let \(g\ge0\) be supported on that terminal set.
If a finite-valued function \(F\ge0\) satisfies
\(g+TF\le F\), then for every integer \(n\ge0\),
\[
 \sum_{j=0}^{n}T^jg\le F .
\]
Consequently the nonnegative sum of all finite first-hit path
weights times their terminal rewards is finite and at most \(F\).
It suffices that \(F=g\) on terminal states and \(TF\le F\)
elsewhere.
\end{lemma}

\begin{proof}
Write \(H_n=\sum_{j=0}^nT^jg\). Each value is a finite sum,
because the operator is finitely branching and \(n\) is finite.
Nonnegativity implies that \(T\) preserves pointwise inequalities.
For \(n=0\), \(H_0=g\le g+TF\le F\). If \(H_n\le F\), then
\[
 H_{n+1}=g+TH_n\le g+TF\le F.
\]
This proves the assertion by induction.

Expanding \(T^jg\) as a finite sum over its successive edges
shows that it counts length-\(j\) paths ending at a state with
terminal reward. A path visiting a terminal state earlier has a
zero outgoing factor, so only first-hit paths have a nonzero
contribution. Conversely each such path appears once, with its
edge-weight product and its terminal reward. The nonnegative sum
over all finite lengths is by definition \(\sup_n H_n\), bounded
by the finite number \(F\) at each initial state. At a terminal
state the row is zero, so \(F=g\) implies the supersolution
inequality there; outside the terminal set \(g=0\), so it is just
\(TF\le F\). This proves the sufficiency statement, including
terminal states of reward zero.
\end{proof}

\begin{definition}\label{d:stationary-certificate}
For a case \((X,d_0)\), fix \(R=70\),
\(B_R=\{-R,\ldots,R\}^2\), \(V=2^{16}\), and \(D_3=2^{48}\).
Supply integers \(Z_{\varphi,q}>0\) and
\(U_{\varphi,q,x}\ge0\), \(x\in B_R\).
Set \(\widetilde U_{\varphi,q,x}=U_{\varphi,q,x}\) inside \(B_R\)
and \(Z_{\varphi,q}\) outside, and \(F=\widetilde U/V\), \(z=Z/V\).
For each individual transition coefficient \(t\) of the exact
stationary or free operator, use the proved rational ceiling
\(\widehat t=p(t)/D_3\), \(p(t)=\lceil D_3t\rceil\);
zero coefficients may be omitted.
The certificate checks all the following:
\[
 \begin{gathered}
 \sum_{\text{free edges }e\text{ from }(\varphi,q)}
       p(t_e)Z_{\operatorname{dest}(e)}
       \le D_3Z_{\varphi,q},\qquad Z_{{\rm C},q},Z_{{\rm D},q}\ge V,\\
 0\le U_{\varphi,q,x}\le Z_{\varphi,q},\qquad
 U_{\varphi,q,c}=0,\qquad U_{\varphi,q,v}=Vg_\varphi,\\
 \sum_{\text{stationary edges }e\text{ from }(\varphi,q,x)}
       p(t_e)\widetilde U_{\operatorname{dest}(e)}
       \le D_3 U_{\varphi,q,x}
       \quad(x\in B_R\setminus\{c,v\}).
 \end{gathered}
\]
All tests are integer inequalities on the regenerated automaton,
directions, phases and lattice sites, not on a trusted supplied matrix.
\end{definition}

An individual transition here fixes a direction and a next phase;
the response pairs belonging to that phase change are summed before
its coefficient is rounded. The definition uses ceilings only for
the finite check; the operators themselves retain their exact
rational coefficients.

For explicit reconstruction of those coefficients, encode directions
by \(0,1,2,3=\mathrm E,\mathrm N,\mathrm W,\mathrm S\).
For an admitted non-revisit edge with incoming direction \(a\)
and outgoing direction \(d\), put
\(\{f,g\}=\{0,1,2,3\}\setminus\{a+2,d\}\), with addition
modulo four. For distinct directions define
\(H(e,e')=1\) for opposite directions and \(H(e,e')=h_C\)
for perpendicular directions. If \(\delta\) specifies the direction
from the site to \(c\), use \(\delta=4\) when \(c\) is not a
neighbour. For an edge not entering \(c\), its response reward is
\[
 K(a,d,\delta)=
 \begin{cases}
 b'_\tau H(g,d),&\delta=f,\\
 b'_\tau H(f,d),&\delta=g,\\
 a_\tau H(f,g)+b_\tau[H(f,d)+H(g,d)],&\delta\notin\{f,g\},
 \end{cases}
 \qquad \tau=\begin{cases}{\rm s},&d=a,\\{\rm t},&d\ne a.\end{cases}
\]
The change-of-phase coefficient is \(x_\tau K/3\). For the free
operator use \(\delta=4\). Diagonal coefficients and revisit
coefficients are those in Definition~\ref{d:phaseweights};
at known sites apply the stated killing and \(1/5\) rules.
If \(t=A/B\ge0\) with integers \(A\ge0,B>0\), its ceiling is
\[
 p(t)=\left\lfloor\frac{D_3A+B-1}{B}\right\rfloor,
 \qquad D_3A\le Bp(t),
\]
and, when \(A>0\), \(B(p(t)-1)<D_3A\). These integer relations
prove \(t\le p(t)/D_3\), including \(t=0\) with \(p(t)=0\).
In particular the diagonal ceiling numerators, in the order
\(x_{\rm s}/3,x_{\rm t}/3,1/3,1/5\), are
\[
 97632666219444,\quad100347585280092,\quad
 93824992236886,\quad56294995342132 .
\]
All response ceilings are specified by the same rational formula
with the exact constants of Lemma~\ref{l:log-mean-bounds} and the
exact dyadic \(h_C\); a decimal replacement of \(h_C\) is not used.
\certpending{h3-transition-ceilings}

\begin{lemma}\label{l:stationary-local-global}
Every certificate of Definition~\ref{d:stationary-certificate}
defines a finite-valued \(F\) on the entire lattice with
\(F=z\) outside \(B_R\), \(F\le z\), and
\[
 g+T_{X,d_0}F\le F .
\]
For an inside site this follows from its checked integer row and
coefficient ceilings. For an outside site every successor value
is at most its corresponding \(z\), and the exact stationary row
is dominated by the free row; the checked free inequality
implies \(\mathsf Wz\le z\).
Thus \(F\) bounds every stopped continuation as in
Lemma~\ref{l:Bellman-soundness}.
\end{lemma}

\begin{proof}
There are finitely many phases and window classes and finitely many
stored sites. Their integer values, and each positive integer
\(Z_{\varphi,q}\), are finite, so the extension defines a
nonnegative finite-valued \(F\) at every state. The caps prove
\(F\le z\) everywhere; outside the square equality holds.

Divide a checked free integer row by \(D_3V>0\). Since every
coefficient is at most its ceiling and every \(Z\) is nonnegative,
the result implies
\[
 (\mathsf Wz)_{\varphi,q}
 \le\frac1{D_3V}\sum_{e}p(t_e)Z_{\operatorname{dest}(e)}
 \le Z_{\varphi,q}/V=z_{\varphi,q}.
\]
At an inside nonterminal site the identical argument, with
\(\widetilde U\) at each successor, gives \(TF\le F\).
This includes every edge leaving the square; its successor is
assigned the specified \(Z/V\), with no periodic identification.

We justify also the free-row comparison used outside. Ignoring
positions, a non-special stationary edge has the same diagonal
coefficient as the corresponding free edge. Its response coefficient
is no greater: when one non-outgoing neighbour is the centre the
surviving coefficient \(b'_\tau\) is at most the free coefficient
\(b_\tau\), and the other free pair contributions are nonnegative.
Otherwise the response coefficients agree. Removing an edge into
the centre can only decrease a row. At a known site the diagonal
coefficient is reduced to \(1/5\), which is no larger than any
free diagonal coefficient, all response coefficients are removed,
and additional edges may be killed. At a terminal site all outgoing
coefficients vanish. Thus every stationary row is dominated, after
positions are forgotten, by the free row with the same phase and
window, retaining directions and next phases separately.

For an outside site, each successor value is at most its
corresponding \(z_{\psi,q'}\). The preceding comparison therefore
gives
\[
 (TF)(\varphi,q,x)\le(\mathsf Wz)_{\varphi,q}
                       \le z_{\varphi,q}=F(\varphi,q,x).
\]
Both terminals \(c,v\) are inside \(B_{70}\). Their checked
values give \(F=g\) there, with zero outgoing rows; at all other
states \(g=0\). This proves \(g+TF\le F\) everywhere, and
Lemma~\ref{l:Bellman-soundness} applies. The terminal lower bounds
on \(Z\) are also consistent with the target caps and ensure that
the target reward never exceeds its corresponding outside cap.
\end{proof}

\begin{definition}\label{d:prefix-bound}
For each case \((X,d_0)\), let \(S_{X,d_0}\) be the exact
\(h\)-weighted kernel sum over trails in \(\mathcal T_n^X\) of lengths
\(1\le n\le L=12\)
with first direction \(d_0\), using their prescribed short-trail orderings.
Let \(\mathcal P_{X,d_0,L}\) be all true length-\(L\) prefixes
avoiding \(c\) and \(v\), with first direction \(d_0\).
For \(P\) in this set, let \(p_\varphi(P)\) be the sum of its
exact phase-history products ending in phase \(\varphi\), starting
with weights one in A and B. Use actual trail revisits, actual dead
neighbours and the type-dependent rewards, with the long-trail
straight-first phase weights.
Write \(x(P)\) for its endpoint and \(q(P)\) for its last-window class.
For the potential of its case define the rational number
\[
 J_{X,d_0}=S_{X,d_0}
       +\sum_{P\in\mathcal P_{X,d_0,12}}\sum_\varphi
                    p_\varphi(P)F(\varphi,q(P),x(P)).
\]
A hit at time \(12\) is in \(S_{X,d_0}\), not in the prefix sum.
\end{definition}

Here is a finite recursive definition of the enumeration and of all
its rational summands. It also fixes precisely what is to be
evaluated in Proposition~\ref{p:stationary-check}. A node at time
\(j\le12\) records its full direction word, current position \(x\),
used unoriented bonds \(E\), and counts \(N(y)\) of earlier
departures from every encountered site. Initially these are
\(j=0,x=u,E=\varnothing,N=0\), with incoming direction
\(u-c=\mathrm W\). Its four phase weights are
\((p_{\rm A},p_{\rm B},p_{\rm C},p_{\rm D})=(1,1,0,0)\).
There are three mutually exclusive terminal cases for a node:
\begin{enumerate}
\item If \(x=v\), add its exact short-trail value to \(S_{X,d_0}\)
and produce no children. Test this case before the length cutoff.
\item If \(j=12\) and \(x\ne v\), add its phase-vector dot product
with \(F(\cdot,q(P),x)\) to the prefix sum, and produce no children.
\item Otherwise, try each of the four directions \(d\). At
\(j=0\) retain only \(d=d_0\). At every time discard a move to
\(c\), a move along a used unoriented bond, and the initial
reverse direction. For every remaining move, increment \(N(x)\),
adjoin its bond to \(E\), move to \(x+d\), and recurse.
\end{enumerate}
All children use copies of the parent's data. At a noninitial node
the reverse of its incoming direction is already excluded by the
used-bond test. The true kind is a revisit exactly when the
preincrement value \(N(x)\) is positive. On a first departure
compute \(\kappa=\kappa_h\) using the predeparture counts: a
neighbour \(y\) is dead exactly when \(y=c\) or \(N(y)\ge2\).
The outgoing neighbour of an allowed true extension is not dead.
Use the actual three-, two-, or zero-pair rule of
Definition~\ref{d:kernel}, without a ceiling.

The phase-vector update for a revisit divides all four components
by five. For a straight first departure it is
\[
 (p_{\rm A},p_{\rm B},p_{\rm C},p_{\rm D})\longmapsto
 \left(\frac{x_{\rm s}p_{\rm A}}3,
       \frac{x_{\rm s}p_{\rm B}}3,
       \frac{x_{\rm s}p_{\rm C}}3,
       \frac{p_{\rm D}+x_{\rm s}\kappa p_{\rm B}}3\right);
\]
for a turning first departure it is
\[
 (p_{\rm A},p_{\rm B},p_{\rm C},p_{\rm D})\longmapsto
 \left(\frac{x_{\rm t}p_{\rm A}}3,
       \frac{p_{\rm B}}3,
       \frac{p_{\rm C}+x_{\rm t}\kappa p_{\rm A}}3,
       \frac{p_{\rm D}}3\right).
\]
These formulas use the old components on their right sides.
For a hit at time \(n\le12\), recover the list of actual first
departures \((x_i,\kappa_h(i))\), sort it by the key in
Definition~\ref{d:kernel-order}, including its path-index tie rule,
and evaluate
\[
 3^{-n}(3/5)^{\operatorname{rev}(w)}
       \sum_k\kappa_h(\sigma_w(k))\prod_{l\le k}x_{\sigma_w(l)}.
\]
The phase weights of a short hit are not substituted for this
short-trail expression.

We record why this recursion is exact. Inductively, every generated
node is a true bond trail avoiding \(c\), with no previous hit of
\(v\), except for a terminal node whose endpoint is its first
hit. The used set and departure counts agree with the word by
their update rules. Conversely, every such word of positive length
has its unique shorter parent, and its last direction passes
exactly the tests above. Induction on length therefore proves
completeness without duplication, including all first hits at
length twelve. At any revisited noninitial site each previous
arrival/departure pair uses two distinct bonds. A third departure
would require six bonds. At \(u\), a third departure would require
five, whereas only three avoid \(c\). Thus no branch can depart
from a twice-departed neighbour; equivalently every allowed outgoing
neighbour is live. These observations justify both the departure
counts and the actual dead-site rule at every node. More explicitly,
a twice-departed noninitial site has used all four of its incident
bonds, so a further arrival there is impossible. A twice-departed
\(u\) has used all three bonds available after exclusion of \(c\).
This proves the asserted outgoing-neighbour property even for a
move that would otherwise be the last step of a branch.

Induction on length using the displayed updates proves that the
phase vector is exactly the sum of phase-history products. Keeping
the full word and counts makes each \(\kappa_h(i)\) its value at
the time of that departure; later changes of a neighbour's status
do not change an earlier reward. The short sorting is an exact
comparison of rationals: for positive rewards compare
\[
 (x_i-1)x_j\kappa_h(j)\quad\hbox{and}\quad
 (x_j-1)x_i\kappa_h(i).
\]
Zero rewards have key \(+\infty\). If equal-key entries are grouped
in an implementation, their value is unchanged by exchanging
them: with preceding product \(P\), the difference between the
two orders is
\[
 P\bigl[x_j\kappa_h(j)(x_i-1)
         -x_i\kappa_h(i)(x_j-1)\bigr]=0.
\]
The product following both entries is also unchanged. This proves
agreement with the path-index tie rule even for a grouped evaluation.

There are finitely many branches, of depth at most twelve.
Every coefficient, every potential value and every summand is
rational. For example a common denominator for all possible phase
edge coefficients is
\[
 \Delta=12139925673538045870080000.
\]
It is divisible by the denominators of \(1/3,1/5,x_{\rm s}/3,
x_{\rm t}/3\) and of every possible
\(x_\tau\kappa_h/3\) from the two types and the four choices of
dead non-outgoing neighbours. This is a finite rational divisibility
check with the displayed constants. \certpending{h3-prefix-denominator}
Multiplying phase weights by
\(\Delta^j\) makes every update integral. If \(A\) is the sum
of the resulting length-twelve phase numerators times the
corresponding stored \(U\) entries, its exact tail contribution is
\(A/(\Delta^{12}V)\). Adding the independently evaluated rational
short part specifies \(J_{X,d_0}\) by finite integer arithmetic.

For a compact endpoint lookup, every length-twelve endpoint satisfies
\(|x-u|_1\le12\) and
\((x_1-u_1)+(x_2-u_2)\equiv12\pmod2\): each nearest-neighbour
step changes this parity and changes distance by at most one.
All these sites lie in \(B_{70}\). A compact table must copy the
entry with spatial index
\((x_1+70)141+(x_2+70)\) from the full potential, for every phase
and referenced window class. This identity, checked for every
referenced coordinate, justifies the lookup. Parity is used only
to identify possible endpoints, and introduces no numerical factor.

\begin{lemma}\label{l:prefix-domination}
For every case whose integer certificate passes, the \(h\)-weighted
contribution of all kernel trails of that first direction,
over every length, is at most \(J_{X,d_0}\).
This holds uniformly for every finite length cutoff, with the same
upper bound. The short part uses the exact order through \(12\);
all longer trails split uniquely into a surviving length-\(12\)
prefix and a dominated stopped suffix.
\end{lemma}

\begin{proof}
Fix a finite total-length cutoff \(R_0\). Trails of length at most
\(\min(R_0,12)\) contribute a sub-sum of \(S_{X,d_0}\), because
their weights are nonnegative and their orderings are exactly the
specified short orderings. This includes any hit at time twelve.

If \(R_0>12\), every remaining true trail has a unique first twelve
steps \(P\), which avoid \(c\) and \(v\), and a suffix of positive
length at most \(R_0-12\), stopping at its first hit of \(v\).
Its prescribed ordering is straight-first, so
Lemma~\ref{l:phase-expansion} applies to its full contribution.
Split each phase history after those twelve departures and group
by its phase \(\varphi\). The sum of the prefix factors is exactly
\(p_\varphi(P)\). Lemma~\ref{l:stationary-domination} bounds the
sum of its true suffix factors and terminal rewards by the
corresponding stationary stopped sum from
\((\varphi,q(P),x(P))\). This comparison is valid for the chosen
cutoff and with the full true prefix as initial history.
Lemmas~\ref{l:stationary-local-global} and
\ref{l:Bellman-soundness} bound that stopped sum by
\(F(\varphi,q(P),x(P))\), independently of the cutoff.

There are finitely many length-twelve prefixes. For each, multiplying
by its nonnegative phase weights and then summing yields the prefix
term in \(J_{X,d_0}\). Even if two prefixes have the same class,
endpoint and phase, their separate prefix weights remain in the
sum, so quotienting loses no multiplicity. Adding the short bound
gives the assertion for \(R_0\). For \(R_0\le12\), including
\(R_0=0\), the nonnegative prefix term may simply be added to the
short bound. Taking the supremum over all finite cutoffs proves
the all-length assertion.
\end{proof}

\begin{lemma}\label{l:kernel-reflection}
For row B, reflection \((x_1,x_2)\mapsto(x_1,-x_2)\) fixes
\(c,u,v\), interchanges the first directions N and S, and preserves
every trail weight, type, dead pair, shift, \(h\)-reward and
prescribed ordering. Hence the true B/N and B/S kernel contributions
are equal at every length and every finite cutoff.
No corresponding reduction is asserted for row C.
\end{lemma}

\begin{proof}
Let \(r(x_1,x_2)=(x_1,-x_2)\). In row B it fixes the three
specified sites and the incoming direction \(\mathrm W\), and
interchanges \(\mathrm N\) and \(\mathrm S\). It maps each
unoriented nearest-neighbour bond bijectively to such a bond. Thus
\(w\mapsto rw\) is a bijection, inverse to itself, between the
two first-direction classes, preserving length, centre avoidance,
first-hit stopping and the set of path indices of earlier
departures. At every index it preserves straightness, turning,
revisit status and every departure count.

It maps the three non-arrival neighbours to the corresponding
neighbours of the reflected departure, and fixes the centre as a
set. Hence a neighbour is dead if and only if its image is dead.
It preserves the outgoing neighbour and may interchange the two
other slots, whose rules are symmetric. The actual assigned
coefficients \(a_\tau,b_\tau,b'_\tau\) therefore agree for each
corresponding live pair. Opposite pairs remain opposite and
perpendicular pairs remain perpendicular, so their types and
\(h\)-weights agree. Since \(r\) is an isometry and fixes \(u,v\),
\[
 \max_{z\in p}\min\{d(z,u),d(z,v)\}
 =\max_{z\in rp}\min\{d(z,u),d(z,v)\}.
\]
Thus shifts agree as well.

Each first departure has the same rational ordering key on both
trails. Its path index is unchanged, so the short ordering including
tie breaking is preserved. The long straight-first ordering is
preserved for the same reason. All products and response sums
in Definition~\ref{d:kernel} consequently agree. Summing the
bijection at a fixed length, and then over any finite length cutoff,
proves the assertion. No equality of automaton classes or of
stationary potentials under \(r\) is needed.
\end{proof}

\begin{proposition}\label{p:stationary-check}
For the five cases B/N, B/W, C/N, C/W, C/S, there exist
integer arrays \(Z,U\) satisfying every check in
Definition~\ref{d:stationary-certificate}, with
\(m=10,L=12,R=70\) and the specified scales.
Their regenerated quotient has \(3741\) states and the spatial
square has \(19881\) sites. Their exact rational prefix evaluations obey
\[
 \begin{array}{c|ccccc}
 (X,d_0)&(B,{\rm N})&(B,{\rm W})&(C,{\rm N})&(C,{\rm W})&(C,{\rm S})\\ \hline
 10^9J_{X,d_0}\ \le&
 351207554&290599163&527872005&380754375&273048088
 \end{array}
\]
Each displayed upper bound denotes an exact integer inequality
after clearing the rational prefix denominator.
\end{proposition}

\begin{proof}[Finite data and evaluation contract]
The five witnesses are the files
\[
 \texttt{stationary\_}\!X\texttt{\_m10\_R70\_d}j
                  \texttt{\_coarse.bin.gz},
 \qquad (X,j)=(B,1),(B,2),(C,1),(C,2),(C,3),
\]
in the archived data directory. Here \(1,2,3\) mean
\(\mathrm N,\mathrm W,\mathrm S\). After decompression each file
has nine little-endian signed 64-bit header integers
\[
 [10,12,70,3741,19881,v_1,v_2,16,j],
\]
where \((v_1,v_2)=(1,0)\) for B and \((0,1)\) for C.
It then has \(4\cdot3741\) unsigned 32-bit integers for \(Z\),
followed by \(4\cdot3741\cdot19881\) unsigned 32-bit integers
for \(U\), and no other payload. Phases are ordered A,B,C,D,
then class numbers in the order of
Proposition~\ref{p:automaton10}; for \(U\), the last index is
\((x_1+70)141+(x_2+70)\). This identifies each supplied literal
with a mathematical array index. In particular,
\(19881=141^2\), rather than a periodic spatial state count.
\certpending{h3-array-indexing}

The finite assertions on these arrays are exactly every integer
inequality and equality in Definition~\ref{d:stationary-certificate}.
For each of the \(4\cdot3741\) phase/class indices, require
\(Z>0\), its free row inequality, and the C,D lower terminal
bounds. \certpending{h3-free-rows}
For each of the \(4\cdot3741\cdot19881\) stored
phase/class/site indices, require \(0\le U\le Z\) and the specified
center and target equalities. \certpending{h3-caps-and-terminals}
Require the nonterminal stationary row inequality when appropriate.
The coefficients are regenerated
from the exact formulas following that definition, not inferred
from floating values stored in a program.

For clarity, all four direction choices, both diagonal and
change-of-phase possibilities, and every admitted successor are
included in those rows. At a known site remove the initial bond
directions and every detected-revisit direction, use \(1/5\) on
the remaining same-phase edges, and include no response edge.
At the centre and target check the terminal values instead of an
outgoing row. For an outside successor substitute its \(Z\) entry.
The quotient transitions and representative incoming directions
are justified by the raw reconstruction in
Proposition~\ref{p:automaton10}. In particular a class with mixed
incoming directions can be used only on revisit transitions, for
which no response direction is needed.

There are \(4\cdot3741\cdot(19881-2)=297469356\) nonterminal
local component inequalities per file, as well as all the free,
cap and terminal checks just specified. These are exact integer
inequalities over the five supplied arrays; gzip checksums or the
termination status of a native checker are not substitutes for them.
\certpending{h3-local-rows}

Evaluate the finite recursion following
Definition~\ref{d:prefix-bound} against these same arrays, retaining
their exact class and coordinate indices. Completeness of that
recursion and the phase-vector update have already been proved
there. The resulting finite enumeration checks, in the column
order of the proposition, are
\[
 \begin{array}{c|rrrrr}
 &B/{\rm N}&B/{\rm W}&C/{\rm N}&C/{\rm W}&C/{\rm S}\\ \hline
 \text{nodes}&139985&161929&105025&155673&145485\\
 \text{surviving prefixes}&89150&103257&66947&99168&92748\\
 \text{short first hits}&937&922&715&1036&861.
 \end{array}
\]
\certpending{h3-prefix-enumeration}
Counts alone are not the prefix assertion: the check must also
evaluate every rational summand. If the exact short sum is
\(s/d\), \(d>0\), and the nonnegative integer prefix accumulator
is \(A\) with common denominator \(\Delta^{12}V\) as specified
there, the displayed bound \(10^9J\le K\) means exactly
\[
 10^9\bigl(s\Delta^{12}V+Ad\bigr)
                  \le Kd\Delta^{12}V .
\]
Use respectively
\(K=351207554,290599163,527872005,380754375,273048088\).
This gives the finite rational prefix tests without a decimal
rounding convention. \certpending{h3-prefix-values}
A compact lookup is accepted only with the
entry-by-entry identity with the full potential described after
Definition~\ref{d:prefix-bound}.
\certpending{h3-prefix-lookup}
The archived exact fractions in the five prefix records identify
the values of these finite expressions; the expressions and the
cleared inequalities above, together with all the array checks,
constitute the stationary certificates. Their proofs are checked by
the Lean kernel in the accompanying formalization, with pending
checks marked above.
\end{proof}

\begin{lemma}\label{l:stationary-assembly}
With the five certified potentials and their exact prefix bounds,
\[
 (M_1h)_B\le2J_{B,{\rm N}}+J_{B,{\rm W}}\le\rho_0h_B,\qquad
 (M_1h)_C\le J_{C,{\rm N}}+J_{C,{\rm W}}+J_{C,{\rm S}}
                                      \le\rho_0h_C ,
\]
where \(h_C=5404319552844595/4503599627370496\) and
\(\rho_0=993015/10^6<1\).
The same inequalities hold for every finite true-trail cutoff;
therefore (H3\(_h\)) holds for the type-dependent kernel.
\end{lemma}

\begin{proof}
Every true trail from \(u\) has first direction N, W or S: direction
E enters the forbidden centre. For row B,
Lemma~\ref{l:kernel-reflection} identifies the true N and S sums
at every finite cutoff; apply Lemma~\ref{l:prefix-domination}
to B/N and B/W. For row C apply that lemma separately to its
three first-direction classes, with no symmetry reduction.
This proves the first inequalities in the statement, uniformly in
the total-length cutoff.

The exact component ceilings of Proposition~\ref{p:stationary-check}
give
\[
 2J_{B,{\rm N}}+J_{B,{\rm W}}
 \le\frac{2\cdot351207554+290599163}{10^9}
 =\frac{993014271}{10^9}
 <\frac{993015000}{10^9}=\rho_0h_B.
\]
For C their sum is
\[
 J_{C,{\rm N}}+J_{C,{\rm W}}+J_{C,{\rm S}}
 \le\frac{1181674468}{10^9}
 <\frac{119}{100}
 <\frac{993015}{10^6}
             \frac{5404319552844595}{4503599627370496}.
\]
The middle inequality is \(1181674468<1190000000\).
For the last inequality, use the exact rational comparisons
\[
 h_C>1199/1000,\qquad
 993015\cdot1199=1190624985>1190000000.
\]
The first comparison is an exact comparison with the displayed
dyadic fraction. These two inequalities imply
\(\rho_0h_C>119/100\). All comparisons involve positive integer
denominators. In particular they give the desired C-row bound.
Finally \(0<\rho_0<1\) since \(0<993015<1000000\), and
\(h_B,h_C>0\).

The same upper bounds hold for every finite cutoff. By
Definition~\ref{h:H3}, this proves (H3\(_h\)). Equivalently, taking
suprema of the nonnegative row sums gives a finite \(M_1h\)
with the asserted bound. Only a finite prefix was used in the
numerical evaluation; the Bellman inequality controls all suffix
horizons simultaneously.
\end{proof}




\section{Operator certificates for the local inequalities}\label{s:local-operator}

Fix a covered pattern \(\pi\in\mathcal P_{\rm ext}\) from
Definition~\ref{d:pattern}.  Write \(o\) for its out-neighbour and
\(f,g\) for the other two sites in
\(\{\mathsf e,\mathsf N,\mathsf W\}\), whether or not they are
removed.  For \(0<c<1\) set
\[
 F_c(m)=Z_{\rm cur}(\pi;m)-cZ_{\rm next}(\pi;m),
 \qquad m\in(\mathbb S^2)^{20}.
\]
All spin functionals in this appendix are finite sums with the laws
of Definition~\ref{d:moment}.  The proposition at the end gives a
sufficient certificate for this pattern's inequality in
Definition~\ref{h:H1}; it does not assert optimality of \(c\).
Empty products and tensor products mean \(1\) and \(\mathbb C\),
respectively, and the Pauli operator of a constant when there are
no vector variables is that scalar on \(\mathbb C\).

\subsection{Passage and leaf reductions}

\begin{lemma}\label{l:Iw}
Let \(Z_w\) be the finite template functional after deleting
the passage site and its incident factors, retaining all other
present bonds and all covered slot factors. Put \(P_{ab}=n_a\cdot n_b\)
for nonremoved sites and zero if either is removed.
For the out-neighbour \(o\) and the other two neighbours \(f,g\),
\[
 Z_{\rm cur}=Z_w[I],\qquad Z_{\rm next}=Z_w[J]>0,\qquad
 I=1+\tfrac13(P_{fg}+P_{fo}+P_{go}),\quad J=1+\tfrac15P_{fg}.
\]
For \(0<c<1\), \(F_c=Z_{\rm cur}-cZ_{\rm next}=Z_w[I-cJ]\), with
\[
 I-cJ=(1-c)+\tfrac13(P_{fo}+P_{go})+(\tfrac13-\tfrac c5)P_{fg}.
\]
Both partition polynomials are affine in each slot vector and
invariant under simultaneous orthogonal transformations of the slots.
\end{lemma}

\begin{proof}
For purposes of expanding a product, set \(n_x=0\) whenever \(x\)
is removed; this adds neither a variable nor a site law to the
functional.  By Definition~\ref{d:pattern}, the present passage
bonds in the current template go precisely to the nonremoved sites
among \(f,g,o\).  In the next template the passage law is
\(\lambda_{\mathrm D}\), its bond to \(o\) is deleted, and its
other present bonds are unchanged.  Consequently the factors
obtained by summing over the passage spin are, respectively,
\[
 \mathcal M\bigl((1+n\cdot n_f)(1+n\cdot n_g)
                         (1+n\cdot n_o)\bigr),\qquad
 \lambda_{\mathrm D}\bigl((1+n\cdot n_f)(1+n\cdot n_g)\bigr).
\]
The first and third moments vanish.  The second moment is
\(\delta_{ij}/3\) for \(\lambda_{\mathrm U}\) and
\(\delta_{ij}/5\) for \(\lambda_{\mathrm D}\), by
Lemma~\ref{l:moments}.  Expansion therefore gives exactly \(I\)
and \(J\), including when one or both of \(f,g\) are removed.
Finite summation can be performed in any order, proving both
partition identities.  Subtraction gives the asserted formula for
\(I-cJ\).

The network defining \(Z_w[1]\) has at most four incident factors
at each site, so \(Z_w[1]>0\) by Lemma~\ref{l:finitefacts}.
All its weights are nonnegative.  The remaining spins have norm
at most one, so \(P_{fg}\ge-1\), also with the removed-site
convention.  Thus
\[
 Z_{\rm next}=Z_w[J]\ge\frac45 Z_w[1]>0.
\]

Each slot occurs in just one linear factor.  All spin variables
in either original template occur in at most four factors:
the passage has at most three, each tip has its three slots and
at most one bond, each corner has its two slots and at most two
bonds, and each neighbour of the passage has at most four bonds.
For \(R\in\mathrm O(3)\), simultaneous transformation of slots
and spin variables leaves every scalar product unchanged.
The site-by-site polynomial covariance in
Lemma~\ref{l:moments}, with these separate degree bounds, then
shows that transforming the slots alone leaves each partition
polynomial unchanged.  This uses covariance of the indicated
polynomials, without asserting invariance of the finite laws on
arbitrary functions.
\end{proof}

\begin{lemma}\label{l:leaf}
In a covered template, integrating a nonremoved tip or corner
with no present internal bond gives a positive factor depending
only on its own slots.
A tip with its one internal bond to \(x\) gives
\(A(1+(3/5)u\cdot n_x)\), \(A>0,\ |u|\le1\).
A corner with exactly one internal bond to \(x\) gives
\(A(1+\rho u\cdot n_x)\), \(A>0,\ |u|\le1\), where
\(\rho=1/2\) for status \(\mathrm U\) and \(\rho=1/3\)
for status \(\mathrm D\).
\end{lemma}

\begin{proof}
All slot vectors in this proof are unit vectors.  A tip has law
\(\lambda_{\mathrm U}\), three slots, and at most one internal
bond, by Definition~\ref{d:pattern}.  If its slots are
\(a_1,a_2,a_3\), their normalizer is
\[
 A=\mathcal M\left(\prod_{j=1}^3(1+a_j\cdot n)\right)
   =1+\frac13\sum_{j<l}a_j\cdot a_l\ge\frac12,
\]
since \(0\le|a_1+a_2+a_3|^2=3+2\sum_{j<l}a_j\cdot a_l\).
With no bond this is the required positive factor.  With a bond
to \(x\), put
\[
 V=\mathcal M\left(n\prod_{j=1}^3(1+a_j\cdot n)\right).
\]
Summing over the tip gives \(A+V\cdot n_x\).  The normalized
mean bound in Lemma~\ref{l:35} gives \(|V|\le3A/5\), so
\(u=5V/(3A)\) has norm at most one and gives the claimed formula.
The moment in \(V\) has degree four, within the stated bound.

For a nonremoved corner let its two slots be \(a,b\), and put
\(t=a\cdot b\in[-1,1]\).  Write \(A\) for its slot normalizer
and \(V\) for its unnormalized slot mean.  Lemma~\ref{l:moments}
gives
\[
 \begin{array}{c|cc}
   &A&V\\ \hline
  \mathrm U&1+t/3&(a+b)/3\\
  \mathrm D&1+t/5&(a+b)/5 .
 \end{array}
\]
In particular \(A\ge2/3\) in status \(\mathrm U\), and
\(A\ge4/5\) in status \(\mathrm D\).  This proves the assertion
for corners with no internal bond.  If there is exactly one bond
to \(x\), its sum is again \(A+V\cdot n_x\).  As
\(|a+b|^2=2+2t\), in the two cases respectively we have
\[
 A^2-4|V|^2=\frac{(1-t)^2}{9}\ge0,
 \qquad
 A^2-9|V|^2=\frac{(1-t)(7-t)}{25}\ge0.
\]
Since \(A>0\), these inequalities give \(|V|\le A/2\) and
\(|V|\le A/3\).  Set \(u=V/(\rho A)\) with the corresponding
\(\rho\).  Every factor and vector constructed here depends only
on the slots at that leaf, as required.
\end{proof}

\begin{definition}\label{d:reduction}
Delete tips and corners with no present internal bond, dropping
their common positive factors. Replace every bonded tip by a
single vector and radius \(3/5\).
In the reduced representation also replace each one-bond corner
by a single vector of radius \(1/2\) or \(1/3\) according to
Lemma~\ref{l:leaf}; in the unreduced representation retain its spin
and two slots. Other bonded corners retain their two slots.
Let \(u=(u_1,\ldots,u_k)\) list these remaining vectors, and define
\(\widehat Z_{\rm cur},\widehat Z_{\rm next}\) by these replacements,
\(\widehat F_c=\widehat Z_{\rm cur}-c\widehat Z_{\rm next}\).
A group is one single vector or the two slots of a retained corner.
\end{definition}

Thus a single vector attached to a retained spin \(n_x\) means
the factor \(1+r_i u_i\cdot n_x\), where \(r_i\) is its prescribed
radius.  Each retained corner contributes
\((1+u_i\cdot n_z)(1+u_{i'}\cdot n_z)\) and retains its original
status law.  The functional after the passage sum consists of
these factors, the remaining internal bonds, and \(I\) or \(J\)
from Lemma~\ref{l:Iw}.  This specifies the polynomials for all
\(u_i\in\mathbb R^3\); no realizability condition on a replacement
vector is part of their definition.

\begin{corollary}\label{c:reduce}
In either representation, \(\widehat F_c\) is real, affine in
each vector and invariant under simultaneous \(\mathrm O(3)\).
If \(\widehat F_c\ge0\) on \((\mathbb S^2)^k\), then
\(F_c\ge0\) on \((\mathbb S^2)^{20}\).
For each original slot assignment there are a common \(K>0\)
and vectors \(\widehat u\) in the closed unit ball such that
\(Z_{\rm cur}=K\widehat Z_{\rm cur}(\widehat u)\) and
\(Z_{\rm next}=K\widehat Z_{\rm next}(\widehat u)\).
\end{corollary}

\begin{proof}
The slot factors and the replacement factors are real and affine
in each of their vector variables.  No vector is used twice.
Finite summation therefore preserves reality and separate
affineness.  In a term containing a pair insertion from \(I\) or
\(J\), each endpoint is a nonremoved non-arrival neighbour of the
passage.  Deleting its passage bond leaves at most three other
factors, and the insertion raises its degree by one.  Other
neighbours have degree at most three after passage deletion, and
retained corners have at most two slots and two bonds.  Replacing
a leaf changes a bond at its surviving endpoint to a single
linear factor and so does not increase that degree.  Hence all
summed spin variables have separate degree at most four.
Every factor is a scalar product or a constant with a real
radius.  Lemma~\ref{l:moments}, applied site by site as in
Lemma~\ref{l:Iw}, proves simultaneous \(\mathrm O(3)\) covariance
of both reduced partition polynomials, and thus of
\(\widehat F_c\).

Fix an original assignment \(m\in(\mathbb S^2)^{20}\).  No two
tips or corners are nearest neighbours in the template.  Thus
each sum in Lemma~\ref{l:leaf} can be performed with all other
spins fixed, without changing the slots of another leaf.  These
leaf factors contain no passage spin.  The same successive
summations therefore give, in the two partition functions, the
same product \(K>0\) of the leaf normalizers and the same list
\(\widehat u\) of replacement vectors and retained slots.
Lemma~\ref{l:leaf} bounds each replacement vector by one, and
the retained slots are unit vectors.  Consequently
\[
 Z_{\rm cur}=K\widehat Z_{\rm cur}(\widehat u),\qquad
 Z_{\rm next}=K\widehat Z_{\rm next}(\widehat u),\qquad
 F_c(m)=K\widehat F_c(\widehat u),\qquad |\widehat u_i|\le1.
\]
Removed sites have no factors, and an empty product of deleted
leaf normalizers is \(K=1\).

For completeness, if \(|v|\le1\), choose \(b=v/|v|\) when
\(v\ne0\), and choose \(b=(0,0,1)\) when \(v=0\).  Then
\[
 v=\frac{1+|v|}{2}b+\frac{1-|v|}{2}(-b).
\]
Successively using affineness in each vector expresses a value
of \(\widehat F_c\) on the product of balls as a sum of its
values on unit vectors with nonnegative product weights summing
to one.  For no vector variables this is the single value of the
constant polynomial.  Positivity on unit vectors therefore
implies positivity at \(\widehat u\), and the last displayed
identity proves the assertion.
\end{proof}

\subsection{Pauli operators and their symmetries}

We use the explicit matrices
\[
 \sigma_1=\begin{pmatrix}0&1\\1&0\end{pmatrix},\qquad
 \sigma_2=\begin{pmatrix}0&-\mathrm i\\\mathrm i&0\end{pmatrix},\qquad
 \sigma_3=\begin{pmatrix}1&0\\0&-1\end{pmatrix}.
\]
They are Hermitian, have square \(I\), and satisfy
\(\operatorname{Tr}\sigma_j=0\) and
\(\operatorname{Tr}(\sigma_j\sigma_l)=2\delta_{jl}\)
for \(1\le j,l\le3\).  Direct multiplication also gives
\(\sigma_j\sigma_l+\sigma_l\sigma_j=2\delta_{jl}I\).

\begin{definition}\label{d:Pauli}
For a real polynomial affine in each vector, expand uniquely
\[
 F(u)=\sum_{a\in\{0,1,2,3\}^k}f_a\prod_i u_i^{a_i},
 \qquad u_i^0=1,\qquad
 \mathcal O(F)=\sum_a f_a\bigotimes_i\sigma_{a_i},\quad\sigma_0=I.
\]
For \(\mathcal I\subset\{1,\ldots,k\}\), invert those vectors,
writing \(u^{\mathcal I}\) and \(F^{\mathcal I}(u)=F(u^{\mathcal I})\).
For \(|a|\le1\), put \(\rho(a)=(I+a\cdot\sigma)/2\).
Set \(S_z=\frac12\sum_i\sigma_3^{(i)}\) and
\(S_\pm=\sum_i(\sigma_1^{(i)}\pm\mathrm i\sigma_2^{(i)})/2\).
\end{definition}

Here \(u_i^j\) for \(j=1,2,3\) denotes a coordinate, not a power.
The expansion is unique because these are distinct monomials
in the \(3k\) real coordinates.  Separate affineness permits
at most one coordinate from each vector.  Its coefficients are
real, so \(\mathcal O(F)\) is Hermitian.  The map
\(F\mapsto\mathcal O(F)\) is real linear, commutes with finite
sums over auxiliary variables, and takes products on disjoint
sets of vector variables to tensor products of operators.

\begin{lemma}\label{l:inversion}
For \(|u_i|\le1\), each \(\rho(u_i)\) is positive semidefinite
of trace one, and
\[
 F(u)=\operatorname{Tr}\!\left[
      \mathcal O(F^{\mathcal I})
                 \bigotimes_i\rho((u^{\mathcal I})_i)\right].
\]
Thus a single inversion \(\mathcal I\) with
\(\mathcal O(F^{\mathcal I})\ge0\) implies \(F\ge0\)
on the product of closed unit balls, in particular on unit vectors.
\end{lemma}

\begin{proof}
The anticommutation identities give
\((a\cdot\sigma)^2=|a|^2I\).  For \(a\ne0\), the eigenvalues
of this traceless Hermitian matrix are \(|a|,-|a|\); for
\(a=0\) both are zero.  The eigenvalues of \(\rho(a)\) are
therefore \((1+|a|)/2,(1-|a|)/2\), which are nonnegative for
\(|a|\le1\).  Its trace is one.  The displayed trace identities
for the Pauli matrices give
\[
 \operatorname{Tr}(\sigma_j\rho(a))=a^j
       \quad(0\le j\le3),\qquad a^0=1.
\]
For any separately affine polynomial \(G\), multiplicativity
of the trace on tensor products now gives
\[
 \operatorname{Tr}\left[\mathcal O(G)
                          \bigotimes_i\rho(v_i)\right]
       =\sum_a g_a\prod_i v_i^{a_i}=G(v).
\]
Take \(G=F^{\mathcal I}\) and \(v=u^{\mathcal I}\).
Inverting the same coordinates twice is the identity, proving
the asserted formula.

The tensor product of the densities is positive semidefinite,
as follows by tensoring their orthonormal eigenbases.  If
\(\mathcal O(F^{\mathcal I})\) is positive semidefinite, choose
an orthonormal eigenbasis \((b_j)\) of it with eigenvalues
\(\lambda_j\ge0\).  For any positive semidefinite \(R\),
\[
 \operatorname{Tr}(\mathcal O(F^{\mathcal I})R)
       =\sum_j\lambda_j\langle b_j,Rb_j\rangle\ge0.
\]
Apply this with the product density.  The formulas also hold
for an empty tensor product, whose trace is one.
\end{proof}

\begin{lemma}\label{l:Ostruct}
For \(\mathcal I\) a union of vector groups and
\(\mathcal O=\mathcal O(\widehat F_c^{\mathcal I})\),
the matrix in the standard tensor basis is real symmetric.
It commutes with \(U^{\otimes k}\) for every \(U\in\mathrm{SU}(2)\)
and with \(S_z,S_+,S_-\).
For a retained \(\mathrm U\)-corner pair, if \(\Pi\) projects that
pair onto \(\operatorname{Sym}^2\mathbb C^2\), then
\(\mathcal O=\Pi\mathcal O\Pi\).
For a retained \(\mathrm D\)-corner pair, \(\mathcal O\) commutes
with exchange of its two factors.
\end{lemma}

\begin{proof}
Put \(G=\widehat F_c^{\mathcal I}\).  An inversion commutes with
a simultaneous orthogonal transformation, so Corollary~\ref{c:reduce}
implies \(G(Ru_1,\ldots,Ru_k)=G(u)\) for \(R\in\mathrm O(3)\).
Its coefficients are real and hence \(\mathcal O\) is Hermitian.
Entrywise conjugation fixes \(\sigma_0,\sigma_1,\sigma_3\) and
negates \(\sigma_2\).  Thus \(\overline{\mathcal O(G)}\) is the
Pauli operator of the polynomial obtained by reflecting the
second coordinate of every vector.  This is \(G\) itself by
its orthogonal covariance.  The operator is therefore real,
and Hermiticity makes it symmetric.

For \(U\in\mathrm{SU}(2)\), conjugation by \(U\) preserves the
real vector space of traceless Hermitian \(2\times2\) matrices
and its inner product \(\tfrac12\operatorname{Tr}(AB)\).
The Pauli matrices form an orthonormal basis for that space.
Consequently there is a real orthogonal matrix \(R\) with
\[
 U\sigma_jU^*=\sum_{l=1}^3R_{lj}\sigma_l.
\]
Extend it by \(R_{00}=1\), \(R_{0j}=R_{j0}=0\) for \(j>0\).
Expanding every tensor factor in the definition of the Pauli
operator gives
\[
 U^{\otimes k}\mathcal O(G)(U^*)^{\otimes k}
       =\mathcal O\bigl(G(R^{\mathsf T}u_1,\ldots,
                                      R^{\mathsf T}u_k)\bigr)
       =\mathcal O(G).
\]
Only orthogonality of \(R\) is needed here.

The concrete generator commutators follow by a polynomial
coefficient comparison.  For \(j=1,2,3\) and real \(t\), the
matrix
\[
 U_j(t)=\frac{I-\mathrm i t\sigma_j}{\sqrt{1+t^2}}
\]
is unitary of determinant one.  The just-proved commutation,
after multiplication by \((1+t^2)^{k/2}\), shows that
\(\mathcal O\) commutes with \((I-\mathrm i t\sigma_j)^{\otimes k}\)
for every real \(t\).  Its coefficient of \(t\) is
\(-\mathrm i\sum_i\sigma_j^{(i)}\), so \(\mathcal O\) commutes
with all three sums.  Their linear combinations give precisely
\(S_z,S_+,S_-\).  When \(k=0\) all these sums are zero.

Consider a retained corner pair \(i,i'\) at \(z\).
Since \(\mathcal I\) is a union of groups, its two slot factors
after inversion have the same sign \(\eta\in\{1,-1\}\):
\[
 (1+\eta u_i\cdot n_z)(1+\eta u_{i'}\cdot n_z).
\]
With all spins fixed, all other factors are independent of
\(u_i,u_{i'}\).  Real linearity of the Pauli map and its tensor
product property show that every summand of \(\mathcal O(G)\),
with the pair factors ordered first, has the form
\[
 (I+\eta n_z\cdot\sigma)\otimes
 (I+\eta n_z\cdot\sigma)\otimes H_n
\]
times a real scalar, for a matrix \(H_n\) on the other factors.
There is no positivity assumption on that scalar or on \(H_n\).
If \(z\) has status \(\mathrm U\), every value of \(n_z\) is
unit.  The matrix \(I+\eta n_z\cdot\sigma\) is then twice a
rank-one orthogonal projection, by the eigenvalue calculation
in Lemma~\ref{l:inversion}.  Its tensor square has range in
the symmetric square: if its one-dimensional range is spanned
by \(v\), the tensor square has range spanned by \(v\otimes v\).
It is Hermitian, so multiplication by the symmetric projection
on either side leaves it unchanged.  This proves
\(\Pi\mathcal O=\mathcal O=\mathcal O\Pi\), and hence
\(\mathcal O=\Pi\mathcal O\Pi\).

For status \(\mathrm D\), the value \(n_z=0\) is also allowed.
For every \(n_z\), including zero, exchanging the two factors
commutes with the tensor square of \(I+\eta n_z\cdot\sigma\).
The same holds after finite summation.  This proves the asserted
exchange commutation.  Each argument in this proof applies
separately to \(\widehat Z_{\rm cur}^{\mathcal I}\) and
\(\widehat Z_{\rm next}^{\mathcal I}\), as well as to their
linear combination: both polynomials have the covariance and
the corner slot factors used above.
\end{proof}

\subsection{Exchange sectors and the central weight}

\begin{definition}\label{d:exchange}
Let \(V\) be the tensor product of \(\mathbb C^2\) for singles,
\(\operatorname{Sym}^2\mathbb C^2\) for retained \(\mathrm U\)-pairs,
and \(\mathbb C^2\otimes\mathbb C^2\) for retained \(\mathrm D\)-pairs.
For each assignment \(\epsilon\) of exchange signs to the
\(\mathrm D\)-pairs, let \(V_\epsilon\) replace each of those
pair factors by its symmetric or exterior square.
If there are \(t\) singles, let \(\kappa=0\) for even \(t\),
\(\kappa=1/2\) for odd \(t\), and
\(V_\epsilon^\circ=\ker(S_z-\kappa I)|_{V_\epsilon}\).
The standard sector basis uses \(e_\pm\) for singles and
\(e_+\otimes e_+\), \(e_+\otimes e_-+e_-\otimes e_+\),
\(e_-\otimes e_-\) for symmetric pairs, with
\(e_+\otimes e_--e_-\otimes e_+\) for antisymmetric pairs.
\end{definition}

We realize all these spaces as subspaces of
\((\mathbb C^2)^{\otimes k}\) with its usual Hermitian inner
product; the pair basis vectors are deliberately unnormalized.
The symmetric and exterior squares are orthogonal, so
\(V=\bigoplus_\epsilon V_\epsilon\) is an orthogonal direct sum.
The central-sector basis is the subset of the indicated tensor
basis having total weight \(\kappa\).  These basis vectors are
orthogonal, have positive squared norms that are powers of two,
and have coordinates in \(\{0,1,-1\}\) in the full tensor basis.

The three total spin operators preserve each such tensor
sector.  Indeed, on a pair put
\[
 v_+=e_+\otimes e_+,\quad
 v_0=e_+\otimes e_-+e_-\otimes e_+,\quad
 v_-=e_-\otimes e_-,\quad
 w=e_+\otimes e_--e_-\otimes e_+.
\]
The pair's \(S_z\)-weights on \(v_+,v_0,v_-,w\) are
\(1,0,-1,0\), respectively, and its nonzero ladder actions are
\[
 S_-v_+=v_0,\quad S_-v_0=2v_-,\qquad
 S_+v_-=v_0,\quad S_+v_0=2v_+.
\]
Both ladders kill \(w\); all ladder actions not displayed on
these four vectors are zero.  On a single factor the weights
are \(1/2,-1/2\), and \(S_-e_+=e_-\), \(S_+e_-=e_+\),
with the other two ladder actions zero.  The total operators
are sums of these local actions.  In particular every weight
in a tensor sector is in \(\kappa+\mathbb Z\).  On each sector,
direct multiplication of the single-factor matrices gives
\[
 S_z^*=S_z,\qquad S_+^*=S_-,\qquad
 [S_z,S_\pm]=\pm S_\pm,\qquad [S_+,S_-]=2S_z.
\]
Operators on different tensor factors commute, so the same
identities hold for their sums and their sector restrictions.

\begin{lemma}\label{l:sector}
Let \(V'\) be any tensor sector formed from singles and
symmetric or exterior pair factors as above.
If a Hermitian \(\mathcal O'\) on \(V'\) commutes with
\(S_z,S_+,S_-\), then it preserves the central-weight space
\(V'^\circ\), and its restriction there has the same set of
eigenvalues as \(\mathcal O'\).
In particular \(\mathcal O'\) is positive definite if and only if
that restriction is positive definite.
\end{lemma}

\begin{proof}
Commutation with \(S_z\) immediately implies preservation of
\(V'^\circ\).  Every eigenvalue of the restriction is therefore
an eigenvalue on \(V'\).  Conversely, let \(\lambda\) be an
eigenvalue of \(\mathcal O'\), and let its nonzero eigenspace
be \(E_\lambda\).  All three spin operators preserve
\(E_\lambda\).  The restriction of \(S_z\) to this space is
Hermitian, so there is a nonzero \(v\in E_\lambda\) with
\(S_zv=mv\).  The explicit tensor weights above show that
\(m\in\kappa+\mathbb Z\).

For any vector of weight \(m\), the adjoint and commutator
identities give
\[
 \|S_-v\|^2-\|S_+v\|^2
       =\langle v,[S_+,S_-]v\rangle=2m\|v\|^2.
\]
Thus \(S_-v\ne0\) if \(m>0\), and \(S_+v\ne0\) if \(m<0\).
The commutators with \(S_z\) show that these vectors have weights
\(m-1\) and \(m+1\), respectively.  Commutation with
\(\mathcal O'\) keeps every nonzero ladder vector in
\(E_\lambda\).

If \(m>\kappa\), then \(m-\kappa\) is a positive integer.
Apply \(S_-\) that many times.  Every weight before the final
step is strictly positive, so no step vanishes, and the final
vector has weight \(\kappa\).  If \(m<\kappa\), then
\(\kappa-m\) is a positive integer.  Every weight before the
final step of successive raising is negative: it is at most
\(-1\) for \(\kappa=0\), and at most \(-1/2\) for
\(\kappa=1/2\).  Hence raising \(\kappa-m\) times also gives a
nonzero vector of weight \(\kappa\).  If \(m=\kappa\), take
\(v\) itself.  In all cases \(\lambda\) occurs on
\(V'^\circ\).

The two Hermitian operators consequently have the same set of
eigenvalues.  A finite-dimensional Hermitian operator is
positive definite exactly when all its eigenvalues are positive,
by its orthonormal eigenbasis, proving the last assertion.
This argument asserts nothing about eigenvalue multiplicities.
\end{proof}

\subsection{Integer positivity and reconstruction}

\begin{lemma}\label{l:gram-residual}
Let \(M\) be a symmetric integer matrix of size \(n\).
Let \(g,d,S\) be positive integers and \(L\) an integer
lower-triangular \(n\times n\) matrix; no invertibility is required.
Define the symmetric integer residual
\[
 E=dS^2M-gLL^{\mathsf T}.
\]
If, for every row,
\[
 E_{ii}>0,\qquad E_{ii}>\sum_{j\ne i}|E_{ij}|,
\]
then \(E\) and \(M\) are positive definite over \(\mathbb R\)
and, as Hermitian matrices, over \(\mathbb C\).
The case \(L=0\), \(g=d=S=1\), is the direct strict
diagonal-dominance certificate.
All hypotheses can be checked with integer sums, products,
absolute values and comparisons, without division or determinants.
\end{lemma}

\begin{proof}
Let \(z\in\mathbb C^n\).  Symmetry and reality of \(E\) imply
\begin{align*}
 z^*Ez
 &=\sum_iE_{ii}|z_i|^2+
           2\sum_{i<j}E_{ij}\operatorname{Re}(\overline z_i z_j)\\
 &\ge\sum_iE_{ii}|z_i|^2-
           2\sum_{i<j}|E_{ij}|\,|z_i|\,|z_j|\\
 &\ge\sum_i\left(E_{ii}-\sum_{j\ne i}|E_{ij}|\right)|z_i|^2.
\end{align*}
The last step uses \(2ab\le a^2+b^2\) for nonnegative \(a,b\).
Every parenthesized coefficient is positive, so \(z^*Ez>0\)
for every nonzero \(z\).  Also
\[
 dS^2 z^*Mz=z^*Ez+g z^*LL^{\mathsf T}z
             =z^*Ez+g\|L^{\mathsf T}z\|^2>0
             \qquad(z\ne0).
\]
Since \(dS^2>0\), this proves positivity of \(M\) over
\(\mathbb C\), and restriction proves it over \(\mathbb R\).
The argument also covers size zero, when positive definiteness
is vacuous.  No rank condition on \(L\) was used.

In particular, lower triangularity lets one form its Gram entries
as the integer sums
\[
 (LL^{\mathsf T})_{ij}
       =\sum_{a=1}^{\min(i,j)}L_{ia}L_{ja}
       \quad(1\le i,j\le n).
\]
The hypotheses and residual entries therefore require only the
listed integer operations and comparisons.  Setting \(L=0\)
and \(g=d=S=1\) leaves \(E=M\), giving the direct certificate.
\end{proof}

\begin{lemma}\label{l:integercheck}
For \(c=p/q\in(0,1)\) rational with positive integers \(p,q\),
choose a positive integer \(D_0\) clearing the coefficients of
\(3D_0\widehat Z_{\rm cur}\) and \(5D_0\widehat Z_{\rm next}\).
In the integer-coordinate central-sector basis, let \(A_\epsilon\)
and \(A'_\epsilon\) be the Hermitian-form matrices of their
inverted Pauli operators. These are symmetric integer matrices.
The form of \(\mathcal O(\widehat F_c^{\mathcal I})\) on the sector
is positive definite if and only if
\[
 M_\epsilon=5qA_\epsilon-3pA'_\epsilon
\]
is positive definite. The primary verification uses the integer
Gram/residual witness of Lemma~\ref{l:gram-residual}.
As an optional fallback only, positivity is equivalent to positivity of all its
leading principal minors. More explicitly, for a size-\(t\) matrix
set \(A^{(0)}=M_\epsilon\), \(p_{-1}=1\), and for \(0\le j<t\)
set \(p_j=A^{(j)}_{jj}\) and update the trailing entries by
\[
 A^{(j+1)}_{il}
    =\frac{p_jA^{(j)}_{il}-A^{(j)}_{ij}A^{(j)}_{jl}}{p_{j-1}}
       \quad(i,l>j).
\]
As long as preceding pivots are nonzero, each division is exact
in the integers and \(p_j\) is the leading principal minor of
order \(j+1\). Acceptance requires every division exact and every
pivot positive.
\end{lemma}

\begin{proof}
We first specify a reconstruction directly from the pattern;
this also gives an explicit choice of denominator.  Let \(X\)
be the set of spins still summed after the passage and the
chosen leaf reductions, and let \(B\) be their remaining present
internal bonds.  For each remaining vector \(u_i\), write
\(x(i)\in X\) for the spin to which it is attached.  Put
\(r_i=1\) for a retained corner slot and give a single vector
its radius \(r_i\in\{3/5,1/2,1/3\}\) from
Definition~\ref{d:reduction}.  The common remaining weight is
\[
 W(n,u)=\prod_{\{x,y\}\in B}(1+n_x\cdot n_y)
                     \prod_{i=1}^k(1+r_i u_i\cdot n_{x(i)}).
\]
If \(\mathcal T\) denotes the product of the original status
laws on \(X\), Lemma~\ref{l:Iw} and the definition of the
representation give
\[
 \widehat Z_{\rm cur}(u)=\mathcal T[W(n,u)I(n)],\qquad
 \widehat Z_{\rm next}(u)=\mathcal T[W(n,u)J(n)].
\]
All nonremoved sites among \(f,g,o\) remain in \(X\); they are
never periphery leaves.  Isolated spins may be retained in
\(X\), contributing their zeroth moment one.

For a word \(w=(j_1,\ldots,j_s)\) of coordinate indices in
\(\{1,2,3\}\), define the local moment
\(\mu_{\mathrm{st}}(w)=\lambda_{\mathrm{st}}
(n_{j_1}\cdots n_{j_s})\).  Only words of length at most four
will occur.  Their exact values, by Lemma~\ref{l:moments}, are
\[
 \begin{array}{c|cc}
 w&\mu_{\mathrm U}(w)&\mu_{\mathrm D}(w)\\ \hline
 \varnothing&1&1\\
 (i)\text{ or }(i,j,l)&0&0\\
 (i,j)&\delta_{ij}/3&\delta_{ij}/5\\
 (i,j,l,m)&
 \dfrac{\delta_{ij}\delta_{lm}+\delta_{il}\delta_{jm}
                         +\delta_{im}\delta_{jl}}{15}&
 \dfrac{\delta_{ij}\delta_{lm}+\delta_{il}\delta_{jm}
                         +\delta_{im}\delta_{jl}}{25}.
 \end{array}
\]
In particular every \(\mathrm U\) moment has denominator dividing
\(15\), and every \(\mathrm D\) moment has denominator dividing
\(25\).

Here is an entrywise formula for the entire coefficient tensor.
For \(\nu\in\{\mathrm{cur},\mathrm{next}\}\), let
\(\mathcal H_\nu\) contain an empty insertion of weight one.
For \(\nu=\mathrm{cur}\), also include \((v,w,j)\) of weight
\(1/3\) for each pair \(vw\in\{fg,fo,go\}\) whose endpoints are
nonremoved and each \(1\le j\le3\).  For
\(\nu=\mathrm{next}\), include just the nonremoved pair \(fg\)
with these three coordinate choices, each of weight \(1/5\).
Denote the insertion weight by \(\omega_\nu(h)\).
Fix a vector multi-index \(a\in\{0,1,2,3\}^k\) and bond
indices \(b\in\{0,1,2,3\}^{B}\).  At \(x\in X\), form a word
\(w_x(a,b,h)\) by collecting:
the index \(a_i\) for each \(i\) with \(x(i)=x\) and
\(a_i\ne0\); the index \(b_e\) for each incident bond \(e\)
with \(b_e\ne0\); and the insertion coordinate \(j\) if
\(h=(v,w,j)\) and \(x\) is one of \(v,w\).
The order of these indices is immaterial.  Expansion of each
factor as a sum of its constant and three coordinate terms gives
\[
 z_{\nu,a}=
 \left(\prod_{i=1}^k r_i^{\mathbf1_{\{a_i\ne0\}}}\right)
 \sum_{b\in\{0,1,2,3\}^{B}}\ \sum_{h\in\mathcal H_\nu}
 \omega_\nu(h)\prod_{x\in X}
          \mu_{\mathrm{st}(x)}\bigl(w_x(a,b,h)\bigr),
 \qquad
 \widehat Z_\nu(u)=\sum_a z_{\nu,a}\prod_i u_i^{a_i}.
\]
Indeed, each choice of terms from the product determines just
such \(a,b,h\); conversely each such triple determines one term,
with the indicated scalar radius and insertion weights.  The
product functional of that term is the product of its local
moments.  This proves the formula by finite distributivity.
The degree check in Corollary~\ref{c:reduce} shows explicitly
that every \(w_x\) has length at most four.

Write \(r_i=\alpha_i/\beta_i\) with positive integers
\(\alpha_i,\beta_i\), using \(\beta_i=1\) for retained slots.
If \(u\) and \(v\) are the numbers of \(\mathrm U\) and
\(\mathrm D\) sites in \(X\), a permissible positive denominator
is
\[
 D_0=15^u25^v\prod_{i=1}^k\beta_i.
\]
Each summand in the coefficient formula, after multiplication
by \(3D_0\) for the current polynomial or \(5D_0\) for the next,
is an integer: a vector radius is used at most once, one local
moment is used at each site, and the only remaining insertion
denominator is \(3\) or \(5\), respectively.  Any smaller
positive \(D_0\) is equally valid if the two coefficient
integrality conditions are verified.  Thus the existence and
positivity of a denominator follow from the pattern rather than
from stored matrix data.

For clarity, the remaining Pauli and basis transformations also
have exact entrywise formulas.  Set
\[
 T_{{\rm cur},a}=3D_0z_{{\rm cur},a},\qquad
 T_{{\rm next},a}=5D_0z_{{\rm next},a},\qquad
 \chi_{\mathcal I}(a)=(-1)^{|\{i\in\mathcal I:a_i\ne0\}|}.
\]
Index the full tensor basis by
\(s,t\in\{+,-\}^k\).  For either polynomial the inverted,
scaled Pauli operator has entries
\[
 (P_\nu)_{st}=\sum_a
       T_{\nu,a}\chi_{\mathcal I}(a)
                         \prod_{i=1}^k(\sigma_{a_i})_{s_i t_i}.
\]
Let \(B_\epsilon\) be the rectangular matrix whose columns are
the integer-coordinate basis vectors of \(V_\epsilon^\circ\)
in Definition~\ref{d:exchange}, expanded in this full tensor
basis, in any fixed order.  The requested matrices are exactly
\[
 A_\epsilon=B_\epsilon^*P_{\rm cur}B_\epsilon,
 \qquad A'_\epsilon=B_\epsilon^*P_{\rm next}B_\epsilon.
\]
The moment formula, inversion signs, explicit Pauli entries, and
explicit basis columns therefore reconstruct every matrix entry
by finite sums and products from the pattern and the chosen
representation.  No matrix or pivot list is an additional
mathematical input.

Since the \(T_{\nu,a}\) are integers and the Pauli matrices
have entries in \(\mathbb Z[\mathrm i]\), the \(P_\nu\) have
entries in \(\mathbb Z[\mathrm i]\).  The reality and symmetry
proof of Lemma~\ref{l:Ostruct} applies separately to each
polynomial, so the \(P_\nu\) are real symmetric.  Their entries
are thus integers.  The columns of \(B_\epsilon\) are real
integer vectors, proving that \(A_\epsilon,A'_\epsilon\) are
symmetric integer matrices.  Coefficient reflection in the
second coordinate alternatively shows that only terms with
an even number of \(\sigma_2\) factors can be nonzero, so the
same construction can be evaluated with real integer entries.

Linearity of the Pauli map gives the exact identity
\begin{align*}
 M_\epsilon
 &=5qA_\epsilon-3pA'_\epsilon\\
 &=15D_0q\,B_\epsilon^*
                \mathcal O(\widehat F_c^{\mathcal I})B_\epsilon.
\end{align*}
The columns of \(B_\epsilon\) form a basis of the central sector,
so \(B_\epsilon\) is an isomorphism from its coordinate space
onto that sector.  The last matrix on the right is the matrix
of the restricted Hermitian form, even though the basis is
unnormalized.  Since \(15D_0q>0\), its form is positive definite
exactly when \(M_\epsilon\) is positive definite.  This is the
matrix to which Lemma~\ref{l:gram-residual} applies in the primary
verification.

We finish by proving the optional determinant test and its
fraction-free implementation.  For any real symmetric matrix,
positive definiteness is equivalent to positivity of every
leading principal minor.  One elementary proof is induction on
the size.  If the first diagonal entry is \(a>0\), write the
matrix as
\[
 M=\begin{pmatrix}a&r^{\mathsf T}\\r&C\end{pmatrix},
 \qquad Q=C-rr^{\mathsf T}/a.
\]
Completing the square gives
\[
 (x,y)^{\mathsf T}M(x,y)
       =a(x+r^{\mathsf T}y/a)^2+y^{\mathsf T}Qy.
\]
Thus \(M\) is positive definite exactly when \(a>0\) and \(Q\)
is positive definite.  A block row elimination of determinant
one shows that every leading minor of size \(r\ge2\) of \(M\)
is \(a\) times the leading minor of size \(r-1\) of \(Q\).
Induction proves both directions, including size one; size zero
has no conditions.  This also proves the Hermitian-form version
for a real matrix, since for \(z=x+\mathrm i y\) its Hermitian
quadratic form is \(x^{\mathsf T}Mx+y^{\mathsf T}My\).

For the recurrence in the statement, index rows and columns
by \(0,\ldots,t-1\), and let \(\Delta_j\) be the determinant
of the leading \(j\times j\) block of \(M_\epsilon\), with
\(\Delta_0=1\).  For \(i,l\ge j\) the induction invariant is
\[
 A^{(j)}_{il}=
  \det M_\epsilon[(0,\ldots,j-1,i),(0,\ldots,j-1,l)].
\]
The ordered lists here specify rows and columns in that order;
for \(j=0\) this is the original entry.  In particular
\(p_j=\Delta_{j+1}\) and \(p_{j-1}=\Delta_j\).
Assume the preceding pivots are nonzero.  The bordered-minor
identity is
\begin{align*}
 &A^{(j)}_{jj}A^{(j)}_{il}
                    -A^{(j)}_{ij}A^{(j)}_{jl}\\
 &\hspace{1cm}=\Delta_j\,
  \det M_\epsilon[(0,\ldots,j,i),(0,\ldots,j,l)]
                       \qquad(i,l>j).
\end{align*}
For \(j=0\) it is the \(2\times2\) determinant formula.  For
\(j>0\), the leading block \(B_j\) of size \(j\) is invertible
because \(\det B_j=\Delta_j\ne0\).  In the submatrix with the
two added row and column indices, eliminate the entries below
\(B_j\).  If the resulting \(2\times2\) Schur block has entries
\(q_{ab}\), each bordered \((j+1)\)-minor is
\(\Delta_j q_{ab}\), while the full \((j+2)\)-minor is
\(\Delta_j(q_{jj}q_{il}-q_{ij}q_{jl})\).
Substitution proves the displayed identity.  The quotient by
\(p_{j-1}=\Delta_j\) is therefore the indicated next bordered
minor, an integer.  This proves exact divisibility, the induction
invariant, and the pivot formula throughout the recurrence.
All positive pivots consequently imply all positive leading
minors, and conversely a positive definite matrix yields exact
divisions and positive pivots at every step.  These determinant
calculations are needed only for this optional fallback.
\end{proof}

\begin{proposition}\label{p:certificate}
Let \(c=p/q\in(0,1)\), with \(p,q\in\mathbb Z_{>0}\).
Choose either allowed representation, an inversion which is a
union of groups, and a justified positive denominator \(D_0\).
For every supported central exchange sector reconstruct
\(M_\epsilon=5qA_\epsilon-3pA'_\epsilon\) as in
Lemma~\ref{l:integercheck}.
Suppose that for each sector there are positive integers
\(g_\epsilon,d_\epsilon,S_\epsilon\) and an integer lower-triangular
matrix \(L_\epsilon\) of its dimension such that
\[
 E_\epsilon=d_\epsilon S_\epsilon^2M_\epsilon
                  -g_\epsilon L_\epsilon L_\epsilon^{\mathsf T},
 \qquad
 (E_\epsilon)_{ii}>
          \sum_{j\ne i}|(E_\epsilon)_{ij}|\quad\text{for every }i.
\]
Then \(F_c(m)\ge0\) for every \(m\in(\mathbb S^2)^{20}\).
The operator may vanish on the discarded antisymmetric
\(\mathrm U\)-pair subspaces; positive definiteness is required
only on the supported exchange sectors.
\end{proposition}

\begin{proof}
Let \(\mathcal O=\mathcal O(\widehat F_c^{\mathcal I})\).
For every sector the residual is real symmetric, since its
reconstructed \(M_\epsilon\) and its Gram matrix are symmetric.
Its strict row inequality forces
\((E_\epsilon)_{ii}>0\), because the sum of absolute values is
nonnegative.  All hypotheses of Lemma~\ref{l:gram-residual}
therefore hold, and every \(M_\epsilon\) is positive definite.
Lemma~\ref{l:integercheck}, including its positive scale
\(15D_0q\), makes the form of \(\mathcal O\) positive definite
on every \(V_\epsilon^\circ\).

Let \(\Pi_V\) be the product of the symmetric projections for
all retained \(\mathrm U\)-pairs, tensored with the identities
on all remaining groups.  These projections act on disjoint
factors and commute.  Lemma~\ref{l:Ostruct} gives
\(\mathcal O=\Pi_V\mathcal O\Pi_V\), with
\(\operatorname{Ran}\Pi_V=V\).  In particular \(\mathcal O\)
vanishes on \(V^\perp\) and maps \(V\) into itself.
Its commutation with every retained \(\mathrm D\)-pair exchange
also makes every \(V_\epsilon\) invariant.  The concrete spin
operators preserve these sectors, as checked after
Definition~\ref{d:exchange}; Lemma~\ref{l:Ostruct} shows that
the restriction of \(\mathcal O\) to each sector commutes with
their restrictions.  Lemma~\ref{l:sector} now promotes positive
definiteness on \(V_\epsilon^\circ\) to positive definiteness
on \(V_\epsilon\).

The direct sum \(V=\bigoplus_\epsilon V_\epsilon\) is orthogonal.
For any full tensor vector \(z\), write
\(\Pi_Vz=\sum_\epsilon z_\epsilon\) with
\(z_\epsilon\in V_\epsilon\).  Invariance and orthogonality give
\[
 \langle z,\mathcal Oz\rangle
       =\sum_\epsilon\langle z_\epsilon,
                                  \mathcal Oz_\epsilon\rangle\ge0.
\]
Thus \(\mathcal O\) is positive semidefinite on the full tensor
space, allowing precisely the discarded antisymmetric
\(\mathrm U\)-pair subspaces to have zero action.
Lemma~\ref{l:inversion} implies \(\widehat F_c\ge0\) on the
product of closed unit balls, hence on unit vectors.
Corollary~\ref{c:reduce} gives \(F_c(m)\ge0\) for every original
twenty-slot assignment, proving the proposition.
\end{proof}



\begin{thebibliography}{99}

\bibitem{AKLT88}
I.~Affleck, T.~Kennedy, E.~H.~Lieb, and H.~Tasaki,
\emph{Valence bond ground states in isotropic quantum antiferromagnets},
Comm.\ Math.\ Phys.\ \textbf{115} (1988), 477--528.

\bibitem{KLT88}
T.~Kennedy, E.~H.~Lieb, and H.~Tasaki,
\emph{A two-dimensional isotropic quantum antiferromagnet with unique disordered ground state},
J.\ Statist.\ Phys.\ \textbf{53} (1988), 383--415.

\end{thebibliography}

\end{document}
